%% LyX 2.3.1-1 created this file.  For more info, see http://www.lyx.org/.
%% Do not edit unless you really know what you are doing.
\documentclass[english,pdftex, a4paper, 12pt]{article}
\usepackage[T1]{fontenc}
\usepackage{geometry}
\geometry{verbose,tmargin=1in,bmargin=1in,lmargin=1.25in,rmargin=1.25in}
\usepackage{verbatim}
\usepackage{amsmath}
\usepackage{amsthm}
\usepackage{amssymb}
\usepackage{graphicx}
\usepackage{setspace}
\usepackage[authoryear]{natbib}
\onehalfspacing

\makeatletter

%%%%%%%%%%%%%%%%%%%%%%%%%%%%%% LyX specific LaTeX commands.
\newcommand{\noun}[1]{\textsc{#1}}

%%%%%%%%%%%%%%%%%%%%%%%%%%%%%% Textclass specific LaTeX commands.
\theoremstyle{plain}
\newtheorem{prop}{\protect\propositionname}
\theoremstyle{plain}
\newtheorem{lem}{\protect\lemmaname}

%%%%%%%%%%%%%%%%%%%%%%%%%%%%%% User specified LaTeX commands.
\PassOptionsToPackage{cmyk}{xcolor}

%%%%%%%%%%%%%%%%%%%%%%%%%%%%%%%%%%%%%%%%%%%%%%%%%%%%%%%%%%%%%%%%%%%%%%%%%%%%%%%
% BABEL and LANGUAGES %%%%%%%%%%%%%%%%%%%%%%%%%%%%%%%%%%%%%%%%%%%%%%%%%%%%%%%%%
%%%%%%%%%%%%%%%%%%%%%%%%%%%%%%%%%%%%%%%%%%%%%%%%%%%%%%%%%%%%%%%%%%%%%%%%%%%%%%%
% \usepackage{listings}                   % it is a source code printer for LATEX
                                          % \lstset{language=Python}
                                          % \lstinputlisting{source.py}   % command used to pretty-print stand alone files
\usepackage[english]{babel}               % [french, frenchb, english, ]
    % http://forum.mathematex.net/latex-f6/les-puces-avec-babel-t4256.html
    % http://www.grappa.univ-lille3.fr/FAQ-LaTeX/11.1.html


%%%%%%%%%%%%%%%%%%%%%%%%%%%%%%%%%%%%%%%%%%%%%%%%%%%%%%%%%%%%%%%%%%%%%%%%%%%%%%%
% FONTS and ENCODING %%%%%%%%%%%%%%%%%%%%%%%%%%%%%%%%%%%%%%%%%%%%%%%%%%%%%%%%%%
%%%%%%%%%%%%%%%%%%%%%%%%%%%%%%%%%%%%%%%%%%%%%%%%%%%%%%%%%%%%%%%%%%%%%%%%%%%%%%%
%
% See:
% https://tex.stackexchange.com/questions/59702/suggest-a-nice-font-family-for-my-basic-latex-template-text-and-math-i-am
%

\usepackage{lmodern}        % Latin Modern family of fonts. Very much like Computer Modern, but with many more glyphs 
                            % (e.g., for characters with accents, glyphs, cedillas, etc)
\usepackage[T1]{fontenc}    % fontenc is oriented to output, that is, what fonts to use for printing characters. 
                            % https://tex.stackexchange.com/questions/44694/fontenc-vs-inputenc 
                            % https://tex.stackexchange.com/questions/664/why-should-i-use-usepackaget1fontenc

% Change some fonts or the whole font family (i.e. serif, sans serif, monospace, and 'math')
    % \usepackage[varg, cmintegrals, cmbraces, ]{newtxtext,newtxmath}  % Other options: libertine, uprightGreek (U.S.) or slantedGreek (ISO), etc...
    % \usepackage{tgtermes}                                            % Only serif ("TeX-Gyre" text)
    % \usepackage{kpfonts}                                             % "Kepler" fonts
      \usepackage{mathpazo}                                          % Based on Hermann Zapf's Palatino font
    % \usepackage{txfonts}                                             % More than a decade old
    % \usepackage{pslatex}                                             % Obsolete?
    % \usepackage{mathptmx}
    %  - \usepackage[scaled=.90]{helvet}
    %  - \usepackage{courier}
	%\usepackage{mathpple}


% \usepackage{textcomp}     % required for special glyphs
% \usepackage{bm}           % load after all math to give access to bold math
\usepackage[utf8]{inputenc} % inputenc allows the user to input accented characters directly from the keyboard; 
                            % utf8x : much broader but less compatible ; latin1 : old?
                            % https://tex.stackexchange.com/questions/44694/fontenc-vs-inputenc

% See:
% https://tex.stackexchange.com/questions/59626/nicely-force-66-characters-per-line
%
% pslatex is a very obsolete package and that its descendant mathptmx is rather inadequate for serious typesetting involving math.
% If you don't need mathematics, other choices based on (Linotype) Times Roman are
%  - tgtermes
%  - newtxtext (based on txfonts, but with corrected metrics) (with its companion math package newtxmath)
%
%
% See:
% http://www.latex-community.org/forum/viewtopic.php?f=8&t=6637
%
% (times, helvet, courier)
% pslatex and txfonts produce (almost) same resutls.
% pslatex supposedly obsolete
% txfonts supposedly up-to-date
%
%
% See:
% ftp://ftp.rrzn.uni-hannover.de/pub/mirror/tex-archive/info/l2tabu/english/l2tabuen.pdf
% or 
% ftp://ftp.dante.de/tex-archive/info/l2tabu/english/l2tabuen.pdf
% in
% 2.3.3 pslatex.sty
%
% pslatex uses a Courier font scaled too narrowly.
% Its main disadvantage is that it does not work with T1 and TS1 encodings.
% So replace:
% \usepackage{pslatex} or \usepackage{txfonts}
% by all three:
% - \usepackage{mathptmx}
% - \usepackage[scaled=.90]{helvet}
% - \usepackage{courier}
%
%
% See:
% http://xpt.sourceforge.net/techdocs/language/latex/latex32-LaTeXAndFonts/single/
% or http://thirteen-01.stat.iastate.edu/wiki/LaTeXFonts
% or http://www.tex.ac.uk/tex-archive/info/beginlatex/html/chapter8.html
%
% When changing fonts, you can change all of the default fonts at once with the following commands:
% 
% Command     Changes the defaults to
% 
% times       Times, Helvetica, Courier
% pslatex     same as Times, but uses a specially narrowed Courier. This is preferred over Times because of the way it handles Courier.
% newcent     New Century Schoolbook, Avant Garde, Courier
% palatino    Palatino, Helevetica, Courier
% palatcm     changes the Roman to Palatino only, but uses CM mathematics
% kpfonts     "Kepler" fonts. A very nicely evolved set of fonts also based originally on Palatino, but with many special features.
%
%
% See:
% https://tex.stackexchange.com/questions/59702/suggest-a-nice-font-family-for-my-basic-latex-template-text-and-math-i-am
%
% There are, of course, many other font packages, most of which provide "only" text-mode fonts.
% Among these are the "TeX-Gyre" font families: 
%  - Termes (a Times Roman clone), 
%  - Pagella (a Palatino clone), and 
%  - Schola (a Century Schoolbook clone); 
% one would load the packages tgtermes, tgpagella, and tgschola, respectively, to access these fonts.
% However, as these are text fonts, you still need to choose a suitable math font.
% 
% Still another possibility you may want to look into is the Linux Libertine font family, to be loaded via the libertine-legacy package.
% If you like this text font and wish to employ the newtxmath package, be sure to load the newtxmath package with the libertine option set;
% doing so will set up a special set of math-mode fonts that harmonizes well with the libertine text fonts.
% 
%
% See also:
% https://tex.stackexchange.com/questions/56876/times-new-roman-fonts-and-maths-without-mathptmx
%
%
% For a comparison, in:
% /home/christophe/Personal/Truc_Et_Astuce_Informatik/LaTeX/comparison_font_types/,
% see: 
% computer.pdf  lmodern.pdf  pslatex.pdf  test_font_type.pdf  three_replacements.pdf  txfonts.pdf
%


%%%%%%%%%%%%%%%%%%%%%%%%%%%%%%%%%%%%%%%%%%%%%%%%%%%%%%%%%%%%%%%%%%%%%%%%%%%%%%%
% AMS MATH %%%%%%%%%%%%%%%%%%%%%%%%%%%%%%%%%%%%%%%%%%%%%%%%%%%%%%%%%%%%%%%%%%%%
%%%%%%%%%%%%%%%%%%%%%%%%%%%%%%%%%%%%%%%%%%%%%%%%%%%%%%%%%%%%%%%%%%%%%%%%%%%%%%%
% \usepackage{amsmath}      % loads amstext, amsbsy, amsopn but not amssymb
                            % equation stuff (eqref, subequations, equation, align, gather, flalign, multline, alignat, split...)
% \usepackage{amsfonts}     % may be redundant with amsmath
% \usepackage{amssymb}      % may be redundant with amsmath
% \numberwithin{equation}{section}  % reset equation counters at start of each "section" and prefix numbers by section number
% \numberwithin{figure}{section}    % reset figure   counters at start of each "section" and prefix numbers by section number


%%%%%%%%%%%%%%%%%%%%%%%%%%%%%%%%%%%%%%%%%%%%%%%%%%%%%%%%%%%%%%%%%%%%%%%%%%%%%%%
% LAY OUT %%%%%%%%%%%%%%%%%%%%%%%%%%%%%%%%%%%%%%%%%%%%%%%%%%%%%%%%%%%%%%%%%%%%%
%%%%%%%%%%%%%%%%%%%%%%%%%%%%%%%%%%%%%%%%%%%%%%%%%%%%%%%%%%%%%%%%%%%%%%%%%%%%%%%
%
% See:
% https://tex.stackexchange.com/questions/59626/nicely-force-66-characters-per-line
% (must be after pslatex, tgterms, etc...)
%
% a) (but works mostly for a4paper, and changes top and bottom margin too...)
% \usepackage[DIV=calc]{typearea}
%
% or
%
% b) (but you have to choose the value and the margin ratio depending on the class...)
% \newlength{\alphabet}
% \settowidth{\alphabet}{\normalfont abcdefghijklmnopqrstuvwxyz}
% \usepackage{geometry}
% \geometry{%
% textwidth=2.5\alphabet,% (Note: 2.5 * 26 = 65)
% hmarginratio={2:3}}    % (Problem: geometry uses 2:3 as default for twoside and 1:1 for oneside,
%                        % independently of what the class thinks about the margins)

% \usepackage{layout}       % use \layout in the tex file to see the values
% \usepackage{layouts}      % it extends the functionality of layout, allowing you to do much, much more
                            % some commands: \pagelayout, \pagevalues, \pagedesign, ...
% \usepackage[cm]{fullpage} % set 'default' full page
% \usepackage{geometry}     % very customizable margins. Under some (rare) circumstances, should be loaded after hyperref
% \usepackage{anysize}      % \marginsize{left}{right}{top}{bottom}
% \usepackage{pdflscape}    % include landscape layout pages (automatically rotate pages in pdf file for easier reading)
% \usepackage{multicol}     % for multi column environment
\usepackage{lipsum}         % to fill in with arbitrary text
\widowpenalty = 4000        % help suppress widows,  default = 4,000 (?), from 0 to 10 000 (from 300 to 1 000 recommended, 10 000 not recommended)
\clubpenalty  = 4000        % help suppress orphans, default = 4,000 (?), from 0 to 10 000 (from 300 to 1 000 recommended, 10 000 not recommended)
\usepackage[final, babel]{microtype} % many good lay-out/justification effects, see:
                                     % texblog.net/latex-archive/layout/pdflatex-microtype/


%%%%%%%%%%%%%%%%%%%%%%%%%%%%%%%%%%%%%%%%%%%%%%%%%%%%%%%%%%%%%%%%%%%%%%%%%%%%%%%
% GRAPHICX %%%%%%%%%%%%%%%%%%%%%%%%%%%%%%%%%%%%%%%%%%%%%%%%%%%%%%%%%%%%%%%%%%%%
%%%%%%%%%%%%%%%%%%%%%%%%%%%%%%%%%%%%%%%%%%%%%%%%%%%%%%%%%%%%%%%%%%%%%%%%%%%%%%%
% \usepackage[final]{graphicx} % options = [final]  = all graphics displayed, regardless of draft option in class
                               % options = [pdftex] = necessary (?) if import PDF files
                               % no option : when importing ps- and eps-files (?)
% \graphicspath{{../images/}}  % tell LaTeX where to look for images
% \DeclareGraphicsExtensions{.pdf, .PDF, .jpg, .JPG, .jpeg, .JPEG, .png, .PNG, .bmp, .BMP, .eps, .ps}
\usepackage{float}                      % Improved interface for floating objects ; add [H] option


%%%%%%%%%%%%%%%%%%%%%%%%%%%%%%%%%%%%%%%%%%%%%%%%%%%%%%%%%%%%%%%%%%%%%%%%%%%%%%%
% HYPERREF (last) then HYPCAP %%%%%%%%%%%%%%%%%%%%%%%%%%%%%%%%%%%%%%%%%%%%%%%%%
%%%%%%%%%%%%%%%%%%%%%%%%%%%%%%%%%%%%%%%%%%%%%%%%%%%%%%%%%%%%%%%%%%%%%%%%%%%%%%%
%
% See:
% https://tex.stackexchange.com/questions/1863/which-packages-should-be-loaded-after-hyperref-instead-of-before
%
\usepackage[
pdftex, 
final,                      % if you do    want to have clickable-colorful links
pdfstartview = FitV,
colorlinks = true,
citecolor = -myblue,
linktocpage  = false,       % ToC, LoF, LoT place hyperlink on page number, rather than entry text
breaklinks   = true,        % so long urls are correctly broken across lines
% pagebackref  = false,     % add page number in bibliography and link to position in document where cited
]{hyperref}

% \usepackage{cleveref} % enhance cross-referencing, allow full formatting, commands:
                        % \cref, \Cref, \crefrange, \cref{eq2,eq1,eq3,eq5,thm2,def3}
                        % supposedly better than \autoref as provided by hyperref

% \usepackage[all]{hypcap} % when link to float (using hyperref), link anchors to beginning (instead of below) float
%%%%%%%%%%%%%%%%%%%%%%%%%%
\usepackage{graphicx}
\usepackage{booktabs}
\newcommand{\specialcell}[2][c]{%
  \begin{tabular}[#1]{@{}c@{}}#2\end{tabular}}
\usepackage{dcolumn}
\usepackage{lscape}
\usepackage{changepage}
\usepackage[cmyk, dvipsnames]{xcolor}
\definecolor{myblue}{RGB}{245,156,74}
\usepackage{floatpag}
\usepackage{adjustbox}
\usepackage{parskip}
\setlength{\parskip}{5pt}

\hypersetup{
colorlinks=true,
urlcolor=blue,
}
\usepackage{bbm}

\makeatother
\newcounter{parentnumber}
\newtheorem*{namedthm*}{\thistheoremname}
\newcommand{\thistheoremname}{} % initialization
\newenvironment{namedthm}[1]
  {\renewcommand{\thistheoremname}{#1}\begin{namedthm*}}
  {\end{namedthm*}}

\newtheorem{obsv}{ \scshape Regularity}
%%%%%%%%%%%%%%%%%%%%%%%%%%%%%%%%%%%%%%%%%%%%%%%%%%%%%%%%%%%%%%%%%%%%%%%%%%%%%%%
% TITLE PAGE %%%%%%%%%%%%%%%%%%%%%%%%%%%%%%%%%%%%%%%%%%%%%%%%%%%%%%%%%%%%%%%%%%
%%%%%%%%%%%%%%%%%%%%%%%%%%%%%%%%%%%%%%%%%%%%%%%%%%%%%%%%%%%%%%%%%%%%%%%%%%%%%%%
\title{Discrete Trade}
\author{Ahmad Lashkaripour}
\date{\today}
\usepackage{adjustbox}
\usepackage{multirow}
\usepackage[justification=centering]{caption}

\renewcommand\labelenumi{(\roman{enumi})}
\renewcommand\theenumi\labelenumi

\usepackage[margin=15pt,font=small,labelfont={bf,sf}]{caption}
\usepackage{microtype}
\usepackage{placeins}
\usepackage{enumitem}
\usepackage{bbm}
 \usepackage{relsize}

\makeatother

\usepackage{babel}
\providecommand{\lemmaname}{Lemma}
\providecommand{\propositionname}{Proposition}

\begin{document}
\maketitle
\begin{abstract}
\emph{International quality specialization} (IQS) and\emph{ pricing-to-market}
(PTM) are two of the most studied phenomena in international trade.
I present new evidence that PTM and IQS are significantly more pronounced
in discrete industries that involve indivisible consumption goods.
In light of this observation, I develop a generalized model of international
trade that accommodates both discrete and continuous industries. I
argue that IQS and PTM in discrete industries are magnified by $(i)$
affordability constraints, and $(ii)$ cross-country differences in
the price of non-traded services---both of which are less relevant
in continuous industries. This prediction finds support from micro-level
trade data. I then map the discrete model to macro-level data and
show that the unique forces underlying the model shed fresh light
on the gains from trade and the \textquotedblleft big-push\textquotedblright{}
effects of globalization. %
\begin{comment}

\paragraph*{Things to do: }

(i) include the term ``endogenous comparative advantage'' or ``endogenous
organization of prodction''? 

(ii) highlight the positive cross-country correlation between service
expenditure and pice in the introduction?

(iii) Discuss the model's implications for the intra-national organization
of economic activity.
\end{comment}
\end{abstract}

\section{Introduction}

\begin{comment}
\begin{itemize}
\item Control for differentiated goods based on the Rauch classification
\item Revise assumptions: drop bound on quality. State assumption about
buying.
\item Revise Lemmas 2 and 3
\item Discuss why continuous goods are not a substitute for services, but
rather an input in services.
\item When buying the service provided by discrete goods, consumers have
two options: $(i)$ pay the full price and benefit from the flow of
services, or $(ii)$ rent the service by paying a (near-zero) fraction
of the discrete good's price plus the labor involved in offering the
service.
\end{itemize}
\end{comment}
Around 50\% of all U.S. import transactions from 1989 to 2015 involved
\emph{discrete} goods such as ``Transport Equipment'', ``Electronics
\& Appliances,'' and ``Machinery.'' A defining feature of these
goods is that they are \textbf{\emph{indivisible}} and durable. That
is, consumers can purchase discrete goods, like cars, smartphones,
or washer-dryers, in only integer quantities. However, once purchased,
these goods supply a continuous flow of transportation, communication,
or laundry services.\footnote{Researchers often think of discrete goods as providing a flow of services:
Cars, for instance, provide a flow of transportation services, or
washer-dryer units provide a flow of dry-cleaning services---see
\citet{benhabib1991homework}, \citet{cooley1995economic}, and \citet{fernandez2011consumption},
among many others, for this analogy.}

Indivisibility of this kind has sharp implications for consumer choice.
Most importantly, consumers cannot buy a fraction of multiple discrete
goods. They are instead left with two options: $(a)$ pay the price
of one full unit, which is akin to paying a fixed acquisition cost;
or $(b)$ purchase a service that can substitute for such goods. They
can also combine these two options. For instance, a consumer may purchase
a personal vehicle for day-to-day use and spend the rest of their
transportation budget on more-specialized transportation services.
Either way and unlike what standard models assume, consumers cannot
just buy a small fraction of different varieties of cars for different
transportation demands.\footnote{There is always the option of renting or leasing multiple varieties
of a discrete good. Here, I treat that option as a type of service.
But even then, one should bear in mind that in most low-income countries
this option is not widely-accessible (see Banerjee (2003)). Also,
this option is less plausible for non-portable discrete goods such
as home appliances.}

These practical issues have received formal attention in several fields,
including \emph{industrial organization} and \emph{macroeconomics}.
Given the prevalence of discrete trade, one may also expect careful
modeling of these same issues in mainstream theories of international
trade. However, on the contrary, the issue of \emph{discreteness}
remains largely untreated in mainstream trade models.

A possible argument for this lack of attention is that discreteness
plays a less critical role when studying macro-level trade patterns.
But a basic analysis of product-level US trade data challenges this
argument. In particular, comparing discrete industries to continuous
(non-discrete) counterparts, such as ``Food,'' ``Chemicals,''
or ``Textile \& Apparel,'' indicates that
\begin{enumerate}
\item There is significantly more cross-national quality/price heterogeneity
within discrete industries;
\item There is a greater degree of international quality specialization
(IQS) in discrete industries; and 
\item There is a greater degree of pricing-to-market (PTM) in discrete industries.
\end{enumerate}
These differences remain significant, even after we control for other
characteristics like the intended final use or degree of product differentiation.
At face value, the above regularities suggest that in discrete industries,
trade patterns such as IQS and PTM are amplified by forces that are
dormant in continuous industries. Such forces are also inevitably
overlooked by mainstream theories that treat all industries or goods
as continuous. 

Motivated by this evidence, I develop a simple model of discrete trade
that pursues two objectives. First, I use the model to elucidate how
\emph{discreteness} may affect macro-level trade patterns (like IQS
and PTM) in ways that are difficult to understand from the lens of
standard trade models. Second, I use the discrete model to determine
how accounting for discrete trade revises our understanding of the
gains from trade---an issue that has received considerable attention
in the contemporary trade literature.

My theory accommodates both continuous and discrete industries. Continuous
industries are modeled in a fairly standard fashion. My modeling of
discrete industries, however, is grounded in three basic assumptions.
$(i)$ Discrete goods are indivisible, wherein purchasing them involves
paying a unit price that is akin to a fixed acquisition cost. $(ii)$
The consumption flow derived from a personally-acquired discrete good
is increasing in its quality and can be substituted with analogous
services. $(iii)$ Quality and quantity are not isomorphic: a firm
cannot produce a lower quality version of their product by employing
fewer inputs under the same production technology. Instead, any supply-side
adjustment to output quality involves paying a fixed R\&D cost.

Using the above model, I show that in \emph{discrete} industries,
IQS and PTM are magnified by two forces that are less relevant in
\emph{continuous} industries:
\begin{enumerate}
\item Cross-national heterogeneity in the price of non-traded services,
and
\item Affordability constraints, whereby consumers with a limited budget
cannot afford certain discrete varieties.
\end{enumerate}
Both of these forces operate in a somewhat similar fashion. They lead
to the \emph{market-level} demand elasticity and demand for quality
to diverge across rich and poor markets, even when \emph{individual-level}
preferences are homothetic and identical across rich and poor consumers.

The role of force $(i)$ in prompting cross-national demand differences
can be explained as follows. In each national market, discrete goods
compete head-to-head with services; but services are significantly
more expensive in high-wage economies (i.e., the Balassa-Samuelson
effect). As a result, the opportunity cost of acquiring a discrete
good is lower in high-wage markets, and this is true across all income
groups. These considerations give rise to two macro-level demand patterns.
First, the overall demand for discrete goods is \emph{less} price-elastic
in high-wage markets. Second, there is also more demand for high-quality
(high-cost) discrete varieties in these markets.

To give an example, consider a premium hair-clipper that is a discrete
tradable good. In Mumbai, India, the supplier of the hair-clipper
competes with \$5 haircut services, whereas in Oslo, Norway, it competes
with \$80 haircut services.\footnote{Source: \url{http://www.businessinsider.com/how-expensive-haircuts-are-around-the-world-2015-9}}
As a result, Indian consumers (across all income groups) will exhibit
a higher degree of price sensitivity when purchasing the hair-clipper.
The supplier will internalize this when setting their price or choosing
their output quality for these two markets.

Contrast this with continuous goods like shampoo or hair products.
Since these goods perish upon consumption, they cannot be directly
substituted with services. Instead, one has to (indirectly) pay for
them even if they use the barber service. The same goes for the gas
used in transportation services or the food used in restaurant services.
One can circumvent paying for the full price of a vehicle or a premium
grill, by taking a cab or going to a restaurant. But irrespective
of whether they use such services or operate their personal vehicle/grill,
they have to pay for the continuous perishable input, which is food
or gas in this example.

Force $(ii)$ prompts cross-national demand differences in a different
way. Since discrete goods cannot be purchased in fractional quantities
and come with a pre-customized quality, they can be unaffordable to
low-budget consumers. Consider, for example, the Toyota Corolla which,
in 2017, came with a price tag that was 8-times the annual wage rate
in India. In such circumstances, \emph{one} \noun{Toyota Corolla}
can be prohibitively expensive for a significant portion of Indian
consumers---especially given that most low-income consumers do not
have access to financing options (\citet{banerjee2003contracting}).
Correspondingly, a marginal increase in price can rule out a non-trivial
measure of potential customers.

The elasticity of demand in discrete industries, therefore, depends
on the hazard rate at which a price increase renders the product unaffordable
to marginal, low-budget consumers. This hazard rate is itself decreasing
in the discrete good's overall \emph{degree of affordability}, i.e.,
its price relative to the wage rate in the economy.\footnote{This relationship holds for any empirically-relevant income distribution,
since they exhibit an Increasing Generalized Failure Rate---see \citet{lariviere2001selling}
and \citet{lariviere2006note}.} These considerations, once again, give rise to two macro-level demand
patterns. First, the overall demand for discrete goods is \emph{less}
price-elastic in high-wage markets, where the good is on average more
affordable. Second, there is more demand for high-quality (high-cost)
varieties in these markets.

The cross-national demand differences triggered by Forces $(i)$ and
$(ii)$, can magnify the extent of PTM and IQS in \emph{discrete}
industries quite significantly. First, they prompt the suppliers of
discrete goods to set higher markups in high-wage markets where their
good is more affordable and also competes with expensive services.
Second, the \emph{affordability-driven} and \emph{service price-driven}
local demand for quality in high-wage economies prompts high-quality
specialization in these economies through scale effects. 

To test the relevance of these forces, I combine data on export prices,
export sales, and non-traded service prices for 147 countries across
5,231 6-digit (HS6) industries in the year 2011. As in the prior literature,
I document that in all industries export price levels depend systematically
on market size, income per capita, and geography. In discrete industries,
however, export prices increase with three additional factors that
are less-relevant in continuous industries: $(i)$ the price of non-traded
services in the importing market, $(ii)$ the price of non-traded
services in the exporting economy, and $(iii)$ the exporting economy's
overall cost efficiency in the industry of interest. While the empirical
significance of these additional factors confirms the predictions
of the model, the estimated effects are also quite profound. For example,
a 10\% increase in the price of services in an importing market is
associated with a 2.5\% increase in the price of \emph{discrete} imports.

On a broader level, accounting for discrete trade can revise our estimates
for the gains from trade. When the traded sector is modeled as discrete
rather than continuous, the predicted gains from trade are on average
28\% larger but more unequally distributed across countries. These
differences are driven by two factors. The first factor is somewhat
mechanical: The gains from discrete trade depend on the \emph{quantity}
share of foreign goods in national consumption, whereas the gains
from continuous trade depend on the \emph{value }share of foreign
goods in national expenditure (\citet{Arkolakis}). 

The second factor is more profound. The discrete model predicts scale-driven
gains from trade that are absent in the continuous model. To elaborate,
in the continuous model, the scale of national production in the traded
sector is invariant to trade.\footnote{More specifically, the national scale of production in the traded
sector is determined exclusively by population size. This result will
be revised if we model the traded sector as a blend of multiple industries
with different scale elasticities---see \citet{costinot2014trade}
and \citet{kucheryavyy2016grounded}.} In the discrete model, however, opening to trade induces firm-relocation
from large to small and low-income economies. This type of relocation
allows firms to collect a higher markup on marginal cost when selling
to large and rich international markets. It also expands the scale
of production in small and low-income economies at the expense of
others. 

Consider Canada, for instance, which is located close to the richer
and larger U.S. market. From the lens of the discrete model, trade-induced
scale expansion accounts for 5.6\% of Canada's real GDP. Intuitively,
trade openness induces firms to relocate from the U.S. to Canada,
using it as a platform to export to the U.S. at a higher markup. Given
the sheer size of the U.S. economy, these developments provide a sizable
boost for Canada's real GDP. They also relate to the \emph{``big-push''}
effects of economic development, wherein the growth of an economy
may be hindered by its small local market (\citet{rosenstein1943problems,lewis1954economic,murphy1989industrialization}).
But through firm-relocation, international trade can push such an
economy out of this supposed growth trap.

Aside from contributing to ongoing debates about the size and nature
of the gains from trade, the present paper contributes to our understanding
of IQS and PTM. Existing theories of IQS and PTM often model industries
as continuous.\footnote{See \citet{Schott2004,Hallak2006,choi2009explaining,Baldwin,lugovskyy2012geography,feenstra2014international,antoniades2015heterogeneous,jaimovich2015love,alcala2016specialization}
for papers studying IQS using continuous good frameworks, as well
as \citet{atkeson2008pricing,hummels2009international,bekkers2015balassa,chen2014quality,simonovska2015income,bertoletti2018international}
for papers studying PTM using continuous good frameworks.} As a result, they naturally emphasize non-homothetic preferences
and cross-national technology differences as the drivers of PTM and
IQS. Prior theories of IQS and PTM that admit discrete goods (e.g.,
\citet{Flam1987,murphy1997quality,verhoogen2008trade,Khandelwal2010,Fajgelbaum2011,dingel2016determinants,auer2018quality}),
have often $(i)$ abstracted from affordability constraints, and $(ii)$
do not include a non-traded service sector.\footnote{For instance, in the seminal work of \citet{Fajgelbaum2011}, IQS
is driven solely by non-homothetic preferences for quality, through
the home-market effect. Another closely related framework by \citet{matsuyama2000ricardian}
emphasizes a hierarchical demand system where different sectors produce
goods of different \emph{priority} as opposed to different quality. } As a result, from the lens of these models, discreteness in itself
does not contribute to IQS or PTM. Accordingly, these models cannot
explain why IQS and PTM are more pronounced in discrete industries.

\section{\label{sec:Suggestive-Evidence}Suggestive Evidence}

I uncover four new regularities, which suggest that discrete goods
account for a significant portion of global trade; and that macro-level
trade patterns such as international quality specialization (IQS)
and pricing-to-market (PTM) occur more intensively in discrete product
categories. These regularities are documented using the publicly-available
US import and export databases compiled and updated by \citeauthor{Feenstra2002}.
The data covers 36 years of US trade, reporting the f.o.b. (free on
board) value and physical quantity of district-level imports (exports)
from (to) various countries during the 1989-2015 period. Each observation
classifies the goods based on the 10-digit Harmonized System (HS10)
classification and reports the ``unit'' in which quantity is measured.\footnote{There is a small subset of observations for which the unit of measurement
is missing. Also, during the 1990-2006 period, the HS10 codes underwent
multiple revisions. So, to concord HS10 codes across revisions, I
use the mapping developed by \citet{pierce2012concordance}.} The import-side data, for instance, reports quantity in 48 different
units, with well over 70\% of the observations reporting quantity
in terms of either ``count'' or ``kilogram.'' 

My analysis in this section is complicated by two basic challenges:
First, \emph{discreteness} may be correlated with other product characteristics
such as degree of differentiation or intent of final use. Second,
\emph{discreteness} may be correlated with hidden firm-level characteristics
that are impossible to control for in the aggregate-level US data.
To address the former issue, I use the classification by \citet{Rauch1999}
to control for product differentiation and the classification of Broad
Economic Categories (BEC) to control for the intent of final use.
To address the latter concern, I cross-check all the benchmark results
with firm-level evidence. This evidence is documented using firm-level
Colombian import and export data, which report trade statistics corresponding
to individual importing and exporting firms at the HS10 level of product
aggregation during the 2007-2013 period---see Appendix \ref{sec:The-Firm-Level-Colombian}
for a more detailed description.

\vspace{0.1in}\begin{obsv}A significant portion of foreign trade
involves discrete goods.\end{obsv} \vspace{0.1in}

I classify an HS10 category as \emph{discrete} if it involves countable
goods.\footnote{A very small minority of HS10 codes report quantity in multiple units.
These observations are dropped throughout my analysis.} As reported in Table \ref{tab: Discrete Industries}, discrete HS10
products typically belong to the ``Machinery,'' ``Electrical \&
Optical Equipment,'' and ``Transport Equipment'' industries. More
than 90\% of import transactions in these industries involve discrete
goods like passenger cars, television sets, or washer/dryers. In many
widely-used classifications, these industries are also often classified
as durable. 

In comparison, industries like ``Food'', ``Chemicals'', and ``Textiles''
consist of mostly non-discrete HS10 categories. They feature products
such as coffee beans, rubbing alcohol, or yarns that typically report
quantity in kilograms. Considering that weight is a relatively continuous
unit of measurement and to fix ideas, I refer to such HS10 categories
as \emph{continuous}. 

A possible concern with the above classification is that some HS10
product codes involve countable but inexpensive goods such as flashlights.
The low unit price in these categories can render indivisibility or
discreteness redundant. So, when documenting the subsequent regularities,
I check robustness under a more conservative classification, wherein
HS10 categories are coded as discrete if they are both $(i)$ countable
and $(ii)$ exhibit an average unit price above \$1000.\footnote{In the interest of space, Regularities 3 and 4 are documented and
reported only under the benchmark classification. However, they can
be stated as is under my more conservative classification of discrete
products.} 

\begin{table}[!tbh]
\caption{\label{tab: Discrete Industries}Discrete versus continuous industries
in data}

\vspace{-.1in}
\begin{adjustwidth}{-0.5in}{0in}
\noindent\resizebox{1.2\textwidth}{!}{
\begin{tabular}{lcccccc}
\toprule                                    

  & \multicolumn{2}{c}{\textbf{US Imports}}  & \phantom{abc}  & \multicolumn{2}{c}{\textbf{US Exports}} 
\\ 
                       \cmidrule{2-3}                                           \cmidrule{5-6} 
Industry  &  \specialcell{\% Share of \\ discrete goods}&  \specialcell{\% Share of industry \\ in total imports} & & \specialcell{\% Share of \\ discrete goods}&  \specialcell{\% Share of industry \\ in total exports}
\\
\toprule 
Machinery				&	97.8	   &	12.3		&&	96.0	&	14.2	\\

Electrical \& Optical Equipment		&	96.8    &	 13.6		&&	94.3	&	14.4	\\

                       
Transport Equipment				&	93.0	   &	19.2		&&	78.8	&	19.3	\\

N.E.C. \& Recycling				&	51.0    &	2.8		&&	52.2	&	0.8	\\ 
Rubber \& Plastic	   &	38.0	   &	 1.6		&&	24.3	&	2.1	\\ 
Minerals				 &	17.5	   &	1.1		&&	20.7	&	0.8	\\ 
Paper					&	16.6	   &	2.5		&&	27.4	&	4.3	\\ 

Textiles, Leather \& Footwear 				&	12.5	   &	10.1		&&	5.0	&	3.4	\\

Wood					 &	9.3	   &	 1.5		&&	7.1	&	0.9	\\
Basic \& Fabricated Metals		     &	6.7	&	6.6	      	&&	6.2	&	6.4	\\
Agriculture \& Mining  					&	2.3	&    13.5      	&&    2.4	&	8.7    \\ 
Chemicals									&	2.1	& 	7.9	        &&	1.6	&	15.6	\\  Food										 &	0.0	   &	4.0		&&	1.0	&	6.7	\\  
Petroleum				&	0.0	   &	 3.4		&&	0.0	&	2.3	\\ 

 Total					&	47.9	   &	100.0 	&&	45.1	&	100.0 \\
\bottomrule
\end{tabular}}
\end{adjustwidth}	

\noindent\begin{minipage}[t]{1\columnwidth}%
\begin{spacing}{0.8}
\emph{\footnotesize{}Note}{\footnotesize{}: All percentage shares
are value-weighted. The sectoral classification is from the WIOD.
Observations that report quantity in units of ``count'' are classified
as discrete.}
\end{spacing}
%
\end{minipage}
\end{table}
\vspace{0.1in} \begin{obsv}
There is significantly more international price heterogeneity within discrete product categories.
\end{obsv} \vspace{0.1in}

I establish the above regularity using \emph{quantity-weighted} unit
prices, which I denote by $\bar{p}_{ji,kt}$ (origin country $j$--US
district $i$--HS10 product $k$--year $t$). Table \ref{tab:Price-Heterogeneity}
reports basic statistics describing the cross-national heterogeneity
in $\bar{p}_{ji,kt}$ within \emph{HS10 product$\times$district$\times$year}
cells. For discrete product categories, the median cell exhibits a
75/25 percentile price spread of $12.0$. For \emph{continuous} product
categories, however, the median cell exhibits a 75/25 price spread
of only $2.1$. The difference is even greater if we adopt a more
conservative definition of discreteness, whereby product categories
are classified as discrete if they are not only countable but also
exhibit an average price per good of at least $\$1000$. In that case,
the median 75/25 price spread stands at $29.3$. The same exact ranking
is borne out if we use the within-product coefficient of variation
as our measure of price heterogeneity.

The lower panel in Table \ref{tab:Price-Heterogeneity} demonstrates
that similar patterns hold even after we control for the importing
firm's characteristics. This panel reports the (within cell $ikt$)
heterogeneity in $\bar{p}_{ji,kt}$, where $i$ now indexes a Colombian
importing firm. To elaborate, firm $i$ may purchase HS10 product
$k$ in year $t$ from various origin countries indexed by $j$. The
discrete products purchased by the same firm in the same year exhibit
significantly more cross-supplier price heterogeneity than non-discrete
products. 

\begin{table}
\caption{\label{tab:Price-Heterogeneity}Price Heterogeneity: \emph{Discrete
vs. Continuous HS10 Categories }}

\vspace{-.15in}
\begin{center}
\begin{adjustwidth}{0in}{0in}
\noindent\resizebox{1\textwidth}{!}{
\begin{tabular}{lcccccc}
\addlinespace
\toprule                                    

  & & \multicolumn{2}{c}{Discrete}  & \phantom{abc}  & \multirow{1}{*}{Continuous} 
\\ 
                       \cmidrule{3-4}                                           
     & & {\small All} & {\small $P_{avg}$>\$1000}  & & 
\\
\toprule   
         
          & &   \multicolumn{4}{c}{\textbf{District-Level Data}  (\emph{United States})}  \\
\cmidrule{3-6}                                   

  75/25 price spread (median)     & & 12.00 & 29.32           && 2.06  \\                        
\addlinespace
 CV of unit price (median)            & & 221\% & 267\%            && 131\% \\
\addlinespace
\addlinespace

Observations                    & &  9,069,000   & 3,613,000    &&    7,520,000 \\
\addlinespace
\toprule

& &   \multicolumn{4}{c}{\textbf{Firm-Level Data}  (\emph{Colombia})}  \\
\cmidrule{3-6}                                   

  75/25 price spread (median)     & & 4.35 & 4.88           &&  1.16  \\                        
\addlinespace
 CV of unit price (median)            & & 89\% & 93\%            && 10\% \\
\addlinespace
\addlinespace
bservations                    & &  2,407,000   & 1,304,000    &&    881,000 \\

\bottomrule
\end{tabular}}
\noindent\begin{minipage}[t]{1\columnwidth} \begin{spacing}{0.8} {\footnotesize \emph{Note}: All statistics are value-weighted. The statistics in the top panel are documented using district-level US import data from 1989 to 2015. The statistics in the bottom panel are documented using firm-level Colombian import data from 2007 to 2013. Continuous goods correspond to observations that report quantity in unit other than count.} \end{spacing} \end{minipage}
\end{adjustwidth}
\end{center}
\end{table}
Trade economists generally attribute such within-product price heterogeneity
to international quality and markup differences (\citet{Schott2004,Hallak2006,Baldwin,Hallak2011}).
Under this interpretation, Regularity 2 suggests that there is significantly
more international quality and markup heterogeneity in discrete product
categories. The following two regularities confirm that the greater
quality/markup heterogeneity in discrete product categories has a
systematic nature, and may be due to the greater degrees of IQS and
PTM.

\vspace{0.1in} \begin{obsv}
North-South quality specialization is more pronounced in discrete product categories.
\end{obsv} \vspace{0.1in}

To analyze quality specialization, I adopt the perspective in \citet{Schott2004}
and \citet{Hummels2005}, whereby the cross-national variation in
export prices to the same market is attributed to quality differences.\footnote{Recent theoretical advances have corroborated this interpretation.
Most notably, \citet{arkolakiselusive} show that (if the firm-level
productivity distribution is Pareto), the distribution of markups
from country $j$ to $i$ is invariant to country $j$'s characteristics.
That being the case, the cross-exporter variation in prices can be
entirely attributed to export quality differences. } Correspondingly, I define the degree of North-South quality specialization
as $\frac{\partial\ln\textrm{export unit price}}{\partial\ln\textrm{origin's GDP p/c}}$.
To determine how product discreteness influences North-South quality
specialization, I run the following regression on the US import data:\footnote{\citet{manova2012export} adopt a similar estimating equation to analyze
the within-product heterogeneity in unit prices, but with firm-level
data.}
\begin{equation}
\log\bar{p}_{ji,kt}=\left(\beta_{w}+\beta_{w}^{\mathcal{D}}\times\mathbbm{1}\left\{ \textrm{Discrete}\right\} \right)\log\textrm{GDP p/c}_{j,t}+\textrm{Controls}_{j,t}+\delta_{ikt}+\epsilon_{jikt},\label{eq: QS}
\end{equation}
The dependent variable, $\bar{p}_{ji,kt}$ denotes the unit price
of country $j$'s exports to US district $i$ in HS10 product category
$k$, in year $t$. On the right-hand side, $\textrm{GDP p/c}_{j,t}$
controls for the \emph{origin }country's GDP per capita; $\mathbbm{1}\left\{ \textrm{Discrete}\right\} $
is a dummy variable that assumes a value of $1$ if product category
$k$ is discrete and is zero otherwise; $\textrm{Controls}_{j,t}$
is composed of a set of standard gravity controls like total GDP,
geo-distance, and Free Trade Agreements; and finally, $\delta_{ikt}$
controls for district-product-year fixed effects. 

In addition to estimating Equation \ref{eq: QS} on the entire sample,
I perform the same estimation while restricting the sample to $(i)$
only final consumption goods, and $(ii)$ only differentiated goods
as classified by \citet{Rauch1999}. Doing so removes the possible
concern that discreteness is correlated with other characteristics
like the intended use or the degree of product differentiation. Also,
since each observation in my country-level data represents a quantity-weighted
average of firm-level variables, the above estimation is subject to
the classic heteroskedasticity problem. To attain consistent estimates,
I follow \citet{kmenta1997elements} and weigh each observation by
the underlying trade quantity.\footnote{When handling grouped data, the weighted estimator where each observation
is weighted by the square root of group size is BLUE. Here, trade
quantity serves as a theory-consistent proxy for group size. Adopting
such a weighting scheme delivers consistent estimates even if group
size is correlated with $\epsilon_{jikt}$---see \citet{solon2015we}.} To account for sample selection, I use the two-step correction developed
by \citet{wooldridge1995selection}.\footnote{Specifically, I estimate an analog of Regression \ref{eq: QS} by
Tobit, but with trade values as the dependent variable. The residual
from this first regression is then used as an additional control when
estimating Equation \ref{eq: QS}---see \citet{harrigan2015export}
for further details.}

In the above regression, the coefficient $\beta_{w}$ reflects the
baseline degree of North-South quality specialization, while $\beta_{w}^{\mathcal{D}}$
reflects the extent to which this kind of specialization is more pronounced
in discrete product categories. The estimation results reported in
Table \ref{tab:Quality-Specialization} point to a positive and statistically
significant $\beta_{w}^{\mathcal{D}}$, indicating that IQS is noticeably
more pronounced in discrete product categories. This remains to be
true if we confine our sample to only final consumption goods or differentiated
goods.\footnote{The Weighted OLS estimator may not correct for the overrepresentation
of observations with extreme values of the explanatory variables (\citet{solon2015we}).
Nonetheless, dropping observations with explanatory variables that
fall above/below the 95/5 percentile does not qualitatively change
any of the conclusions deduced from my analysis. }

\begin{table}[htb]\centering \def\sym#1{\ifmmode^{#1}\else\(^{#1}\)\fi} 
\begin{adjustwidth}{-0.25in}{-0in}
\caption{\label{tab:Quality-Specialization} Discreteness and North-South Quality Specialization (\emph{dependent}: $\log \bar{p}_{j,ikt}$)}
\vspace{-.1in}
\noindent\resizebox{1.1\textwidth}{!}{
\begin{tabular}{l*{6}{D{.}{.}{-1}}} \toprule                    
& \multicolumn{2}{c}{All Goods} & \multicolumn{2}{c}{Final Goods} & \multicolumn{2}{c}{Differentiated Goods} \\ 
\toprule   
         
          &    \multicolumn{6}{c}{\textbf{District-Level Data}  (\emph{United States})}  \\
\cmidrule{2-7}

$\log\textrm{GDP p/c}_{j,t}$              &      0.14\sym{***}&        0.02\sym{***}&        0.11\sym{***}&        0.09\sym{***}&        0.13\sym{***}&        0.01         \\                     &     (0.010)         &     (0.009)         &     (0.007)         &     (0.010)         &     (0.010)         &     (0.015)         \\ \addlinespace 
$\log\textrm{GDP p/c}_{j,t}$$\times$$\textrm{Discrete}_{k}$    
&     ...                &        0.18\sym{***}&       ...              &        0.04\sym{***}&  ...                     &        0.17\sym{***}\\                     &                     &     (0.014)         &                     &     (0.011)         &                     &     (0.017)         \\ \\

Observations  &     \multicolumn{2}{c}{16,636,000}         &     \multicolumn{2}{c}{6,647,000}         & \multicolumn{2}{c}{13,107,000} \\
Fixed Effects  &     \multicolumn{6}{c}{importing district$\times$HS10 prdouct$\times$year} 
         
\\    \addlinespace
\toprule   
         
          &    \multicolumn{6}{c}{\textbf{Firm-Level Data}  (\emph{Colombia})}  \\
\cmidrule{2-7}
$\log\textrm{GDP p/c}_{j,t}$                  &        0.16\sym{***}&        0.12\sym{***}&        0.16\sym{***}&        0.12\sym{***}&        0.17\sym{***}&        0.13\sym{***}\\                     &     (0.010)         &     (0.008)         &     (0.010)         &     (0.008)         &     (0.010)         &     (0.009)         \\ \addlinespace
$\log\textrm{GDP p/c}_{j,t}$$\times$$\textrm{Discrete}_{k}$    & ... &        0.05\sym{***}& ...                     &        0.06\sym{***}&          ...           &        0.04\sym{***}\\                     &                     &     (0.007)         &                     &     (0.008)         &                     &     (0.008)         \\  \\

Observations  &     \multicolumn{2}{c}{3,288,000 }         &     \multicolumn{2}{c}{ 2,533,000  }         & \multicolumn{2}{c}{2,474,000} \\
Fixed Effects  &     \multicolumn{6}{c}{importing firm$\times$HS10 prdouct$\times$year} \\
\bottomrule
\end{tabular}}
\noindent\begin{minipage}[t]{1.1\columnwidth} \begin{spacing}{0.8} {\footnotesize  \emph{Note}: The estimating equation is Equation \ref{eq: QS}. The top panel is produced using district-level US import data from 1989 to 2015. The bottom panel is produced using firm-level Colombian import data from 2007 to 2013. All estimations are conducted with additional  controls for total GDP, distance, and free trade agreements. Each observation is weighted by the corresponding trade quantity.  Standard errors are reported in parentheses and clustered by exporting country and year. \sym{**} \(p<0.05\), \sym{***} \(p<0.01\)} \end{spacing} \end{minipage} 
\end{adjustwidth}
\end{table} 

The bottom panel in Table \ref{tab:Quality-Specialization} documents
the same pattern while controlling for firm-level characteristics.
The results in this panel correspond to an estimation of Equation
\ref{eq: QS} on the firm-level Colombian import data, in which case
index $i$ indexes an importing firm rather than a district. These
results indicate that IQS is more pronounced in discrete product categories
even after we control for \emph{importing firm-product-year }fixed
effects and divide goods by degree of differentiation and intent of
final use.

The above claims rest on an implicitly-assumed link between quality
and unit price, which can be complicated by cross-national productivity
differences. Specifically, the results in Table \ref{tab:Quality-Specialization}
can be alternatively interpreted as high-income countries having a
comparative cost advantage in continuous industries. I discuss this
issue more elaborately in Section \ref{sec:Empirical-Evidence}, where
I analyze the role of discreteness while controlling for cross-country
productivity differences as well as other confounding factors. 

\vspace{0.1in} \begin{obsv}
Pricing-to-market is more pronounced in discrete product categories.
\end{obsv} \vspace{0.1in}

To analyze PTM, I adopt the perspective in \citet{alessandria2011pricing}
and \citet{simonovska2015income}, whereby the cross-market variation
in prices from the same exporter is attributed to markup differences.\footnote{Alternatively, the same pattern can be attribute to firm varying their
quality across markets. In that case, Regularity 4 can be stated as
firms vary their product quality more intensively

discrete industries. } Accordingly, I define the degree of PTM as $\frac{\partial\ln\textrm{export unit price}}{\partial\ln\textrm{destination's GDP p/c}}$.
To determine how discreteness influences PTM, I run the following
regression on the US export data:
\begin{align}
\log\bar{p}_{ji,kt} & =\left(\tilde{\beta}_{w}+\tilde{\beta}_{w}^{\mathcal{D}}\times\mathbbm{1}\left\{ \textrm{Discrete}\right\} \right)\log\textrm{GDP p/c}_{i,t}+\textrm{Controls}_{i,t}+\delta_{jkt}+\epsilon_{jikt}\label{eq: PTM}
\end{align}
The dependent variable, $\bar{p}_{ji,kt}$, denotes the unit price
of US district $j$'s exports to country $i$ in HS10 category $k$,
in year $t$. On the right-hand side, $\textrm{GDP p/c}_{i,t}$ controls
for the \emph{destination} country's GDP per capita, $\mathbbm{1}\left\{ \textrm{Discrete}\right\} $
is a discrete product dummy; $\textrm{Controls}_{i,t}$ is composed
of a set of standard gravity controls like total GDP, geo-distance,
and Free Trade Agreements; and finally, $\delta_{jkt}$ controls for
district- product-year fixed effects. Also, as before, I estimate
Equation \ref{eq: PTM} separately for all goods as well the subset
of final consumption goods and differentiated goods.

In the above regression, $\tilde{\beta}_{w}$ reflects the baseline
degree of PTM, while $\tilde{\beta}_{w}^{\mathcal{D}}$ reflects the
extent to which PTM is more pronounced in discrete product categories.
The estimation results reported in Table \ref{tab:Pricing-to-Market}
point to a positive and statistically significant $\tilde{\beta}_{w}^{\mathcal{D}}$,
indicating that PTM is in fact noticeably more pronounced in discrete
product categories. This also remains to be true if we confine our
sample to only final consumption goods or differentiated goods.

\begin{table}[htb]\centering \def\sym#1{\ifmmode^{#1}\else\(^{#1}\)\fi} 
\begin{adjustwidth}{-0.25in}{-0in}
\caption{\label{tab:Pricing-to-Market} Discreteness and Pricing-to-market (\emph{dependent}: $\log \bar{p}_{ji,kt}$)}
\vspace{-.1in}
\noindent\resizebox{1.1\textwidth}{!}{
\begin{tabular}{l*{6}{D{.}{.}{-1}}} \toprule                    
& \multicolumn{2}{c}{All Goods} & \multicolumn{2}{c}{Final Goods} & \multicolumn{2}{c}{Differentiated Goods} \\ 
\toprule   
         
          &    \multicolumn{6}{c}{\textbf{District-Level Data}  (\emph{United States})}  \\
\cmidrule{2-7}

$\log\textrm{GDP p/c}_{i,t}$       &          0.04\sym{***}&        0.02\sym{***}&        0.05\sym{***}&        0.01\sym{***}&        0.09\sym{***}&        0.05\sym{***}\\                     &     (0.007)         &     (0.004)         &     (0.007)         &     (0.004)         &     (0.015)         &     (0.011)         \\ \addlinespace 
$\log\textrm{GDP p/c}_{i,t}$$\times$$\textrm{Discrete}_{k}$    
&            ...         &        0.06\sym{***}&        ...             &        0.13\sym{***}& ...                     &        0.06\sym{***}\\                     &                     &     (0.018)         &                     &     (0.026)         &                     &     (0.024)         \\ \\
Observations  &     \multicolumn{2}{c}{7,855,000 }         &     \multicolumn{2}{c}{1,663,00 }         & \multicolumn{2}{c}{5,832,608} \\
Fixed Effects  &     \multicolumn{6}{c}{exporting district$\times$HS10 prdouct$\times$year} 
         
\\    \addlinespace
\toprule   
         
          &    \multicolumn{6}{c}{\textbf{Firm-Level Data}  (\emph{Colombia})}  \\
\cmidrule{2-7}
$\log\textrm{GDP p/c}_{i,t}$                  
&        0.04\sym{***}&        0.02\sym{***}&        0.05\sym{***}&        0.01\sym{***}&        0.09\sym{***}&        0.05\sym{***}\\                     &     (0.007)         &     (0.004)         &     (0.007)         &     (0.004)         &     (0.015)         &     (0.011)         \\ \addlinespace
$\log\textrm{GDP p/c}_{i,t}$$\times$$\textrm{Discrete}_{k}$    
&    ...                 &        0.06\sym{***}&           ...          &        0.13\sym{***}& ...                     &        0.06\sym{***}\\                     &                     &     (0.018)         &                     &     (0.026)         &                     &     (0.024)         \\  \\

Observations  &     \multicolumn{2}{c}{668,000 }         &     \multicolumn{2}{c}{ 387,000  }         & \multicolumn{2}{c}{567,000} \\
Fixed Effects  &     \multicolumn{6}{c}{exporting firm$\times$HS10 prdouct$\times$year} \\
\bottomrule
\end{tabular}}
\noindent\begin{minipage}[t]{1.1\columnwidth} \begin{spacing}{0.8} {\footnotesize  \emph{Note}: The estimating equation is Equation \ref{eq: PTM}. The top panel is produced using district-level US export data from 1989 to 2015. The bottom panel is produced using firm-level Colombian export data from 2007 to 2013. All estimations are conducted with additional  controls for total GDP, distance, and free trade agreements. Each observation is weighted by the corresponding trade quantity.  Standard errors are reported in parentheses and clustered by importing country and year. \sym{**} \(p<0.05\), \sym{***} \(p<0.01\)} \end{spacing} \end{minipage} 
\end{adjustwidth}
\end{table} 

The bottom panel in Table \ref{tab:Pricing-to-Market} documents the
same pattern while controlling for firm-level characteristics. The
results in this panel are derived by estimating Equation \ref{eq: PTM}
on the firm-level Colombian export data, in which case index $j$
identifies an exporting firm rather than a district.\footnote{In the Colombia sample numerous firms supply both discrete and continuous
goods. To give an example, firm \#900500347 exports both durometers
(HS9024800000) that are discrete and priced over \$500 per unit as
well as continuous aluminum sheets (HS7616999000).} They indicate that PTM is more pronounced in discrete product categories
even after we control for \emph{exporting firm-product-year }fixed
effects and divide goods by degree of differentiation and intent of
final use.

To take stock, the evidence presented in this section indicate that
IQS and PTM are more pronounced in discrete industries. One possible
reason for this is that the standard drivers of IQS and PTM operate
more forcefully in discrete industries. Another possibility is that
\emph{discreteness} gives rise to alternative driving forces behind
IQS and PTM that are otherwise irrelevant. In what follows, I present
a theoretical framework that highlights two such alternative forces.

\section{\label{sec:Theoretical-Framework}Theoretical Framework}

My main objective in this section is to highlight two alternative
drivers of IQS and PTM in discrete industries. Considering this objective,
my theory purposely abstracts from \emph{taste-driven} or \emph{cost-driven}
IQS and PTM, which have little to do with discreteness. Accordingly,
I model continuous industries in a rather standard fashion. My modeling
of discrete industries and the insights that derive from it, however,
rests on three distinctive assumptions:
\begin{enumerate}
\item Discrete goods are indivisible: purchasing them requires paying the
price of one full unit, which is akin to paying a fixed acquisition
cost.
\item The consumption flow derived from a personally-acquired discrete good
can be substituted with analogous services---e.g., consumers can
forgo buying a truck for moving purposes by renting a moving service.\footnote{This feature can be attributed to the strong overlap between discreteness
and durability (see Section \ref{sec:Suggestive-Evidence}). Unlike
perishable consumption goods, durable consumption goods are non-rival.
That is, one consumer's usage of the good does not exclude others
from using it later. Accordingly, one can use a taxi service without
(effectively) paying for the car but by only paying for the labor
and the perishable fuel used in providing the service. This feature
of durable goods has been widely-invoked in macro-economics literature---see
\citet{benhabib1991homework}, \citet{cooley1995economic}, and \citet{fernandez2011consumption}.}
\item Quality and quantity are not isomorphic. In particular, from the perspective
of a firm, supplying higher quality is different from supplying more
quantity. The former requires incurring a sunk R\&D cost, whereas
the latter only requires employing more inputs without necessarily
altering one's production technology.
\end{enumerate}
For the sake of composition, I also make various parametric assumptions
throughout this section; but these assumptions are less consequential
to the theoretical results that follow---Appendix \ref{sec: Non-Parametric-Model}
shows that the main predictions of the model hold equally in a more
general (less-parametric) environment. 

\subsection{Economic Environment}

\paragraph*{\emph{Countries.}}

The global economy consists of $i=1,...,I$ countries, with $\mathbb{C}$
denoting the set of countries. Country $i$ is populated by $L_{i}$
individuals, each of whom is endowed with $y\in\mathbb{R_{+}}$ units
of effective labor that is the sole factor of production. All individuals
are perfectly mobile across the production of different goods but
are immobile across countries; and are paid an economy-wide wage,
$w_{i}$, per unit of effective labor. 

\paragraph*{\emph{Income Distribution.}}

I assume that the distribution of effective labor, $y$, across individuals
within country $i$ is given by
\[
G_{i}(y)=1-e^{-\zeta_{i}\left(y-\bar{y}_{i}\right)}.
\]
Since each individual residing in country $i$ collects an income
equal to $yw_{i}$, then $G_{i}\left(.\right)$ also governs the distribution
of income in that country. The exponential parametrization of the
income distribution is motivated by evidence in \citet{druagulescu2001evidence}
and \citet{gabaix2009power}. I allow parameters $\zeta_{i}$ and
$\bar{y}_{i}$ to vary arbitrarily across countries; but I impose
that the mean of the distribution equals one in all countries, i.e.,
$\mathbb{E}_{i}(y)=\bar{y}_{i}+\zeta_{i}=1$ for all $y$. 

\paragraph*{\emph{Industries.}}

The economy is composed of multiple industries (indexed by $k$) that
fall under three general categories:
\begin{enumerate}
\item Continuous industries, which involve traded and \emph{infinitely-divisible}
goods, with $\mathbb{K}_{\mathcal{C}}$ denoting the set of all continuous
industries
\item Discrete industries, which involve traded but \emph{indivisible} goods,
with $\mathbb{K}_{\mathcal{D}}$ denoting the set of all discrete
industries.
\item Non-traded service industries that supply a homogeneous output, which
serves as a substitute for discrete goods. 
\end{enumerate}
I hereafter use the \emph{tilde} notation to distinguish variables
corresponding to the service sector. For instance, based on this choice
of notation, $\tilde{p}_{i,k}$ denotes the price of industry $k$-related
services in market $i$.

\paragraph*{\emph{Product Space.}}

Traded goods (whether continuous or discrete) are horizontally and
vertically differentiated. That is, they come in different firm-level
varieties (indexed by $\omega$) and different quality tiers (indexed
by $\varphi$). The set of admissible quality tiers in each industry
is predetermined and denoted by $\varPhi_{k}$. As in \citet{Fajgelbaum2011},
firms sort into (only one) of these predetermined quality tiers and
supply their differentiated variety in that particular quality level. 

In such a setup, the firm index $\omega$ uniquely identifies not
only the supplying firm but also the quality of the variety supplied
by that firm. To be more specific, let $\Omega_{j,k}(\varphi)$ denote
the set of all firms supplying quality $\varphi$ from country $j$
in industry $k$. Then, if $\omega\in\Omega_{j,k}(\varphi)$ it is
automatically implied that $\omega\notin\Omega_{j,k}(\varphi')$ for
any $\varphi'\neq\varphi$. In that regard, the index $\omega$ implicitly
conveys full information about quality as well as industry and country
of origin.

\subsection{Demand and Supply in Continuous Industries}

I start by describing the supply and demand structure in continuous
industries. Given its standard underpinning, the continuous model
can also serve as a useful benchmark to better understand my less-standard
model of discrete industries.\vspace{-0.15in}

\paragraph*{\emph{Demand.}}

The demand structure is similar to \citet{Krugman1980}, except that
there are multiple quality tiers to choose from. In particular, each
individual spends a constant fraction, $e_{k}$, of their income on
industry $k$ goods.\footnote{To be specific, the welfare of individual $\iota$ is a Cobb-Douglas
aggregate of the industry specific utility levels. Namely, $W_{\iota}=\prod_{k\in\mathbb{K}}U_{\iota,k}^{e_{k}}$.} The budget allocated to \emph{continuous} industry $k\in\mathbb{K}_{\mathcal{C}}$
can be spent on buying a fractional quantity of various firm-level
varieties. An individual's utility from consuming a basket $\boldsymbol{q}=\{q_{ji,k}(\omega)\}$
of firm-level varieties is given by a nested-CES function,\footnote{Implicit in the above formulation is that markets are segmented. Without
this assumption, PTM and the possibility of international arbitrage
may lead to zero trade between countries---see \citet{foellmi2017international}.}
\[
U_{k}\left(\boldsymbol{q}\right)=\left(\sum_{\varphi\in\varPhi_{k}}\left[\sum_{j\in\mathbb{C}}\sum_{\omega\in\Omega_{j,k}(\varphi)}\varphi^{\frac{1}{1+\theta_{k}}}\ q_{ji,k}(\omega)^{\frac{\theta_{k}}{1+\theta_{k}}}\right]^{\frac{1+\theta_{k}}{\theta_{k}}\frac{\gamma_{k}}{1+\gamma_{k}}}\right)^{\frac{1+\gamma_{k}}{\gamma_{k}}},\ \ \ k\in\mathbb{K}_{\mathcal{C}}
\]
where $\gamma_{k}$ and $\theta_{k}>0$ reflect the degree of within-
and cross-quality substitutability between different firm-level varieties.
Utility maximization subject to the budget constraints, $\sum_{\varphi}\sum_{j}\sum_{\omega}p_{ji,k}(\omega)q_{ji,k}(\omega)\leq e_{k}w_{i}y_{\iota}$
implies that the share of spending on variety $\omega$ is independent
of personal income and given by,
\begin{equation}
\lambda_{ji,k}(\omega)=\frac{p_{ji,k}(\omega)^{-\theta_{k}}}{\sum_{\ell\in\mathbb{C}}\sum_{\omega'\in\Omega_{\ell,k}(\varphi)}p_{\ell i,k}(\omega')^{-\theta_{k}}}\times e_{i,k}(\varphi),\ \ \ \forall\omega\in\Omega_{j,k}(\varphi)\label{eq: lambda cont.}
\end{equation}
where $e_{i,k}(\varphi)$ is the total share of spending on quality
$\varphi$:
\[
e_{i,k}(\varphi)=\frac{\left(\varphi\sum_{\ell\in\mathbb{C}}\sum_{\omega\in\Omega_{\ell,k}(\varphi)}p_{\ell i,k}(\omega)^{-\theta_{k}}\right)^{\gamma_{k}/\theta_{k}}}{\sum_{\varphi'\in\varPhi_{k}}\left(\varphi'\sum_{\ell\in\mathbb{C}}\sum_{\omega\in\Omega_{\ell,k}(\varphi')}p_{\ell i,k}(\omega)^{-\theta_{k}}\right)^{\gamma_{k}/\theta_{k}}},\ \ \ \forall\varphi\in\varPhi_{k}.
\]
Accordingly, total demand for variety $\omega$ in market $i$ can
be expressed as,
\begin{equation}
q_{ji,k}(\omega)=\lambda_{ji,k}(\omega)e_{k}w_{i}L_{i}/p_{ji,k}(\omega),\label{eq: Demand Contin.}
\end{equation}
where $w_{i}L_{i}$ is total income in market $i$---the fact that
total income equals $w_{i}L_{i}$ derives from my earlier choice of
normalization, $\mathbb{E}_{i}(y)=1$.\vspace{-0.15in}

\paragraph*{\emph{Supply.}}

The supply-side of the economy is similar to \citeauthor{Fajgelbaum2011}'s
\citeyearpar{Fajgelbaum2011} extension of \citet{Krugman1980}. That
is, labor is the sole factor of production. Firms supplying goods
of the same quality (in industry $k$) are symmetric within each country,
and compete under monopolistic competition--as shown in Appendix
\ref{sec: Discrete-Trade-with-Firm-Het} the model can be easily extended
to account for firm heterogeneity. A typical firm $\omega$ supplying
quality $\varphi$ from country $j$ has to incur a sunk entry cost,
$w_{j}f_{k}^{e}(\varphi)$. It also incurs a marginal cost for each
unit of output produced and transported to market $i$ that is given
by
\[
c_{ji,k}(\omega)=\tau_{ji,k}a_{j,k}c(\varphi)w_{j}.\ \ \ \omega\in\Omega_{j,k}(\omega)
\]
In the above cost specification $\tau_{ji,k}$ denotes the iceberg
transport cost; $a_{j,k}$ is a quality-neutral (\emph{country$\times$industry}-specific)
cost shifter; and $c(\varphi)$ is a universal cost shifter that accounts
for high-quality output being more costly to produce in every country.
An important feature of the above parametrization is that it precludes
any Ricardian comparative advantage along the quality margin.

The variable profit collected by a typical firm $\omega$ is, thus,
given by
\[
\pi_{ji,k}(\omega)=\left[p_{ji,k}(\omega)-c_{ji,k}(\omega)\right]q_{ji,k}(\omega),
\]
where $q_{ji,k}(\omega)$ is determined by Equation \ref{eq: Demand Contin.}.
Firms choose their optimum price, $p_{ji,k}(\omega)$, by maximizing
$\pi_{ji,k}(\omega)$, which implies a standard monopoly markup over
marginal cost. The number of firms supplying quality $\varphi$ from
country $j$ (namely, $N_{j,k}(\varphi)$) is then determined by the
\emph{free-entry} condition, wherein firms enter until net profits
are drawn to zero in every quality tier of industry $k$:
\[
\sum_{i\in\mathbb{C}}\pi_{ji,k}(\omega)=w_{j}f_{k}^{e}(\varphi),\ \ \ \forall\omega\in\Omega_{j,k}(\varphi).
\]


\subsection{Demand and Supply in Discrete Industries}

Discrete industries involve goods that are \emph{indivisible} and
substitutable with analogous services. I use a discrete choice framework
to model consumer demand in these industries. On the supply-side,
though, I assume that discrete industries are identical to continuous
counterparts.\vspace{-0.15in}

\paragraph*{\emph{Demand}.}

To draw consumption utility from a discrete variety $\omega\in\Omega_{j,k}(\varphi)$,
consumers have to pay the unit price, $p_{ji,k}(\omega)$, which is
akin to paying a fixed acquisition cost. After the good is acquired,
it delivers a flow of consumption quantity, $q_{k}(\varphi)$, that
is increasing in the good's quality, $\varphi$. The flow of personal
consumption provided by a discrete good (like a car) is substitutable
with analogous services (like a taxi service). Consumers can, therefore,
circumvent paying the acquisition cost altogether, or they can purchase
a cheaper (lower-$\varphi$) discrete variety and complement it with
a range of specialized services.

Suppose individual $\iota$ purchases discrete variety $\omega\in\Omega_{j,k}(\varphi)$
(a firm-level variety of quality $\varphi$ originating from country
$j$) for personal use and combines it with $\tilde{q}_{k}$ hours
of service. The utility they derive from this choice is given by
\[
U_{\iota}(\omega,\tilde{q})=\tilde{q}_{k}+q_{k}(\varphi)+\xi_{\iota}(\omega),\ \ \ \omega\in\Omega_{j,k}(\varphi)
\]
where $q_{k}(\varphi)$ (as defined earlier) denotes the consumption
quantity derived from a personally-acquired unit of good $\omega$;
and $\xi^{\iota}(\omega)$ is an \emph{individual$\times$variety}-specific
utility shifter that accounts for individual $\iota$'s personal taste
for variety $\omega$. Implicit in the above parametrization is the
fact that the consumption derived from a discrete good is perfectly
substitutable with analogous services. Without loss of generality,
I henceforth impose that $q_{k}(\varphi)=\varphi$.

As is common in the literature, I assume that the vector $\boldsymbol{\xi}_{\iota,k}\equiv\left\{ \xi_{\iota}(\omega)\right\} $,
which describes individual $\iota$'s taste with respect to various
industry $k$ varieties, is drawn independently from a General Extreme
Value distribution: $F_{k}(\boldsymbol{\xi}_{\iota,k})=\exp\left(-\sum_{\varphi\in\varPhi_{k}}\left(\sum_{\omega}\exp(\xi_{\iota}(\omega))^{-\theta_{k}}\right)^{\gamma_{k}/\theta_{k}}\right)$,
where $\theta_{k}\geq\gamma_{k}\geq0$ are parameters governing the
heterogeneity of consumer taste within and across quality levels in
industry $k$. This assumption allows me to obtain closed-form demand
functions.

An individual (with endowment $y$) maximizes their utility by choosing
their preferred variety, $\omega$, within industry $k$ subject to
a personal budget constraint. Recall that a fraction $e_{k}$ of an
individual's income is spent on industry $k$ goods/services (e.g.,
transportation goods and services). The budget constraint imposes
that the price of variety $\omega$ not exceed one's budget, $e_{k}w_{i}y$.
Moreover, given the choice of $\omega$, the budget constraint implicitly
determines their choice of service consumption, $\tilde{q}_{k}$.
That is, if an individual chooses variety $\omega\in\Omega_{j,k}(\varphi)$,
their budget constraint,
\[
\widetilde{p}_{i,k}\widetilde{q}_{k}+p_{ji,k}(\omega)=e_{k}w_{i}y,
\]
implicitly pins down $\tilde{q}_{k}=\left[e_{k}w_{i}y-p_{ji,k}(\omega)\right]/\tilde{p}_{i,k}$.
Taking into account the individual's budget constraint and appealing
to the Theorem of Extreme Value, utility maximization then implies
that the share of individuals with endowment $y$ who choose variety
$\omega$ is given by
\begin{equation}
\lambda_{ji,k}(\omega;y)=\frac{\exp\left(p_{ji,k}(\omega)\right)^{-\theta_{k}/\tilde{p}_{i,k}}}{\sum_{\ell\in\mathbb{C}}\sum_{\omega'\in\Omega_{\ell,k}(\varphi,y)}\exp\left(-p_{\ell i,k}(\omega')\right)^{-\theta_{k}/\tilde{p}_{i,k}}}\times e_{i,k}(\varphi;y),\label{eq: lambda disc.}
\end{equation}
if $y\geq\frac{p_{ji,k}(\omega)}{e_{k}w_{i}}$; and is zero otherwise
(see Appendix \ref{sec:Proof-of-Lemma 1} for proof). The share of
individuals with endowment $y$ who choose a variety with quality
$\varphi$ is similarly given by
\[
e_{i,k}(\varphi;y)=\frac{\left[\varphi\ \sum_{\ell\in\mathbb{C}}\sum_{\omega'\in\Omega_{\ell,k}(\varphi,y)}\exp\left(-p_{\ell i,k}(\omega)\right)^{-\theta_{k}/\tilde{p}_{i,k}}\right]^{\frac{\gamma_{k}}{\theta_{k}}}}{\sum_{\varphi'\in\varPhi_{k}}\left[\varphi'\ \sum_{\ell\in\mathbb{C}}\sum_{\omega'\in\Omega_{\ell,k}(\varphi',y)}\exp\left(-p_{\ell i,k}(\omega)\right)^{-\theta_{k}/\tilde{p}_{i,k}}\right]^{\frac{\gamma_{k}}{\theta_{k}}}},
\]
where $\Omega_{j,k}(\varphi,y)\subset\Omega_{j,k}(\varphi)$ denotes
the subset of varieties in set $\Omega_{j,k}(\varphi)$ that are affordable
to individuals with endowment $y$. The total demand facing firm $\omega$
can, therefore, be calculated as the sum of the demand across all
individuals in the economy as follows:
\begin{equation}
q_{ji,k}(\omega)=\left(\int_{y=\frac{p_{ji,k}(\omega)}{e_{k}w_{i}}}^{\infty}\lambda_{ji,k}(\omega;y)dG_{i}(y)\right)L_{i}.\label{eq: Demand Discrete}
\end{equation}
Before moving forward let me draw a brief comparison between the demand
structure in discrete and continuous industries. In both sets of industries,
individual-level preferences are homothetic in the following sense.
If two individuals have different income levels (namely, $y'$ and
$y$) but can choose from the same set of varieties, then
\[
\lambda_{ji,k}(\varphi,y)=\lambda_{ji,k}(\varphi,y').
\]
Moreover, if the above condition is satisfied, it automatically follows
that the two individuals also exhibit the same demand for quality:
$e_{i,k}(\varphi,y)=e_{i,k}(\varphi,y')$. \vspace{-0.15in}

\paragraph*{\emph{Supply.}}

On the supply side, discrete industries are assumed to be identical
to continuous industries. That is, labor is the only factor of production
and firms compete under monopolistic competition. Each firm in country
$j$ incurs a sunk cost, $w_{j}f_{k}^{e}(\varphi)$, to enter industry
$k$ in quality tier $\varphi$. Upon entry, it also incurs a marginal
cost equal to $c_{ji,k}(\omega)=\tau_{ji,k}a_{j,k}c(\varphi)w_{j}$
in order to supply each unit of output to market $i$. Recall that
this cost structure rules out Ricardian comparative advantage along
the quality margin. Each firm sets their optimum price by maximizing
their variable profit,
\[
\pi_{ji,k}(\omega)=\left[p_{ji,k}(\omega)-c_{ji,k}(\omega)\right]q_{ji,k}(\omega),
\]
where $q_{ji,k}(\omega)$ is now given by Equation \ref{eq: Demand Discrete}.
Profit maximization simply entails a standard monopoly markup over
marginal cost. The number of firms in every quality tier of industry
$k$ is determined by the \emph{free-entry} condition, wherein firms
enter until net profits are drawn to zero: $\sum_{i\in\mathbb{C}}\pi_{ji,k}(\omega)=w_{j}f_{k}^{e}(\varphi)$. 

It is important to note that the production of quality is subject
to scale economies in all industries. That is, the average cost of
supplying quality is decreasing in total output. As a result, to circumvent
trade costs, high-quality firms agglomerate in economies where there
is a relatively large local demand for high-quality varieties, leading
to the well-known \emph{home-market} effect.

\subsection{Demand and Supply of Non-Traded Services}

The demand for non-traded services is implicitly determined by the
consumption choice in discrete industries.\footnote{This is a simplifying but innocuous assumption. One can always assume
that the economy features an extra set of non-traded service industries
that are unrelated to discrete industries. Adding this feature, as
long as we maintain the Cobb-Douglas assumption on industry-level
expenditure shares, will not change any of the results in the paper. } I also assume that the service sector is perfectly competitive. In
principle, services may employ discrete goods as an intermediate input
(e.g., taxi services use cars as an input), but I assume that due
to repeated use, the fixed cost paid for acquiring the discrete good
constitutes a negligible fraction of the total cost.\footnote{Specifically, consider an incumbent firm supplying \emph{industry
$k$-related} services in market $i$. Suppose such a firm spends
$\tilde{c}_{i,k}^{\mathcal{D}}$ on purchasing discrete input varieties
and sells in total $\tilde{Q}_{i,k}$ units of services. The implicit
assumption here is that $\tilde{c}_{i,k}^{\mathcal{D}}/\tilde{Q}_{i,k}\approx0$,
so that the average cost for this supplier is approximately equal
to the marginal cost: $\tilde{AC}_{i}=\tilde{a}_{i,k}w_{i}+\frac{\tilde{c}_{i,k}^{\mathcal{D}}}{\tilde{Q}_{i,k}}\approx\tilde{a}_{i,k}w_{i}.$} Consequently, the average cost, effectively, equals the marginal
cost, $\tilde{a}_{i,k}w_{i}$, for service suppliers.

I also assume that the marginal labor cost of supplying services is
the same across all industries, i.e., $\tilde{a}_{i,k}=\tilde{a}_{i,g}$
for all $g$ and $k\in\mathbb{K}_{\mathcal{D}}$. Altogether, these
assumptions imply that the competitive price of services in market
$i$ equals:
\[
\tilde{p}_{i,k}=\tilde{p}_{i}=\tilde{a}_{i}w_{i}\ \ \ \forall k.
\]
The above equation indicates that $\tilde{p}_{i}$ and the local wage
rate, $w_{i}$, are positively but not perfectly correlated.\footnote{The positive relation between $\tilde{p}_{i}$ and $w_{i}$ is contingent
on $\tilde{a}_{i}$'s being sufficiently homogeneous across countries.
This condition serves as the foundation of the celebrated \emph{Blasssa-Samuelson
effect} (\citet{bhagwati1984services,samuelson1994facets}). } We can actually construct many theory-consistent examples where the
price of non-traded services differ across two markets with the same
$w_{i}$. This distinction plays an important role when testing the
model, later in Section \ref{sec:Empirical-Evidence}.

\subsection{Equilibrium}

In equilibrium, each firm $\omega$ that supplies quality $\varphi$
from country $j$ in industry $k$ (i.e., $\omega\in\Omega_{j,k}$)
charges the same price and sells the same quantity of output. Considering
this, I can express all variables as a function of quality rather
than the firm-level index, $\omega$. In that case, $p_{ji,k}(\varphi)$
denotes the price charged in market $i$ by a typical firm $\omega\in\Omega_{j,k}(\varphi)$.
With this choice of notation, I now present the main equilibrium outcomes
and formally define the equilibrium. 

Following Equations \ref{eq: lambda cont.} and \ref{eq: lambda disc.},
the market share of origin $j$ varieties in country $i$ can be expressed
as\footnote{Note that in the case of continuous industries, $\lambda_{ji,k}(\varphi)$
denotes the \emph{value-share} of spending; whereas in the case of
discrete industries, it denotes the \emph{quantity-share}.}
\begin{equation}
\lambda_{ji,k}(\varphi)=\begin{cases}
\frac{N_{j,k}(\varphi)p_{ji,k}(\varphi)^{-\theta_{k}}}{\sum_{\ell\in\mathbb{C}}N_{\ell,k}(\varphi)p_{\ell i,k}(\varphi)^{-\theta_{k}}}\times e_{i,k}(\varphi;y) & \textrm{if }k\in\mathbb{K}_{\mathcal{C}}\\
\int_{y=\frac{p_{ji,k}(\varphi)}{e_{k}w_{i}}}^{\infty}\frac{N_{j,k}(\varphi)\exp\left(p_{ji,k}(\varphi)\right)^{-\frac{\theta_{k}}{\tilde{p}_{i}}}}{\sum_{\ell\in\mathbb{C}_{i}(\varphi,y)}N_{\ell,k}(\varphi)\exp\left(p_{\ell i,k}(\varphi)\right)^{-\frac{\theta_{k}}{\tilde{p}_{i}}}}\times e_{i,k}(\varphi;y)\ dG_{i}(y) & \textrm{if }k\in\mathbb{K}_{\mathcal{D}}
\end{cases}\label{eq: expenditure share}
\end{equation}
where $\mathbb{K}_{\mathcal{C}}$ denotes the \textbf{\emph{set of
continuous industries}}; $\mathbb{K}_{\mathcal{D}}$ denotes the \textbf{\emph{set
of discrete industries}}; and $\mathbb{C}_{i}(\varphi,y)\in\mathbb{C}$
denotes the subset of countries whose quality $\varphi$ varieties
are affordable to income group $y$ in market $i$. The share of quality
$\varphi$ in total consumption is, meanwhile, given by:
\[
e_{i,k}(\varphi;y)=\begin{cases}
\frac{\left(\varphi\sum_{\ell\in\mathbb{C}}N_{\ell,k}(\varphi)p_{\ell i,k}(\varphi)^{-\theta_{k}}\right)^{\gamma_{k}/\theta_{k}}}{\sum_{\varphi'\in\varPhi_{k}}\left(\varphi'\sum_{\ell\in\mathbb{C}}N_{\ell,k}(\varphi')p_{\ell i,k}(\varphi')^{-\theta_{k}}\right)^{\gamma_{k}/\theta_{k}}} & \textrm{\textrm{if }k\ensuremath{\in\mathbb{K}_{\mathcal{C}}}}\\
\frac{\left(\varphi\sum_{\ell\in\mathbb{C}(\varphi,y)}N_{\ell,k}(\varphi)\exp\left(p_{\ell i,k}(\varphi)\right)^{-\frac{\theta_{k}}{\tilde{p}_{i}}}\right)^{\gamma_{k}/\theta_{k}}}{\sum_{\varphi'\in\varPhi_{k}}\left(\varphi'\sum_{\ell\in\mathbb{C}_{i}(\varphi,y)}N_{\ell,k}(\varphi')\exp\left(p_{\ell i,k}(\varphi')\right)^{-\frac{\theta_{k}}{\tilde{p}_{i}}}\right)^{\gamma_{k}/\theta_{k}}} & \textrm{if }k\in\mathbb{K}_{\mathcal{D}}
\end{cases}.
\]
The above expressions immediately imply that the equilibrium demand
elasticity facing individual firms can be stated as (see Appendix
\ref{sec: Proof-Lemma3} for derivation):
\begin{equation}
\varepsilon_{ji,k}(\varphi)=\begin{cases}
\theta_{k}+1 & \textrm{if }k\in\mathbb{K}_{\mathcal{C}}\\
\theta_{k}\frac{p_{ji,k}\left(\varphi\right)}{\tilde{p}_{i}}+\zeta_{i}\frac{p_{ji,k}\left(\varphi\right)}{e_{k}w_{i}} & \textrm{if }k\in\mathbb{K}_{\mathcal{D}}
\end{cases}.\label{eq: Eq epsilon}
\end{equation}
The difference between the demand elasticity faced by firms in discrete
versus continuous industries plays a key role in my theory. So, I
will discuss these differences more elaborately in the following Subsection.
Finally, given the expression for $\lambda_{ji,k}(\varphi)$, aggregate
trade values, $X_{ji,k}(\varphi)\equiv N_{j,k}(\varphi)p_{ji,k}(\varphi)q_{ji,k}(\varphi)$,
are given by
\begin{equation}
X_{ji,k}\left(\varphi\right)=\begin{cases}
\lambda_{ji,k}(\varphi)e_{k}Y_{i} & \textrm{if }k\in\mathbb{K}_{\mathcal{C}}\\
\lambda_{ji,k}(\varphi)\frac{p_{ji,k}(\varphi)}{w_{i}}Y_{i} & \textrm{if }k\in\mathbb{K}_{\mathcal{D}}
\end{cases},\label{eq: EQ trade values}
\end{equation}
where $Y_{i}\equiv w_{i}L_{i}$. Equilibrium corresponds to a vector
of prices, $\boldsymbol{p}\equiv\{p_{ji,k}(\varphi)\}$, wage rates,
$\boldsymbol{w}\equiv\{w_{j}\}$, and the number of firms, $\boldsymbol{N}\equiv\{N_{j,k}(\varphi)\}$,
that satisfy $(i)$ the optimal monopoly pricing condition (MP), $(ii)$
the balanced trade condition (BT), and $(iii)$ the free entry condition
(FE):
\[
\begin{cases}
p_{ji,k}(\varphi)=\frac{\varepsilon_{ji,k}\left(\varphi\right)}{\varepsilon_{ji,k}\left(\varphi\right)-1}\tau_{ji,k}a_{j,k}c(\varphi)w_{j} & \forall\varphi\in\varPhi_{k},k\in\mathbb{K}\ \ \ (\textrm{MP})\\
\sum_{j\in\mathbb{C}}\sum_{k\in\mathbb{K}}\sum_{\varphi\in\varPhi_{k}}X_{ji,k}\left(\varphi\right)=\sum_{\ell\in\mathbb{C}}\sum_{k\in\mathbb{K}}\sum_{\varphi\in\varPhi_{k}}X_{i\ell,k}\left(\varphi\right) & \forall k\in\mathbb{K}\ \ \ (\textrm{BT})\\
N_{j,k}(\varphi)w_{j}f_{k}^{e}(\varphi)=\sum_{i\in\mathbb{C}}X_{ji,k}(\varphi)/\varepsilon_{ji,k}(\varphi) & \forall k\in\mathbb{K}\ \ \ (\textrm{FE})
\end{cases},
\]
By Walras' law, the satisfaction of the BT and FE conditions ensures
that labor markets also clear in each country.

\subsection{Discussion: \emph{The Demand Elasticity Facing Discrete Goods}}

While I assumed that individual preferences for traded goods are homothetic,
the market-level demand elasticity facing discrete goods varies systematically
across markets. This variation is driven by $(i)$ cross-market heterogeneity
in the price of non-traded services, and $(ii)$ affordability constraints
being more-or-less binding in different markets. To elaborate on these
two channels, we can decompose the demand elasticity facing discrete
goods in market $i$ as follows:

\[
\varepsilon_{ji,k}(\varphi)=\underbrace{\theta_{k}\frac{p_{ji,k}\left(\varphi\right)}{\tilde{p}_{i}}}_{\textrm{NT-driven}}\ +\ \underbrace{\zeta_{i}\frac{p_{ji,k}\left(\varphi\right)}{e_{k}w_{i}}}_{\textrm{affordability-driven}}.
\]
In the above expression, the \emph{NT-driven} component accounts for\emph{
}cross-country differences in the price of non-traded services. This
component depends on $(a)$ the degree of product differentiation
$\theta_{k}$, the role of which is well understood, as well as $(b)$
$p_{ji,k}(\omega)/\tilde{p}_{i}$, which is the amount of service
consumption one forgoes to acquire a full unit of variety $\omega$.
Importantly, the demand elasticity is lower the higher the price of
services in market $i$.

The \emph{affordability-driven} component is driven by the fact that
an incremental increase in $p_{ji,k}(\omega)$ can make variety $\omega$
prohibitively expensive for a fraction of existing consumers in market
$i$. The extent of demand loss due to this channel, depends on $(a)$
variety $\omega$'s \emph{degree of affordability} in market $i$,
$p_{ji,k}(\omega)/e_{k}w_{i}$, as well as $(b)$ the hazard rate
of the income distribution, $\zeta_{i}$.\footnote{As shown in Appendix \ref{sec: Non-Parametric-Model}, even when $G_{i}(y)$
does not exhibit a constant hazard rate, the propositions that follow
hold as long as $G_{i}(y)$ satisfies the IGFR property. } 

\section{Quality Specialization and Pricing-to-Market}

In this section, I use my theoretical model to outline two alternative
drivers of pricing-to-market (PTM) and international quality specialization
(IQS) in discrete industries. This section, therefore, sheds light
on why (as documented in Section \ref{sec:Suggestive-Evidence}) PTM
and IQS occur more intensively in discrete industries. 

\paragraph*{\emph{Pricing-to-Market (PTM).}}

Using the demand elasticity specified by Equation \ref{eq: Eq epsilon},
we can immediately calculate the optimal monopoly price for each traded
variety. Doing so, implies the following price formulation in continuous
($k\in\mathbb{K}_{\mathcal{C}}$) and discrete ($k\in\mathbb{K}_{\mathcal{D}}$)
industries:
\[
p_{ji,k}(\varphi)=\begin{cases}
\left(1+1/\theta_{k}\right)c_{ji,k}(\varphi) & \textrm{if }k\in\mathbb{K}_{\mathcal{C}}\\
c_{ji,k}(\varphi)+\vartheta_{i,k}w_{i} & \textrm{if }k\in\mathbb{K}_{\mathcal{D}}
\end{cases},
\]
with $1/\vartheta_{i,k}\equiv\frac{\theta_{k}}{\tilde{a}_{i}}+\frac{\zeta_{i}}{e_{k}}$
denoting a market-specific \emph{structural} parameter. The above
equation indicates that (absent non-homothetic preferences) there
is no scope for PTM in continuous industries. In discrete industries,
however, PTM occurs even though preferences are homothetic; as suppliers
charge systematically higher prices in markets where $(a)$ they compete
with higher-priced services; and ($b)$ their product is more affordable
to the average consumer.

These results can perhaps explain why PTM is more pronounced in discrete
industries. In both discrete and continuous industries, PTM may be
driven by non-homothetic preferences. But in discrete industries,
PTM is magnified by $(a)$ cross-national heterogeneity in the price
of non-traded services, and $(ii)$ affordability constraints; both
of which are less relevant in continuous industries.

\paragraph*{\emph{International Quality-Specialization (IQS).}}

To fix ideas, allow me to first define IQS formally from the lens
of the present model. Let $Q_{ji,k}(\varphi)\equiv N_{j,k}(\varphi)q_{ji,k}(\varphi)$
denote aggregate export quantity and suppose $\varphi_{H}$ and $\varphi_{L}$
respectively denote a \emph{high} and a \emph{low}-quality tier in
industry $k$, i.e., $\varphi_{H}>\varphi_{L}$. As a matter of definition,
an increase in country-level characteristic $x$ prompts high-quality-specialization
if and only if
\[
\frac{\partial}{\partial x_{j}/x_{\ell}}\left\{ \frac{Q_{ji,k}\left(\varphi_{H}\right)/Q_{\ell i,k}\left(\varphi_{H}\right)}{Q_{ji,k}\left(\varphi_{L}\right)/Q_{\ell i,k}\left(\varphi_{H}\right)}\right\} >0,
\]
for (nearly) all $j$, $\ell$, and $i\in\mathbb{C}$. So, to determine
patterns of IQS based on this definition, we need to derive an expression
for the relative exports of $\varphi_{H}$ to $\varphi_{L}$. Doing
so in Appendix \ref{sec: Deriving-Equation for QS}, yields the following
expression:
\begin{equation}
\frac{Q_{ji,k}(\varphi_{H})/Q_{\ell i,k}(\varphi_{H})}{Q_{ji,k}(\varphi_{L})/Q_{\ell i,k}(\varphi_{L})}\approx\begin{cases}
1 & \textrm{if }k\in\mathbb{K}_{\mathcal{C}}\\
\frac{N_{j,k}\left(\varphi_{H}\right)/N_{\ell,k}\left(\varphi_{H}\right)}{N_{j,k}\left(\varphi_{L}\right)/N_{\ell,k}\left(\varphi_{L}\right)}\ e^{\left(\tau_{\ell i}w_{\ell}a_{\ell,k}-\tau_{ji}w_{j}a_{j,k}\right)\frac{\vartheta_{i,k}}{w_{i}}} & \textrm{if }k\in\mathbb{K}_{\mathcal{D}}
\end{cases}.\label{eq: IQS}
\end{equation}
The above expression clearly indicates that there is no scope for
IQS in continuous industries. That is because IQS in these industries
is either cost-driven or taste-driven, both of which have been ruled
out by assumption. By contrast, in the case of discrete industries,
the above expression identifies two alternative drivers of IQS that
are neither cost-driven nor taste-driven. 

The first of these are affordability constraints, which operate (primarily)
through the exponential term in Equation \ref{eq: IQS}. The exponential
terms account for the fact that, due to affordability constraints,
an across-the-board increase in the cost of all varieties supplied
by country $j$ has a disproportionally higher effect on demand for
high-quality (high-cost) varieties. To demonstrate this, we can take
a basic derivative from Equation \ref{eq: IQS} with respect to the
national level cost-shifters, $a_{j,k}$, which implies that:
\[
\frac{\partial}{\partial a_{\ell,k}/a_{j,k}}\left\{ \frac{Q_{ji,k}(\varphi_{H})/Q_{\ell i,k}(\varphi_{H})}{Q_{ji,k}(\varphi_{L})/Q_{\ell i,k}(\varphi_{L})}\right\} >0\ \ \ \textrm{if }k\in\mathbb{K}_{\mathcal{D}}.
\]
Based on the above expression, countries tend to specialize in higher-quality
varieties in industries where they are more cost-efficient. The intuition
is that affordability constraints are more binding for high-quality
varieties. So, an improvement in cost-efficiency (keeping all else
the same) increases market access disproportionally more for high-quality
varieties. Such a link between within-industry IQS and industry-level
cost efficiency is quite unique and has little precedent in the literature.\footnote{The above result offers a possible micro-foundation for the assumption
underlying \citet{alcala2016specialization} that efficiency or productivity
is quality-biased.}

The second driver of IQS in discrete industries is the \emph{home-market
effect}, which operates through the number of firms in Equation \ref{eq: IQS}.
Unlike standard theories, though, the home-market effect here is not
driven by non-homothetic taste for quality. Instead, it is driven
(primarily) by cross-market differences in $\tilde{p}_{i}$. That
is, a higher price of services increases the home-market demand for
quality, leading to high-quality specialization due to scale economies
in production.\footnote{A higher $\tilde{p}_{i}$ lowers the demand elasticity for discrete
goods, which in turn creates relatively more demand for high-quality
(high-cost) product varieties---see \citet{Fajgelbaum2011,matsuyama2015home,dingel2016determinants}
for a similar argument, but in the case where the home-market effect
is driven by non-homothetic preferences over quality. } To demonstrate this point formally, we can take a simple derivate
from the free-entry condition with respect to $\tilde{p}_{i}$, noting
that $\partial\varepsilon_{ji,k}(\varphi)/\partial\tilde{p}_{i}<0$
for all $j$ and $i$. Doing so, implies the following, provided that
international trade costs are sufficiently high:
\begin{equation}
\frac{\partial}{\partial\tilde{p}_{j}/\tilde{p}_{\ell}}\left\{ \frac{N_{j,k}\left(\varphi_{H}\right)/N_{\ell,k}\left(\varphi_{H}\right)}{N_{j,k}\left(\varphi_{L}\right)/N_{\ell,k}\left(\varphi_{L}\right)}\right\} >0\ \Longrightarrow\ \frac{\partial}{\partial\tilde{p}_{j}/\tilde{p}_{\ell}}\left\{ \frac{Q_{ji,k}\left(\varphi_{H}\right)/Q_{\ell i,k}\left(\varphi_{H}\right)}{Q_{ji,k}\left(\varphi_{L}\right)/Q_{\ell i,k}\left(\varphi_{H}\right)}\right\} >0.\label{eq: prop_3 eq3}
\end{equation}
Let me reiterate that the \emph{home-market effect} in itself is not
specific to this model. In fact, \citet{costinot2016more} have already
characterized this result under very general conditions.\footnote{In the present model, the scale elasticity (defined as $\partial\ln Q_{ji,k}\left(\varphi\right)/\partial\ln N_{j,k}\left(\varphi\right)$
) is equal to ``\emph{one}.'' So, scale effects are strong enough
to trigger a \emph{strong} home-market effect (see \citet{costinot2016more}).
That is, all else equal, high-wage countries are not only \emph{gross}
exporters of high-quality but also \emph{net} exporters of high-quality
varieties. } What is instead novel, is that the \emph{home-market effect} is not
driven by non-homothetic taste but by cross-country differences in
the price of non-traded services. Along similar lines, affordability
constraints can also trigger IQS through the home-market effect, given
that high-quality varieties are more affordable and, therefore, more-demanded
in high-wage markets. The following proposition summarizes the above
results.
\begin{prop}
\label{prop: prop 1}Absent Ricardian comparative cost advantage and
non-nomothetic preferences: (a) There is no scope for either quality
specialization nor pricing-to-market in continuous industries; but
(b) quality specialization and pricing-to-market can still occur in
discrete industries due to cross-national differences in the price
of non-traded services and affordability constraints.
\end{prop}
\begin{figure}

\caption{\label{fig:Numerical-Example}Numerical Example: \emph{North-South
Quality Specialization}}

\begin{adjustwidth}{-0in}{0in}

\includegraphics[scale=0.45]{MPEC/Service_driven_IQS}\includegraphics[scale=0.45]{MPEC/Affordability_driven_IQS}

\end{adjustwidth}

\noindent\begin{minipage}[t]{1\columnwidth}%
\begin{spacing}{0.8}
\emph{\footnotesize{}Note:}{\footnotesize{} This graph compares demand
and output across different quality levels ($\varphi=1,..,50$) for
two economies (North and South) that differ only in their }\emph{\footnotesize{}Hicks-neutral}{\footnotesize{}
productivity, i.e., $a_{N}/a_{S}=2$. In the left panel, service prices
vary between countries but affordability constraints are non-binding.
In the right panel, affordability constraints are bindings but there
are no scale economies in production. More details about the simulated
economy are provided in Footnote \ref{fn: numerical example}.}
\end{spacing}
%
\end{minipage}
\end{figure}
Figure \ref{fig:Numerical-Example} elucidates the above arguments
using a simulated economy that features multiple countries and quality
tiers.\footnote{\label{fn: numerical example}The example involves 5 economies producing/consuming
only discrete varieties in 50 different quality tiers: $\varphi=1,...,50$.
Population size in all countries is normalized to one: $L_{j}=1$,
$\forall j$; but countries differ in their productivity: $a_{j}=1+0.25(j-1)$.
The region with the lowest productivity is labeled the \emph{South}
($s$) and the highest productivity region is labeled the \emph{North}
($n$): $a_{s}/a_{n}=2$. I also assume $c(\varphi)=0.08\varphi$,
$f(\varphi)=\varphi^{0.25}$, $\theta=4$, and $\tau_{ji}=2$ for
all $j\neq i$.} It compares two countries in this simulated economy (namely, North,
$n$, and South, $s$) that differ only in a Hicks-neutral cost shifter,
$a_{s}/a_{n}=2$. The left panel assumes all varieties are universally
affordable and IQS is driven by the combination of scale economies
and cross-national differences in $\tilde{p}_{i}$. The right panel
shuts down scale economies, so that IQS is driven solely by affordability
constraints.

A natural follow-up question to Proposition \ref{prop: prop 1} is
whether the alternative forces that drive PTM and IQS in discrete
industries are empirically distinguishable from cost-driven or taste-driven
IQS and PTM? To answer this question, it is useful to define two \emph{aggregate}
trade statistics that are often used to describe standard theories
of IQS and PTM. The first statistic is the industry-level export \emph{unit
price}, which can be defined as
\[
\bar{p}_{ji,k}\equiv\frac{\sum_{\varphi\in\varPhi_{k}}p_{ji,k}(\varphi)Q_{ji,k}(\varphi)}{\sum_{\varphi\in\varPhi_{k}}Q_{ji,k}(\varphi)}.
\]
The second statistic is the industry-level export quality, which can
be analogously defined as:s
\[
\bar{\varphi}_{ji,k}\equiv\frac{\sum_{\varphi\in\varPhi_{k}}\varphi\ Q_{ji,k}(\varphi)}{\sum_{\varphi\in\varPhi_{k}}Q_{ji,k}(\varphi)}.
\]
As proven in Appendix \ref{subsec: Proof Proposition 2}, combining
the above definition with our prior results on PTM and IQS, yields
the following proposition.\vspace{0.1in}
\begin{prop}
\label{prop: 2}$(i)$ Export unit prices are increasing in the price
of non-traded services in the importing market, but only so in discrete
industries:
\[
\begin{cases}
\frac{\partial\bar{p}_{ji,k}(.)}{\partial\tilde{p}_{i}}=0 & \textrm{\textrm{if }k\ensuremath{\in\mathbb{K}_{\mathcal{C}}}}\\
\frac{\partial\bar{p}_{ji,k}(.)}{\partial\tilde{p}_{i}}>0 & \textrm{\textrm{if }k\ensuremath{\in\mathbb{K}_{\mathcal{D}}}}
\end{cases};
\]
$(ii)$ Export unit prices are increasing in the price of non-traded
services in the exporting market, but only so in discrete industries.
\[
\begin{cases}
\frac{\partial\bar{p}_{ji,k}(.)}{\partial\ln\tilde{p}_{j}}=0 & \textrm{\textrm{if }k\ensuremath{\in\mathbb{K}_{\mathcal{C}}}}\\
\frac{\partial\bar{p}_{ji,k}(.)}{\partial\ln\tilde{p}_{j}}>0 & \textrm{\textrm{if }k\ensuremath{\in\mathbb{K}_{\mathcal{D}}}}
\end{cases};
\]
(iii) Export quality is increasing in a country's industry-level degree
of comparative cost advantage, but only so in discrete industries:
\[
\begin{cases}
\frac{\partial\bar{\varphi}_{ji,k}(.)}{\partial a_{j,k}}=0 & \textrm{\textrm{if }k\ensuremath{\in\mathbb{K}_{\mathcal{C}}}}\\
\frac{\partial\bar{\varphi}_{ji,k}(.)}{\partial a_{j,k}}<0 & \textrm{\textrm{if }k\ensuremath{\in\mathbb{K}_{\mathcal{D}}}}
\end{cases}.
\]
\end{prop}
The above proposition presents a set of comparative static results
indicating that \emph{the price of non-traded services }and \emph{cost-efficiency}
have differential effects on aggregate export price/quality levels
in discrete versus continuous industries. In the following section,
I will use the above proposition to empirically test my theory. Before
doing so, let me note that $\tilde{p_{i}}$, $\tilde{p}_{j}$, $a_{j,k}$
affect the price of exports/imports through both the quality and markup
channels. These two channels are independent, but (in the present
setup) they always reinforce one another's effect on prices. That
being the case, I can just rely on readily observable price data to
test the model's predictions. Encouragingly, as I will show next,
the above predictions are all borne out.

\section{\label{sec:Empirical-Evidence}Testing the Predictions of the Model}

Proposition \ref{prop: 2} will serve as the basis for testing my
theory. In particular, I test $(a)$ if the export unit price, $\bar{p}_{ji,k}$,
is positively related to $\tilde{p}_{i}$, $\tilde{p}_{j}$, and $a_{j,k}$;
and also $(b)$ if these relationships are amplified by \emph{discreteness}.
To conduct my analysis, I merge two different datasets. The first
of which is from the 2011 International Comparison Program (ICP) compiled
by the World Bank Development Data Group.\footnote{The ICP data has been released for years 1970, 1973, 1975, 1980, 1985,
1993, 2005, and 2011. Here, I choose 2011 as the benchmark year for
my analysis.} The ICP data reports national expenditures and PPPs for more than
150 \textquotedblleft Basic Headings\textquotedblright{} categories.
This information allows me to construct the aggregate price of non-traded
services, $\tilde{p}_{i}$ for 182 countries.\footnote{I treat the following basic headings in the ICP as non-tradable: 'Housing',
'Education', 'Health', and 'Restaurants and Hotels'. In the main analysis
'Communication' is not treated as a non-tradable heading, but treating
it as one does not alter any of the results qualitatively. Moreover,
similar \citet{sposi2015trade}, I treat \textquotedblleft Change
in inventories and valuables,\textquotedblright{} as unclassified.} 

I merge the \noun{ICP }data with the\noun{ Comtrade-baci} database,
which is the most comprehensive database on export unit prices. The
\noun{baci} data is compiled by the CEPII based on the United Nations
\noun{Comtrade} database but uses a harmonization methodology to reconcile
mirror flows. This harmonization approach provides a more complete
geographical coverage than if only a single direction of \noun{Comtrade}
statistics were to be used---\citet{gaulier2010baci} describe the
\noun{Baci} approach to harmonization, which involves computing a
weighted average of mirror flows. For the year 2011, the \noun{Comtrade-baci}
database reports physical quantities and values of bilateral trade
flows by 6-digit HS product category, allowing me to construct the
corresponding unit price levels. 

Combining the ICP and the BACI datasets, leaves me with a sample of
bilateral trade statistics for 147 importers and 147 exporters across
5,231 HS6 product categories in the year 2011. For each country in
the sample, \citet{head2010erosion} report matching data for GDP,
distance, and population size. So, altogether, for every \emph{exporting
country $j$$\times$importing country $i$$\times$HS6 industry $k$}
(observation $ji,k$), I can construct a quantity-weighted unit price,
$\bar{p}_{ji,k}$, as well as exporter/importer characteristics such
as total GDP, GDP per capita, geo-distance, $\tilde{p}_{j}$, and
$\tilde{p}_{i}$. To test the prediction $(iii)$ from Proposition
\ref{prop: 2}, I also need data on the degree of industry-level comparative
cost advantage. Unfortunately, the technology parameters, $a_{j,k}$,
are not directly observable. So, as is common in the literature, I
infer them from sales data. To this end, I construct the Balassa index
of revealed comparative advantage for each country $j$ in industry
$k$, as follows:\footnote{The results of my analysis are qualitatively robust to adopting alternative
measures of revealed comparative advantage, such as that proposed
by \citet{proudman2000evolving}.}
\[
\textrm{RCA}_{j,k}=\frac{\sum_{i}X_{ji,k}/\sum_{i}\sum_{g}X_{ji,g}}{\sum_{\ell}\sum_{i}X_{\ell i,k}/\sum_{\ell}\sum_{i}\sum_{g}X_{\ell i,g}}
\]
The above measure provides a theory-consistent proxy for $a_{j,k}$;
since from the lens of the model, a lower $a_{j,k}$ (ceteris paribus)
is associated with a higher $RCA_{j,k}$.

The \noun{Baci }data does not explicitly report the units in which
quantity is measured.\footnote{Unlike the U.S. data which identify the unit in which physical quantity
is measured, the \noun{Baci }data only reports the weight-equivalent
of physical quantity for all observations. } So, I use the U.S. import and export data to identify the discrete
product categories, with $\mathbbm{1}_{k\in\mathbb{K_{\mathcal{D}}}}$
denoting the corresponding discrete product dummy.\footnote{By this classification, discrete goods account for 36\% of global
trade in the BACI dataset.} This approach can classify 90\% of observations in the \noun{Baci}.
Alternatively, to preserve all observations, I can classify a HS6
product category as discrete if it belongs to either the ``Machinery,''
``Electrical \& Optical Equipment,'' or ``Transport Equipment''
sectors. Both approaches deliver qualitatively similar results. 

To test the model, I run the following regression on the final sample
that includes 5,249,424 observations:\footnote{Using a linear regression to test the comparative static results outlined
by Proposition \ref{prop: 2}, requires some clarification. In the
words of \citet{reiss2007structural}, the standard reasoning behind
this choice is that \emph{``When the conditional expectation of y
is nonlinear in x,} s\emph{tatistical theory tells us (under certain
sampling assumptions) that a regression provides a best (minimum expected
squared prediction error) linear approximation to the nonlinear conditional
expectation function.'' }However, as \citet{reiss2007structural}
point out, there are basic caveats associated with using a linear
reduced-form specification to test comparative static results.}
\begin{align}
\hspace{-0.2in}\ln\overline{p}_{ji,k}= & \left(\beta_{1}+\beta_{1}^{\mathcal{D}}\cdot\mathbbm{1}_{k\in\mathbb{K_{\mathcal{D}}}}\right)\ln\tilde{p}_{i}\ +\ \left(\beta_{2}+\beta_{2}^{\mathcal{D}}\cdot\mathbbm{1}_{k\in\mathbb{K_{\mathcal{D}}}}\right)\ln\tilde{p}_{j}\nonumber \\
 & \ +\left(\beta_{3}+\beta_{3}^{\mathcal{D}}\cdot\mathbbm{1}_{k\in\mathbb{K_{\mathcal{D}}}}\right)\ln\textrm{RCA}_{j,k}\ +\ Controls_{ji}+\delta_{k}+\epsilon_{ji,k},\label{eq: Baci Regression}
\end{align}
where $\delta_{k}$ controls for HS6 fixed effects; and $Controls_{ji}$
is composed of additional controls including total GDP, GDP per capita,
and geo-distance. %
\begin{comment}
\begin{align*}
\ln\overline{p}_{ji,k}= & \left(\beta_{1}+\beta_{2}\cdot disc\right)RCA_{j,k}+\left(\beta_{3}+\beta_{4}\cdot disc\right)\ln\tilde{p}_{j}+\left(\beta_{5}+\beta_{6}\cdot disc\right)\ln\tilde{p}_{i}\\
+ & \beta_{7}\ln w_{j}+\beta_{8}\ln w_{i}+\beta_{9}\ln Y_{j}+\beta_{10}\ln Y_{i}+\beta_{11}\ln\textrm{Dist}_{ji}+\delta_{k}+\epsilon_{ji,k},
\end{align*}
\end{comment}
{} In addition to the baseline specification, I run the above regression
under alternative specifications, controlling for \emph{exporter,
importer,} and \emph{exporter$\times$importer} fixed effects.\footnote{As in Section \ref{sec:Suggestive-Evidence}, since I am dealing with
grouped data each observation is weighted by trade quantity, which
is a theory-consistent measure of the number of firms grouped into
each observation---refer to Section \ref{sec:Suggestive-Evidence}
for a more elaborate discussion.} 

Based on Proposition \ref{prop: 2}, the estimated coefficients $\beta_{1}^{\mathcal{D}}$,
$\beta_{2}^{\mathcal{D}}$, and $\beta_{3}^{\mathcal{D}}$ should
be positive and statistically significant.\footnote{The estimated coefficients $\beta_{1}$, $\beta_{2}$, and $\beta_{3}$
may be positive given that $\mathbbm{1}_{k\in\mathbb{K_{\mathcal{D}}}}$
does not perfectly identify discrete product categories. More specifically,
an industry where $\mathbbm{1}_{k\in\mathbb{K_{\mathcal{D}}}}=0$
may still contain a small fraction of observations that correspond
to discrete goods. } Encouragingly, the estimation results presented in Table \ref{tab:price-regression}
support theses predictions.

In line with prediction $(i)$ from Proposition \ref{prop: 2}, the
elasticity of export prices with respect to the price of non-traded
services in the destination market is significantly higher in discrete
industries. To provide numbers, consider Column 3 that controls for
exporter$\times$HS6 fixed effects. A 10\% increase in $\tilde{p}_{i}$,
is associated with a 1.53\% increase in the price of exports to market
$i$ in discrete industries. But in continuous industries, the effect
is only 0.06\% and statistically insignificant. To put this number
in perspective, the same estimation indicates that a 10\% increase
in the GDP per capita of country $i$ is associated with an only 1.6\%
increase in the price of exports to that market.\footnote{Adding $\ln\textrm{GDP per capita}_{i}\times\mathbbm{1}_{k\in\mathbb{K_{\mathcal{D}}}}$
as an additional control does not change the sign of the coefficients
reported in Table \ref{tab:price-regression}. However, it increases
the standard errors, which can be indicative of multicollinearity
between $\ln\textrm{GDP per capita}_{i}\times\mathbbm{1}_{k\in\mathbb{K_{\mathcal{D}}}}$
and $\mathbbm{1}_{k\in\mathbb{K_{\mathcal{D}}}}\times\ln\tilde{p}_{i}$.}

\begin{table}[!htb] \def\sym#1{\ifmmode^{#1}\else\(^{#1}\)\fi}
\begin{centering}
\begin{adjustwidth}{-0in}{-0in}
\caption{The determinants of IQS and PTM (\emph{Dependent Variable}: $\ln \bar{p}_{ji,k}$)}\label{tab:price-regression} 
\vspace{-0.1in}
\resizebox{1\textwidth}{!}{%  %
\begin{tabular}{l*{5}{D{.}{.}{-1}}}

\toprule                  
Regressor (log)  &\multicolumn{1}{c}{(1)} &\multicolumn{1}{c}{(2)} &\multicolumn{1}{c}{(3)}  &\multicolumn{1}{c}{(4)} \\ 
\midrule

{\large $\tilde{p}_i$}        & 0.074\sym{*} & ... & 0.006 & ... \\ & (0.039) & & (0.022) & \\ \addlinespace   
{\large $\tilde{p}_i$}$\times$ Discrete & 0.015 & ... & 0.153\sym{***}& 0.161\sym{***}\\ & (0.108) & & (0.021) & (0.021) \\ \addlinespace

{\large $\tilde{p}_j$}            & 0.070 & 0.073\sym{**} & ... & ... \\ & (0.051) & (0.033) & & \\ \addlinespace  
{\large $\tilde{p}_j$} $\times$ Discrete  &       0.061\sym{*}  & 0.067\sym{**} & ... & 0.104\sym{***}\\ & (0.033) & (0.029) & & (0.029) \\ \addlinespace   

{\large $\textrm{RCA}_{j,k}$}        & 0.032\sym{***}& 0.021\sym{***}& ... & 0.029\sym{***}\\ & (0.005) & (0.004) & & (0.005) \\ \addlinespace 
{\large $\textrm{RCA}_{j,k}$} $\times$ Discrete        & 0.091\sym{***}& 0.055\sym{***}& ... & 0.071\sym{***}\\ & (0.015) & (0.011) & & (0.012) \\ \addlinespace 

Controls for GDP \& GDP p/c         & \multicolumn{1}{c}{Y} & \multicolumn{1}{c}{Y} & \multicolumn{1}{c}{Y} &  \multicolumn{1}{c}{...} \\
Controls for Distance        & \multicolumn{1}{c}{Y} & \multicolumn{1}{c}{Y} & \multicolumn{1}{c}{Y} & \multicolumn{1}{c}{...} \\
HS6 FE       & \multicolumn{1}{c}{Y} & \multicolumn{1}{c}{...} & \multicolumn{1}{c}{...} & \multicolumn{1}{c}{Y} \\
Exporter$\times$HS6 FE       & \multicolumn{1}{c}{N} & \multicolumn{1}{c}{N} & \multicolumn{1}{c}{Y} & \multicolumn{1}{c}{N} \\
Importer$\times$HS6 FE       & \multicolumn{1}{c}{N} & \multicolumn{1}{c}{Y} & \multicolumn{1}{c}{N} & \multicolumn{1}{c}{N} \\
Exporter$\times$Importer FE & \multicolumn{1}{c}{N} & \multicolumn{1}{c}{...} & \multicolumn{1}{c}{...} & \multicolumn{1}{c}{Y} \\
Observations        &     \multicolumn{1}{c}{4,734,290}         &    \multicolumn{1}{c}{4,734,290}        &     \multicolumn{1}{c}{4,734,290}  &  \multicolumn{1}{c}{4,734,290}    \\
\bottomrule \addlinespace 
\end{tabular} }
\noindent\begin{minipage}[t]{1\columnwidth} \begin{spacing}{0.8} {\footnotesize \emph{Note:} this table corresponds to estimating Equation \ref{eq: Baci Regression}. Standard errors are reported in parentheses and clustered by destination and origin country. Unit prices are constructed using the BACI database in 2011; the price of non-traded services are constructed using the ICP database in the same year; and each observation in the sample is weighted by trade quantity. \sym{*} \(p<0.10\), \sym{**} \(p<0.05\), \sym{***} \(p<0.01\)}
\end{spacing} \end{minipage}
\end{adjustwidth}
\end{centering}
\end{table}

Prediction $(ii)$ is also borne out, albeit less decisively. Discreteness
has a positive and statistically significant effect on how export
prices vary with the price of non-traded services in the origin country.
In general, this outcome may be driven by services being used as an
input in production. This channel perhaps explains why $\bar{p}_{ji,k}$
and $\tilde{p}_{j}$ are positively correlated, even in continuous
industries. However, in line with the discrete model's prediction,
this positive correlation is significantly more pronounced in discrete
industries; standing around 0.67 percentage points higher (Table \ref{tab:price-regression},
Column 2). 

Finally, prediction $(iii)$ is also borne out: In discrete industries,
a 10\% increase in the origin country's degree of comparative advantage,$\textrm{RCA}_{j,k}$,
is associated with a 0.77\% increase in export prices. In continuous
industries, by comparison, $\textrm{RCA}_{j,k}$ is associated with
a mere 0.21\% increase in export prices (Table \ref{tab:price-regression},
Column 2).\footnote{That $\textrm{RCA}_{j,k}$ matters in continuous industries may be
driven by two factors. First, industries coded as continuous in my
analysis may feature a small number of discrete categories. Second,
that $\textrm{RCA}_{j,k}$ is inherently related to export prices
as it reflects either a comparative \emph{quality} advantage or a
comparative \emph{cost} advantage. In the former case, $\textrm{RCA}_{j,k}$
is positively related to $\overline{p}_{ji,k}$, while in the latter
case it is negatively related to $\overline{p}_{ji,k}$.}%
\begin{comment}
\begin{table}[htbp]\centering \def\sym#1{\ifmmode^{#1}\else\(^{#1}\)\fi} 
\begin{adjustwidth}{-0.35in}{-0.35in}
\hspace{-0.1in}
\caption{The determinants of quality specialization and pricing-to-market}\label{tab:price-regression}  \begin{tabular}{l*{8}{D{.}{.}{-1}}} 
           \multicolumn{5}{c}{\emph{dependent variable}: $\ln \bar{p}_{ji,k}$} \\
\toprule                  
regressor  &\multicolumn{1}{c}{(1)} &\multicolumn{1}{c}{(2)} &\multicolumn{1}{c}{(3)} &\multicolumn{1}{c}{(4)}\\  \midrule
$\ln RCA_{j,k}$   &                     &      0.039\sym{***}&                     &      0.039\sym{***}\\                    		 &                  &  (0.0003)         &                     &  (0.0003)         \\ \addlinespace 
$\ln \tilde{p}_i$ &                     &                     &      0.0585\sym{***}&      0.0597\sym{***}\\                     		&                     &                   &   (0.0015)         &   (0.0015)         \\ \addlinespace 
$\ln w_{j}$             &       0.152\sym{***}&       0.155\sym{***}&       0.151\sym{***}&       0.154\sym{***}\\                     &   (0.0082)         &  (0.0003)         &  (0.0003)         &  (0.0003)         \\ \addlinespace 
$\ln Y_{j}$         &     -0.068\sym{***}&     -0.057\sym{***}&     -0.068\sym{***}&     -0.057\sym{***}\\                     &   (0.0043)         &  (0.00025)         &  (0.0002)         &  (0.0002)         \\ \addlinespace 

$\ln w_{i}$                &      0.062\sym{***}&      0.061\sym{***}&      0.035\sym{***}&      0.033\sym{***}\\                     &   (0.0064)         &  (0.0004)         &  (0.0008)         &  (0.0008)         \\ \addlinespace 
$\ln Y_{i} $                &     -0.029\sym{***}&     -0.031\sym{***}&     -0.030\sym{***}&     -0.031\sym{***}\\                     &   (0.0074)         &  (0.0002)         &  (0.0002)         &  (0.0002)         \\ \addlinespace 
$\ln \textrm{Dist}_{ji}$                  &      0.075\sym{***}&      0.067\sym{***}&      0.076\sym{***}&      0.068\sym{***}\\                     &   (0.0061)         &  (0.0003)         &  (0.0003)         &  (0.0003)         \\  \\
Observations        &     \multicolumn{1}{c}{5,925,922}         &     \multicolumn{1}{c}{5,925,922}        &     \multicolumn{1}{c}{5,925,922}         &     \multicolumn{1}{c}{5,925,922}         \\ 
\bottomrule \multicolumn{5}{l}{\footnotesize Robust standard errors in parentheses}\\ \multicolumn{5}{l}{\footnotesize \sym{***} \(p<0.001\)}\\ \end{tabular} 
\end{adjustwidth}
\end{table}
\end{comment}
{}  One caveat here, though, is that I am not testing the model's
predictions \emph{directly} with information on quality and markups.
The above analysis, therefore, cannot rule out the case where $\tilde{p}_{j}$,
$\tilde{p}_{i}$, and $\textrm{RCA}_{j,k}$ have countervailing effects
on export markups and qualities. Instead, loosely speaking, the analysis
can only identify the joint effect of these variables on \emph{quality}$\times$\emph{markup}.
This limitation is, however, hard to overcome as estimates for quality
and markup are difficult to attain at the scale in which my analysis
is conducted. This limitation notwithstanding, the above results plus
those presented in Section \ref{sec:Suggestive-Evidence} corroborate
the assertion that IQS and PTM are magnified by additional forces
in discrete industries. 

\paragraph*{Restricting Sample to Differentiated or Final Goods.}

In the above analysis, discreteness may be accounting for other omitted
product characteristics such as the degree of differentiation or intent
of final use. To address this issue, I reestimate Equation \ref{eq: Baci Regression}
on two restricted samples: First, a sample that only includes differentiated
HS6 products as classified by \citet{Rauch1999}. Second, a sample
that includes only HS6 products that are intended for final consumption,
based on the Broad Economic Categories (BEC) classification.

The results derived from these restricted samples are reported in
Table \ref{tab:price-regression-robust}. They indicate that, even
within differentiated goods, discreteness has a significant effect
on how export prices, $\bar{p}_{ji,k},$ vary with the price of services
in the origin and destination economies, $\tilde{p}_{j}$, and $\tilde{p}_{i}$.
Likewise, discreteness has a significant effect on how export prices
vary with the origin country's degree of comparative advantage as
measured by $\textrm{RCA}_{j,k}$. The same applies when we restrict
attention to only final consumption goods. Appendix \ref{sec:Additional-Test-the}
presents an additional test in which HS6 products are grouped into
\emph{more-} or \emph{less-differentiated} bins according to the demand
elasticity estimates in \citet{Broda2006}. Estimating Equation \ref{eq: Baci Regression}
separately for these product groups produces similar results.

\begin{table}[!htb] \def\sym#1{\ifmmode^{#1}\else\(^{#1}\)\fi}
\begin{centering}
\begin{adjustwidth}{-0.5in}{-0in}
\caption{Testing the model on only \emph{differentiated} and \emph{final} goods (\emph{Dependent}: $\ln \bar{p}_{ji,k}$)}\label{tab:price-regression-robust} 
\vspace{-0.1in}
\resizebox{1.2\textwidth}{!}{%  %
\begin{tabular}{l*{8}{D{.}{.}{-1}}}

\toprule                    
& \multicolumn{3}{c}{\textbf{Final Goods}} & \phantom{abc}  & \multicolumn{3}{c}{\textbf{Differentiated Goods}} \\
\cmidrule{2-4} \cmidrule{6-8}

Regressor (log)  &\multicolumn{1}{c}{(1)} &\multicolumn{1}{c}{(2)} &\multicolumn{1}{c}{(3)} && &\multicolumn{1}{c}{(1)} &\multicolumn{1}{c}{(2)} &\multicolumn{1}{c}{(3)} \\
\toprule   

{\large $\tilde{p}_i$}  & ... & 0.074\sym{**} & ... && ... & 0.056\sym{*} & ... \\ & & (0.034) & && & (0.034) & \\ \addlinespace 
{\large $\tilde{p}_i$} $\times$ Discrete & ... & 0.179\sym{***}& 0.183\sym{***}&& ... & 0.148\sym{***}& 0.156\sym{***}\\ & & (0.019) & (0.019) & && (0.021) & (0.016) \\ \addlinespace

{\large $\tilde{p}_j$} & -0.087 & ... & ... && 0.079\sym{***}& ... & ... \\ 
& (0.086) & & && (0.039) & & \\ \addlinespace  
{\large $\tilde{p}_j$} $\times$ Discrete & 0.186\sym{***}& ... & 0.155\sym{***}&& 0.052\sym{***}& ... & 0.097\sym{***}\\ & (0.025) & & (0.021) && (0.005) & & (0.019) \\ \addlinespace

{\large $\textrm{RCA}_{j,k}$}& 0.036\sym{***}& ... & 0.050\sym{***} && 0.017\sym{***}& ... & 0.029\sym{***}\\
& (0.006) & & (0.005) && (0.007) & & (0.006) \\ \addlinespace 
{\large $\textrm{RCA}_{j,k}$} $\times$ Discrete & 0.004\sym{***}& ... & 0.012 && 0.057\sym{***}& ... & 0.070\sym{***}\\ & (0.001) & & (0.008) && (0.011) & & (0.010) \\ \addlinespace 
 

Controls for GDP \& GDP p/c         & \multicolumn{1}{c}{Y} & \multicolumn{1}{c}{Y} & \multicolumn{1}{c}{...} &&   \multicolumn{1}{c}{Y} & \multicolumn{1}{c}{Y} & \multicolumn{1}{c}{...} \\
Controls for Distance        & \multicolumn{1}{c}{Y} & \multicolumn{1}{c}{Y} & \multicolumn{1}{c}{...} &&   \multicolumn{1}{c}{Y} & \multicolumn{1}{c}{Y} & \multicolumn{1}{c}{...} \\
HS6 FE       & \multicolumn{1}{c}{...} & \multicolumn{1}{c}{...} & \multicolumn{1}{c}{Y} && \multicolumn{1}{c}{...} & \multicolumn{1}{c}{...} & \multicolumn{1}{c}{Y} \\
Exporter$\times$HS6 FE       & \multicolumn{1}{c}{N} & \multicolumn{1}{c}{Y} & \multicolumn{1}{c}{N} && \multicolumn{1}{c}{N} & \multicolumn{1}{c}{Y} & \multicolumn{1}{c}{N} \\
Importer$\times$HS6 FE       & \multicolumn{1}{c}{Y} & \multicolumn{1}{c}{N} & \multicolumn{1}{c}{N} && \multicolumn{1}{c}{Y} & \multicolumn{1}{c}{N} & \multicolumn{1}{c}{N} \\
Exporter$\times$Importer FE & \multicolumn{1}{c}{...} & \multicolumn{1}{c}{...} & \multicolumn{1}{c}{Y} && \multicolumn{1}{c}{...} & \multicolumn{1}{c}{...} & \multicolumn{1}{c}{Y} \\
Observations                &    \multicolumn{3}{c}{1,413,000}        &&     \multicolumn{3}{c}{3,589,000}       \\
\bottomrule \addlinespace 
\end{tabular} }
\noindent\begin{minipage}[t]{1.2\columnwidth} \begin{spacing}{0.8} {\footnotesize \emph{Note:} this table estimates Equation \ref{eq: Baci Regression} separately for differentiated and final good categories. Differentiated goods are identified based on  \citet{Rauch1999}. Final goods are identified using the BEC classification. Standard errors are reported in parentheses and clustered by origin and destination country.  Unit prices are constructed using the BACI database in 2011; the price of non-traded services are constructed using the ICP database in the same year; and each observation in the sample is weighted by trade quantity. \sym{*} \(p<0.10\), \sym{**} \(p<0.05\), \sym{***} \(p<0.01\)}
\end{spacing} \end{minipage}
\end{adjustwidth}
\end{centering}
\end{table}

\section{\label{sec: GT}The Big Picture: \emph{The Gains from Discrete Trade}}

Does accounting for discrete trade revise our understanding of the
gains from trade? To answer this question, we can appeal to the result
in \citet{McFadden1978} that the expected welfare of each individual
in country $i$ with respect to discrete industry $k\in\mathbb{K}_{\mathcal{D}}$
can be expressed as follows:
\[
W_{i,k}\left(y\right)=\exp(y)\left(\sum_{\varphi\in\varPhi_{k}}\left[\varphi\sum_{j\in\mathcal{J}_{i,k}\left(y,\varphi\right)}N_{j,k}\left(\varphi\right)\exp\left(p_{ji,k}\left(\varphi\right)\right)^{-\frac{\theta_{k}}{\tilde{p}_{i}}}\right]^{\frac{\gamma_{k}}{\theta_{k}}}\right)^{\frac{1}{\gamma_{k}}},
\]
where $\mathcal{J}_{i,k}\left(y,\varphi\right)=\left\{ j\in\mathbb{C}\mid p_{ji,k}\left(\varphi\right)\leq e_{k}w_{i}y\right\} $
denotes the set of countries from which an individual with endowment
$y$ can affordably buy quality $\varphi$ varieties. Correspondingly,
the expected welfare of the same individual under \emph{autarky }(denoted
by $\textrm{A}$) is given by
\[
W_{i,k}^{\textrm{A}}\left(y\right)=\exp(y)\left(\sum_{\varphi\in\varPhi_{k}\left(y\right)}\left[\varphi N_{i,k}\left(\varphi\right)\exp\left(p_{ii,k}\left(\varphi\right)\right)^{-\frac{\theta_{k}}{\tilde{p}_{i}}}\right]^{\frac{\gamma_{k}}{\theta_{k}}}\right)^{\gamma_{k}},\ \ \ k\in\mathbb{K}_{\mathcal{D}}
\]
where $\varPhi_{k}\left(y\right)=\left\{ \varphi\in\varPhi_{k}\mid p_{ii,k}\left(\varphi\right)\leq e_{k}w_{i}y\right\} $
denotes the set of quality levels that can be affordably purchased
from domestic suppliers. Combining the above two equations; using
the exact hat-algebra notation, $\hat{x}\equiv x/x^{\textrm{A}}$;
and noting that the gains from trade in \emph{continuous} industries
are governed by the \citet{Arkolakis} (ACR) formula, yields the following
proposition.%
\begin{comment}
\[
\frac{W_{i,k}^{A}\left(y\right)}{W_{i,k}\left(y\right)}=\sum_{\varphi\in\varPhi_{k}}\frac{\left[f_{k}\left(\varphi\right)N_{i,k}\left(\varphi\right)v_{k}\left(y,\tilde{p}_{i},p_{ii,k}\right)\right]^{\gamma_{k}}}{\sum_{\varphi\in\varPhi_{k}}f_{k}\left(\varphi\right)\left[\sum_{j}N_{j,k}\left(\varphi\right)v_{k}\left(y,\tilde{p}_{i},p_{ji,k}\right)^{\theta_{k}}\right]^{\frac{\gamma_{k}}{\theta_{k}}}}
\]

\[
\frac{W_{i}^{A}}{W_{i}}=\frac{\sum_{\varphi}\left(\frac{N_{i,\varphi}^{A}}{N_{i,\varphi}}\lambda_{ii,\varphi}\right)^{\frac{1}{\theta}}f(\varphi)\left(\sum_{k}N_{k,\varphi}h\left(1-\rho_{ki,\varphi}\right)^{\theta}\right)^{\frac{1}{\theta}}}{\sum_{\varphi}f(\varphi)\left(\sum_{k}N_{k,\varphi}h\left(1-\rho_{ki,\varphi}\right)^{\theta}\right)^{\frac{1}{\theta}}}=\sum_{\varphi=1}^{\bar{\varphi}}\left(\frac{\lambda_{ii,\varphi}}{\hat{N}_{i,\varphi}}\right)^{\frac{1}{\theta}}
\]
\[
\frac{\Delta W_{i,k}\left(y\right)}{W_{i,k}(y)}=1-\left(\sum_{\varphi\in\varPhi_{k}}\left[\lambda_{ii,k}\left(\varphi,y\right)^{\frac{1}{\theta_{k}}}/\hat{N}_{i,k}\left(\varphi\right)^{-\frac{1}{\theta_{k}}}\ \widehat{v_{k}\left(y,\tilde{p}_{i},p_{ii,k}\right)}\right]^{\gamma_{k}}\Lambda_{i,k}\left(\varphi,y\right)\right)^{\frac{1}{\gamma_{k}}}
\]
\end{comment}
\vspace{0.1in}
\begin{prop}
\label{prop: GT}{[}\textbf{The Gains from Trade}{]} The gains from
trade\textup{, $GT_{i}\left(y\right)\equiv\Delta W_{i}\left(y\right)/W_{i}(y)$},
for income group $y$ residing in country $i$ are given by
\[
GT_{i}\left(y\right)=1-\prod_{k\in\mathbb{K}}\left(\sum_{\varphi\in\varPhi_{k}}\lambda_{ii,k}\left(\varphi,y\right)^{\frac{\gamma_{k}}{\theta_{k}}}\widehat{N}_{i,k}\left(\varphi\right){}^{-\frac{\gamma_{k}}{\theta_{k}}}\ e_{i,k}\left(\varphi,y\right)\right)^{\frac{e_{k}}{\gamma_{k}}},
\]
where $e_{i,k}\left(\varphi,y\right)$ denotes the probability of
choosing quality level $\varphi\in\varPhi_{k}$, while $\lambda_{ii,k}\left(\varphi,y\right)$
denotes the probability of purchasing a domestic variety.\vspace{0.1in}
\end{prop}
Based on the above proposition, the gains from trade are materialized
through two distinct channels: $(i)$ the \emph{scale-driven }gains
that depend on $\widehat{N}_{i,k}\left(\varphi\right)$; and $(ii)$
the \emph{import variety gains} that depend on $\lambda_{ii,k}\left(\varphi,y\right)$.
Beholding this decomposition, let me first qualitatively discuss the
distributional consequences of trade. Then, I will formally quantify
the aggregate gains from discrete trade. 

\paragraph*{\emph{The Distributional Gains from Trade.}}

A straightforward corollary of Proposition \ref{prop: GT} is that
the gains from trade favor consumers that are more likely to purchase
foreign varieties, i.e., consumers who exhibit a low $\lambda_{ii,k}\left(\varphi,y\right)$.
In discrete industries, if trade costs are sufficiently large, low-income
individuals cannot afford many imported varieties. Hence, even if
individual-level preferences are homothetic, $\lambda_{ii,k}\left(\varphi,y\right)$
can be decreasing in $y$, and the gains from trade can favor high-income
(high-$y$) individuals.

So, how can we reconcile the above claim with the claim in \citet{fajgelbaum2014measuring}
that the gains from trade favor low-income consumers who spend a greater
fraction of their income on tradables? The answer lies in the different
interpretations of $\lambda_{ii,k}\left(\varphi,y\right)$ in discrete
versus continuous industries. In continuous industries, $\lambda_{ii,k}\left(\varphi,y\right)$
perfectly coincides with the share of expenditure on domestic varieties.
In discrete industries, however, $\lambda_{ii,k}\left(\varphi,y\right)$
can be strictly unrelated to the expenditure share. In particular,
high-income individuals may be more likely to purchase imported varieties;
but since they buy only one unit of each good, such a purchase may
exhaust a smaller fraction of their income.%
\begin{comment}
Affordability constraints also have a second-order \emph{pro-competitive}
effect on the gains from trade. Specifically, the additive makeup
charged by domestic firms in market $i$ is equal to $w_{i}\left[\frac{\theta_{k}}{\tilde{a}_{i}}+\frac{\zeta_{i}}{e_{k}}\rho_{ii,k}(\varphi)\right]^{-1}$,
where $\rho_{ii}\left(\varphi\right)$ denotes the importance of the
marginal consumer as measured by purchasing probability, with $\rho_{ii}\left(\varphi\right)\approx1$
to a first-order. The introduction of imported varieties will increase
the weight of the marginal consumers, as they cannot afford some of
the newly-introduced import varieties that come with a high transport
cost tags. As a result domestic firms are encouraged to reduce their
markup to a second-order to accommodate these marginal consumers. 
\end{comment}
{} Quantifying these distributional gains, though, requires detailed
international expenditure data across various income groups. That
type of data is largely inaccessible to researchers at this point.
Given this underlying challenge, I turn my attention instead to quantifying
the aggregate gains from trade.

\paragraph*{\emph{The Aggregate Gains from Trade.}}

My discussion above indicated that affordability constraints matter
for the distributional consequences of trade. Now, I purposely abstract
from these constraints to quantify the representative gains from trade
for the case where affordability constraints are non-binding and preferences
are \emph{Gorman representable}. In this special case, unobservable
quality differentiation is also irrelevant for the aggregate gains
from trade, and we can disregard it hereafter.\footnote{To elaborate, if preferences are \emph{Gorman representable}, then
$\lambda_{ii,k}(y,\varphi)=\lambda_{ii,k}(y,\varphi')$ for all $\varphi$
and $\varphi'\in\varPhi_{k}$. Consequently, the quality dimension
disappears from the gains from trade equation in this case. This outcome,
though, is an artifact of assuming that $\theta_{k}$ is uniform across
different quality tiers within the same industry. } The representative gains from trade can, therefore, be expressed
as follows:%
\begin{comment}
{[}Suppose $v_{k}=y/p_{ii,k}^{\alpha}\tilde{p}_{i}^{1-\alpha}$ ,
as $\alpha\rightarrow0$ then why does $N_{i}$ remain relevant? because
the taste matters! so we have to discount the taste parameter!{]} 
\end{comment}
\begin{equation}
GT_{i}=1-\prod_{k\in\mathbb{K}}\left(\lambda_{ii,k}/\hat{N}_{i,k}\right)^{\frac{e_{i,k}}{\theta_{k}}}.\label{eq: GT}
\end{equation}
From the lens of the above formula, \emph{discrete} and \emph{continuous}
industries are distinguished by $(i)$ the different interpretations
of $\lambda_{ii,k}$, and $(ii)$ the different characterizations
of $\hat{N}_{i,k}$ in discrete versus continuous industries. If we
treat all industries as \emph{continuous}, the above formula can be
simply interpreted as the celebrated ACR formula.

\subsection{Computing the Gains From Trade}

To transparently highlight the role of discreteness, I use Equation
\ref{eq: GT} to quantify the gains from trade in two extreme cases.
First, as is common in the literature, I treat the traded sector of
the economy as one integrated \emph{continuous} industry. Second,
I opt for the other extreme, treating the traded sector as one integrated
\emph{discrete} industry. Accordingly, I hereafter drop the industry
subscript, $k$, and use $\mathcal{C}$ and $\mathcal{D}$ to respectively
indicate if the traded sector is modeled as continuous or as discrete.\footnote{An alternative approach would be to use industry-level data and treat
select industries as discrete. Under this approach the resulting intuitions
would remain qualitatively the same.} 

When applying Equation \ref{eq: GT} under the presumption that the
traded sector is continuous ($k=\mathcal{C}$), $\lambda_{ii,\mathcal{C}}$
can be calculated as the \emph{value }share of domestic goods in tradable
consumption. $\hat{N}_{i,\mathcal{C}}$ can be derived using the free
entry condition. In particular, free entry implies that $N_{i,\mathcal{C}}=r_{i,\mathcal{C}}L_{i}/(1+\theta)f^{e}$,
with $r_{i,\mathcal{C}}$ denoting the share of economy $i$'s total
revenue generated in the traded sector. From this expression we can
conclude that $\hat{N}_{i,\mathcal{C}}=1$ since $r_{i,\mathcal{C}}^{\textrm{A}}=r_{i,\mathcal{C}}=e_{i,\mathcal{C}}$,
where $e_{i,\mathcal{C}}$ denotes the share of expenditure on traded
goods.\footnote{$r_{i,\mathcal{C}}^{\textrm{A}}=r_{i,\mathcal{C}}=e_{i,\mathcal{C}}$
results from assuming one traded sector and balanced trade.} Plugging $\hat{N}_{i,\mathcal{C}}=1$ back in to Equation \ref{eq: GT},
yields the following formula for the gains from trade under the continuous
interpretation of the economy:
\begin{equation}
GT_{i}^{\textrm{ACR}}=1-\lambda_{ii,\mathcal{C}}^{\frac{e_{i,\mathcal{C}}}{\theta}}.\label{eq: GT CES}
\end{equation}
In the other extreme where the traded sector is treated as entirely
discrete ($\mathcal{D}$), $\lambda_{ii,\mathcal{D}}$ can be calculated
as the \emph{quantity }share of domestic goods in tradable consumption.
$\hat{N}_{i,\mathcal{D}}$ can be once again derived from the free
entry condition. As demonstrated in Appendix \ref{subsec: GT Derivation},
free entry implies that $N_{i,\mathcal{D}}=\sum_{\ell\in\mathbb{C}}\left(\lambda_{i\ell,\mathcal{D}}Y_{\ell}/Y_{i}\right)L_{i}/\theta f^{e}$
and, correspondingly, $N_{i,\mathcal{D}}^{\textrm{A}}=L_{i}/\theta f^{e}$.\footnote{These expressions are derived under the implicit assumption that productivity
in the traded sector is sufficiently high, so that domestic discrete
varieties are affordable to local consumers irrespective of international
trade costs. This assumption ensures that the demand function facing
domestic varieties is continuous and well-defined.} Combining these two expressions, $\hat{N}_{i,\mathcal{D}}=N_{i,\mathcal{D}}/N_{i,\mathcal{D}}^{\textrm{A}}$
is given by
\begin{equation}
\hat{N}_{i,\mathcal{D}}=\sum_{\ell\in\mathbb{C}}\lambda_{i\ell,\mathcal{D}}\frac{Y_{\ell}}{Y_{i}}.\label{eq: scale-driven GT}
\end{equation}
Plugging the above expression back into Equation \ref{eq: GT}, yields
the following formula for the gains from trade under the discrete
interpretation of the economy:
\begin{equation}
GT_{i}^{\textrm{Discrete}}=1-\lambda_{ii,\mathcal{D}}^{\frac{e_{i,\mathcal{D}}}{\theta}}\left(\sum_{\ell\in\mathbb{C}}\lambda_{i\ell,\mathcal{D}}Y_{\ell}/Y_{i}\right)^{-\frac{e_{i,\mathcal{D}}}{\theta}},\label{eq: GT one-sector}
\end{equation}
where $e_{i,\mathcal{D}}$ denotes the share of expenditure on categories
where discrete goods are bundled with non-traded services in the consumption
basket.\footnote{See Appendix \ref{tab:List-of-industries} for a descriptive list
of consumption categories that fit this bill.} 

Evaluating Equation \ref{eq: GT CES} is rather straightforward, but
mapping Equation \ref{eq: GT one-sector} to data involves a couple
of auxiliary steps. First, we need to evaluate $\lambda_{ii,\mathcal{D}}$,
which corresponds to the \emph{quantity} share of domestic consumption.
The \noun{Comtrade--baci} database reports import quantities, but
measuring $\lambda_{ii,\mathcal{D}}$ also requires information on
the quantity of domestic consumption. One way to overcome this measurement
challenge is to infer $\lambda_{ii,\mathcal{D}}$ from the domestic
expenditure shares and prices. To demonstrate this, recall that $X_{ji}$
denotes country $i$'s total spending on tradable varieties from country
$j$, with $p_{ji}$ denoting the corresponding price of these varieties.
Given information on $\{X_{ji}\}$ and $\{p_{ji}\}$, we can express
$\lambda_{ii,\mathcal{D}}=Q_{ii}/\sum_{j}Q_{ji}$ as follows
\[
\lambda_{ii,\mathcal{D}}=\frac{X_{ii}/p_{ii}}{X_{ii}/p_{ii}+\sum_{j\neq i}X_{ji}/p_{ji}}=\frac{X_{ii}/X_{i}}{X_{ii}/X_{i}+\left(1-X_{ii}/X_{i}\right)p_{ii}/\overline{p}_{-ii}},
\]
where $X_{i}=\sum_{j}X_{ji}$, while $\overline{p}_{-ii}\equiv\sum_{j\neq i}\frac{X_{ji}}{X_{i}}p_{ji}$
denotes the (trade-weighted average) unit price of country $i$'s
imports from the rest of the world. To evaluate the above equation,
we need data data on $p_{ii}/\overline{p}_{-ii}$ and $\{X_{ji}\}$.
Data on the former variable is obtainable from the Penn World Tables
(PWT), which report aggregate price levels for imports and domestic
expenditure. The data on tradable expenditure values, $\{X_{ji}\}$
are taken from the 2008 version of the World Input-Output Database
(\noun{WIOD}, \citet{timmer2012world}).\footnote{To make the data consistent with theory, I purge it from trade imbalances
following the methodology in \citet{costinot2014trade}. } Combining the aforementioned data, I can measure $\lambda_{ii,\mathcal{D}}$
for 31 countries. Likewise, measures for $e_{i,\mathcal{C}}$ and
$e_{i,\mathcal{D}}$ can be constructed with data on $\{X_{ji}\}$
by grouping industries into traded and non-traded categories---see
Table \ref{tab:List-of-industries} in the appendix for details.\footnote{Table \ref{tab:GT} reports the corresponding values for $e_{i,\mathcal{C}}$
and $e_{i,\mathcal{D}}$ for each of the 31 countries. Unlike the
continuous model that is calibrated to the share of tradable expenditure,
the discrete model is calibrated to $e_{i,\mathcal{D}}$, which is
the sum of expenditure on discrete tradables and matching services.
In Appendix \ref{sec: Construction-of-Figure}, I check if the discrete
model can replicate the (out-of-sample) share of tradable expenditure
in isolation. Reassuringly, the model performs relatively well on
this front.}

The second term in Equation \ref{eq: GT one-sector}, $\left(\sum_{\ell\in\mathbb{C}}\lambda_{i\ell,\mathcal{D}}Y_{\ell}/Y_{i}\right)^{-\frac{e_{i,\mathcal{D}}}{\theta}},$
is computed using bilateral trade quantities, $Q_{j\ell}$, from the
\noun{Comtrade--baci} database described in Section \ref{sec:Empirical-Evidence}.
Specifically, the export quantity share $\lambda_{i\ell,\mathcal{D}}$
is computed as,
\[
\lambda_{i\ell,\mathcal{D}}=\frac{Q_{i\ell}}{\sum_{j\neq\ell}Q_{j\ell}}\left(1-\lambda_{\ell\ell,\mathcal{D}}\right).
\]
Data on total domestic absorption, $Y_{i}$, are taken from the PWT.
Plugging values for $\lambda_{\ell i,\mathcal{D}}$ and $Y_{i}$ into
Equation \ref{eq: GT one-sector} and combining it with the previously-constructed
data for $\lambda_{ii,\mathcal{D}}$, allows me to compute the gains
from discrete trade for each of the 31 countries in the 2008 WIOD
sample. For the same set of countries, I can compute the gains implied
by the ACR formula by plugging data values for $e_{i,\mathcal{C}}$,
$e_{i,\mathcal{D}}$, and $\lambda_{ii,\mathcal{C}}=X_{ii}/X_{i}$
into Equation \ref{eq: GT CES}. 

\paragraph*{Estimating $\theta$.}

I estimate $\theta$ using the triple difference approach developed
by \citet{caliendo2014estimates}. This approach is fully described
in Appendix \ref{sec:Estimation-of-theta}, and requires data on bilateral
trade shares and tariff rates. Trade shares are taken directly from
the WIOD. Data on bilateral tariff rates are taken from the United
Nations Statistical Division, Trade Analysis and Information System
(UNCTAD-TRAINS, see Appendix \ref{sec:Estimation-of-theta}). The
estimated $\theta$ is reported in Table \ref{tab:Estimated theta}.
The main difference between the continuous and discrete cases is that
under the former the estimating equation reduces to the familiar log-log
specification implied by CES-based models. Under the latter, the estimating
equation adopts a log-level specification---see Appendix \ref{sec:Estimation-of-theta}
for further details. These differences notwithstanding, both approaches
estimate a relatively similar value for $\theta$, which is in line
with exiting estimates from the literature.\footnote{\citet{Broda2006} estimate an import demand elasticity that is analogous
to $\theta+1$. Their product-level estimates exhibit an average value
of 6.} 
\begin{center}
\begin{table}[!tbh]
\caption{\label{tab:Estimated theta}Estimated $\theta$ under discrete and
continuous models }

\centering{}\vspace{-0.1in}
\resizebox{0.8\textwidth}{!}{\begin{tabular}{lccc} 
\toprule 
& Discrete Model & Continuous Model \\ 
\midrule  
Estimated $\theta$ & 4.30 & 3.63 \\
	& (0.53)   &   (0.52) \\ \addlinespace
observations & \multicolumn{2}{c}{163,058} \\
\bottomrule \vspace{-0.2in}
\end{tabular}}
\noindent\begin{minipage}[t]{1\columnwidth} \begin{spacing}{0.8} {\footnotesize \emph{Note:} This table reports the estimated $\theta$ using the triple differences methodology described in Appendix \ref{sec:Estimation-of-theta}. The estimation uses (a) pooled data on trade shares from the 2008 version of the WIOD that span 15 traded industries,  and (b) matched tariff data from the UNCTAD-TRAINS database. The first column treats traded industries as discrete. The second columns treats traded industries as continuous. Standard errors are reported in parentheses.}
\end{spacing} \end{minipage}
\end{table}
\vspace{-0.5in}
\par\end{center}

\subsection{Quantitative Results}

The computed gains from trade are reported in Table \ref{tab:GT}.
In summary, the discrete interpretation of the data predicts gains
from trade that are on average larger (7.6\% versus $5.9\%$) and
more heterogeneous than those implied by the continuous interpretation.
These differences are driven by two factors: First, from the view
of the discrete model, the gains from trade can be large even if a
country spends a small fraction of its income on tradable goods. To
elaborate, consider the US and Mexican economies that spend a similar
fraction of their income on categories where traded goods are bundled
with services. The US economy, though, spends only 27\% of its total
expenditure on \emph{specifically} traded goods compared to 44\% in
Mexico. The continuous interpretation of the economy attributes this
pattern to the US consumers assigning a lower weight to traded goods
in their preferences. This interpretation inevitably predicts relatively
small gains for the US. The discrete model, however, attributes the
same pattern to the income-inelastic nature of discrete goods and
the higher purchasing power of US consumers relative to Mexican Consumers. 

Second, the discrete model admits \emph{scale-driven} gains that are
highly heterogeneous across countries. They favor small countries
(e.g., Belgium and Canada) or low-income countries (e.g., Brazil and
Indonesia) that are located in the proximity of larger and richer
markets . For instance, the ACR formula predicts that trade increases
Brazil's real GDP by only 1.8\%. The discrete model, meanwhile, predicts
gains of around 6.3\% for Brazil, of which 3.6\% is purely scale-driven.
For Canada, the ACR formula predicts gains of around 5.4\%. The discrete
model predicts 15\% in gains, of which 5.6\% is scale-driven. As I
will elaborate below, these gains are driven by firm-relocation into
these countries from larger and richer trading partners like the US.
Correspondingly, the scale-driven gains from trade are negative for
the US and stand around $-3\%$. 

\begin{table}[!htbp]\centering \caption{The gains from trade}\label{tab:GT}
\vspace{-0.1in}
\begin{tabular}{lcccccccccc} \toprule & \multicolumn{4}{c}{} & \multicolumn{1}{c}{} && \multicolumn{2}{c}{$GT_i$ Discrete} \\ \addlinespace[3pt] \cline{8-9} \addlinespace[3pt] & \specialcell{$e_{i,\mathcal{C}}$} & \specialcell{$e_{i,\mathcal{D}}$} & \specialcell{GDP} & \specialcell{$L_i$} & \specialcell{$GT_i$ Continuous} && \specialcell{Total} & \specialcell{Scale-Driven} \\ \midrule Country & \multicolumn{1}{c}{(1)} & \multicolumn{1}{c}{(2)} & \multicolumn{1}{c}{(3)} & \multicolumn{1}{c}{(4)} & \multicolumn{1}{c}{(5)} && \multicolumn{1}{c}{(6)} & \multicolumn{1}{c}{(7)} \\ \midrule AUS & 0.25 & 0.56 & 0.07 & 0.07 & 3.1\% && 8.4\% & 3.1\% \\AUT & 0.32 & 0.64 & 0.03 & 0.03 & 8.1\% && 8.2\% & -5.2\% \\BEL & 0.32 & 0.63 & 0.04 & 0.04 & 12.7\% && 25.0\% & 4.1\% \\BRA & 0.42 & 0.76 & 0.63 & 0.11 & 1.8\% && 6.3\% & 3.6\% \\CAN & 0.31 & 0.68 & 0.11 & 0.10 & 5.4\% && 15.0\% & 5.6\% \\CHN & 0.63 & 0.79 & 4.36 & 0.31 & 3.2\% && 2.5\% & -0.8\% \\CZE & 0.42 & 0.69 & 0.03 & 0.02 & 8.6\% && 13.0\% & 0.6\% \\DEU & 0.35 & 0.67 & 0.27 & 0.26 & 6.3\% && 5.3\% & -4.6\% \\DNK & 0.27 & 0.63 & 0.02 & 0.02 & 7.3\% && 8.7\% & -5.5\% \\ESP & 0.31 & 0.60 & 0.15 & 0.11 & 4.1\% && 3.1\% & -3.8\% \\FIN & 0.35 & 0.64 & 0.02 & 0.02 & 4.9\% && 7.2\% & -0.6\% \\FRA & 0.28 & 0.59 & 0.21 & 0.20 & 4.2\% && 3.3\% & -4.2\% \\GBR & 0.21 & 0.62 & 0.20 & 0.19 & 4.4\% && 4.0\% & -6.3\% \\GRC & 0.33 & 0.65 & 0.04 & 0.02 & 5.3\% && 2.1\% & -6.9\% \\HUN & 0.45 & 0.73 & 0.03 & 0.01 & 11.5\% && 12.4\% & -4.1\% \\IDN & 0.51 & 0.73 & 0.77 & 0.03 & 3.5\% && 8.5\% & 4.3\% \\IND & 0.51 & 0.76 & 3.86 & 0.08 & 3.2\% && 2.7\% & -1.2\% \\IRL & 0.27 & 0.58 & 0.01 & 0.02 & 5.8\% && 5.0\% & -5.5\% \\ITA & 0.34 & 0.65 & 0.19 & 0.16 & 3.7\% && 2.6\% & -3.5\% \\JPN & 0.36 & 0.66 & 0.42 & 0.33 & 2.2\% && 0.6\% & -2.5\% \\KOR & 0.52 & 0.77 & 0.16 & 0.07 & 5.4\% && 3.3\% & -3.3\% \\MEX & 0.44 & 0.70 & 0.38 & 0.07 & 4.9\% && 6.8\% & 0.1\% \\NLD & 0.27 & 0.62 & 0.05 & 0.06 & 8.8\% && 16.0\% & -0.8\% \\POL & 0.39 & 0.69 & 0.13 & 0.04 & 6.0\% && 7.5\% & -1.9\% \\PRT & 0.33 & 0.65 & 0.03 & 0.02 & 6.4\% && 5.9\% & -5.2\% \\RUS & 0.42 & 0.69 & 0.47 & 0.11 & 5.9\% && 9.1\% & 5.1\% \\SVK & 0.41 & 0.71 & 0.02 & 0.01 & 11.5\% && 17.5\% & 0.0\% \\SVN & 0.38 & 0.64 & 0.01 & 0.00 & 11.4\% && 12.5\% & -4.2\% \\SWE & 0.30 & 0.64 & 0.03 & 0.03 & 6.0\% && 8.8\% & -2.3\% \\TUR & 0.42 & 0.73 & 0.23 & 0.05 & 3.9\% && 3.0\% & -2.9\% \\USA & 0.27 & 0.67 & 1.00 & 1.00 & 2.3\% && 1.5\% & -3.0\% \\ \bottomrule \addlinespace 
\end{tabular} 
\noindent\begin{minipage}[t]{1\columnwidth} \begin{spacing}{0.8} {\footnotesize Note: $e_{i,\mathcal{C}}$ denotes the share of expenditure on traded goods in country $i$. $e_{i,\mathcal{D}}$ denotes the share of expenditure on categories where traded goods are bundled with matching services. See Table \ref{tab:List-of-industries} in the appendix for details.} \end{spacing} \end{minipage}
\end{table}

\subsection{Elucidating the Scale-Driven Gains from Discrete Trade }

The scale-driven gains from discrete trade operate through a unique
channel that is akin to a \emph{reverse home-market effect}. After
opening to trade, firms relocate from richer and larger economies
that host their main costumer-base to lower-income countries. Doings
so exposes them to higher trade costs but allows them to collect a
higher markup over marginal cost. To better understand these effects,
consider a world economy with two countries: $i=1,2$. Suppose the
two countries have identical productivity levels and face similar
trade costs, but differ in their population sizes:
\[
a_{1,\mathcal{D}}=a_{2,\mathcal{D}}=a;\ \ \ \ \tilde{a}_{1}=\tilde{a}_{2}=1;\ \ \ \ \tau_{12}=\tau_{21}=\tau>1;\ \ \ \ L_{2}>L_{1}.
\]
In this basic setup, the free entry condition implies that (see Appendix
\ref{subsec: GT Derivation})
\[
N_{i,\mathcal{D}}=\sum_{\ell=1,2}\left(\lambda_{i\ell,\mathcal{D}}\frac{w_{\ell}L_{\ell}}{w_{i}L_{i}}\right)\frac{L_{i}}{\theta f^{e}},\ \ \ \ i=1,2.\ \ \ \ (\textrm{Free Entry})
\]
Moreover, given that $p_{ij,\mathcal{D}}=\tau aw_{i}+w_{j}/\theta$,
the balanced trade condition can be expressed as follows:
\[
\left(\tau a\frac{w_{2}}{w_{1}}+\frac{1}{\theta}\right)\lambda_{21,\mathcal{D}}w_{1}L_{1}=\left(\tau a\frac{w_{1}}{w_{2}}+\frac{1}{\theta}\right)\lambda_{12,\mathcal{D}}w_{2}L_{2}.\ \ \ \ \ (\textrm{Balanced Trade})
\]
Combining the balanced trade and free entry conditions yields the
following expression for the number of firm located in country $i=1,2$
(see Appendix \ref{subsec: Equation 17}):
\begin{equation}
N_{i,\mathcal{D}}=\left(1+\frac{\frac{w_{-i}}{w_{i}}-\frac{w_{i}}{w_{-i}}}{\frac{w_{i}}{w_{-i}}+\frac{1}{a\tau\theta}}\lambda_{-ii,\mathcal{D}}\right)\frac{L_{i}}{\theta f^{e}}.\label{eq: N_i,D}
\end{equation}
Since $N_{i}^{\text{A}}=L_{i}/\theta f^{e}$ denotes the mass of firms
under autarky, the above equation states that opening to trade preserves
the total number of firms, globally; but induces firm-relocation from
the high- to low-wage country. The intuition is that by relocating
to the low-wage country, firms can collect a higher multiplicative
markup over marginal cost, as is evident from the optimal price equation,
$p_{ij,\mathcal{D}}=a\tau w_{i}(1+\frac{1}{a\tau\theta}\frac{w_{j}}{w_{i}})$.\footnote{We can see the same effect by reformulating Equation \ref{eq: scale-driven GT}
as follows:
\[
\hat{N}_{i,\mathcal{D}}=\frac{\sum_{\ell\in\mathbb{C}}\left(1-\mu_{i\ell,\mathcal{D}}^{-1}\right)\omega_{i\ell,\mathcal{D}}}{\sum_{\ell\in\mathbb{C}}\left(1-\mu_{\ell i,\mathcal{D}}^{-1}\right)\lambda_{\ell i,\mathcal{D}}},
\]
where $\mu_{i\ell,\mathcal{D}}\equiv p_{ji,\mathcal{D}}/c_{ji,\mathcal{D}}$
denotes the multiplicative markup over marginal cost, and $\omega_{i\ell,\mathcal{D}}$
denotes the share of market $\ell$ in country $i$'s total exports
($\sum_{\ell\in\mathbb{C}}\omega_{i\ell,\mathcal{D}}=1$). } Furthermore, we can invoke the balanced trade condition to show that
the larger economy, in our example, pays a higher wage, i.e., $w_{2}>w_{1}$.
In that case, Equation \ref{eq: N_i,D} implies that opening to trade
induces firm-relocation from the larger to the smaller country:
\[
\hat{N}_{1,\mathcal{D}}>1;\ \ \ \ \hat{N}_{2,\mathcal{D}}<1.
\]
These effects resemble the \emph{``big-push''} effects of economic
development, as formalized by \citet{rosenstein1943problems}, \citet{lewis1954economic},
and \citet{murphy1989industrialization}. To outline this connection,
note that the scale of production in small or poor economies is constrained
by their inadequate local demand. Opening to trade, however, relocates
production activity towards these economies. In doing so, it provides
them with the necessary pig-push platform to expand.

\section{Concluding Remarks}

\emph{``Why nations trade''} and \emph{``how much they gain from
trade''} are perhaps two of the most central questions in the history
of economic thought. Economists have historically approached these
two questions with trade models that treat all goods as infinitely
divisible. This paper asked how our answers to these old questions
are revised if we account for the discrete nature of traded goods.
In this process, I identified two new drivers of international specialization
and pricing-to-market. I further demonstrated that these previously-overlooked
forces are empirically significant. Perhaps most importantly, I illustrated
that accounting for the discrete or lumpy nature of trade can greatly
modify our estimates for the gains from trade.

In the quest to understand discrete trade, this paper only scratched
the surface. Given the fruitful nature of this preliminary step, several
extensions of the discrete model are in order. First, the model provides
a new perspective on the determinants of market power. In that regard,
it can potentially shed light on the documented rise of global market
power in import-competing industries (see \citeauthor{de2017rise}).
Relatedly, the discrete model can provide a fresh perspective on the
trade slowdown and zeros. To elaborate on the latter subject, there
is a withstanding puzzle as to why the incidence of zero trade is
related to income per capita. Appealing to affordability constraints
can perhaps help resolve this puzzle. Finally, in the interest of
space, the present paper resorted to quantifying the aggregate gains
from discrete trade. Applying the discrete trade framework to micro-level
expenditure data presents a promising avenue for feature research
on the distributional consequences of trade.

\begin{singlespace}
\bibliographystyle{jpe}
\bibliography{Discrete}

\end{singlespace}

\newpage{}

\appendix

\part*{Appendix}

\section{\label{sec:The-Firm-Level-Colombian}The Firm-Level Colombian Export
and Import Data}

The firm-level data used in Section \ref{sec:Suggestive-Evidence}
covers the universe of Colombian import and export transactions for
the 2007--2013 period. The data has been collected and made available
by the National Tax Agency. For each import transaction, it identifies
the exporting and importing firm's id, the 10-digit Harmonized System
(HS10) classification to which the imported or exported goods belong,
as well as the f.o.b. value in US dollars, quantity, and the unit
in which quantity is reported. 

The import database features 7,296 distinct HS10 product categories,
and 226,288 firms from 251 different countries. Each of these import
transactions is conducted by Colombian firm that is identified by
their unique 9-digits tax id. In total there are 95,071 of such importers,
most of which import a wide range of discrete and non-discrete HS10
products. The export database features 6,590 distinct HS10 product
categories sold to 231 different markets. Each of these export transactions
is conducted by Colombian firm that is identified by their unique
9-digits tax id. In total there are 33,075 exporting firms in the
database, many of which export a wide range of discrete and non-discrete
HS10 products to various markets.

\begin{onehalfspace}
The Colombia import and export data report quantity in 10 different
units. The vast majority of entries report quantity either in counts,
\emph{``UNIDADES O ARTICULOS''}, or in kilograms, \emph{``KILOGRAMO.''}
72\% of all observations in the import data (51\% in value terms)
involve goods that report quantity in counts. In comparison, 63\%
of all observations in the export data (12\% in value terms) involve
goods that report quantity in counts. I classify such goods are \emph{discrete}
and classify all other goods as \emph{continuous}.
\end{onehalfspace}

\section{\label{sec: Non-Parametric-Model}Non-Parametric Model}

In this appendix, I present a general non-parametric version of the
model, which features the following elements: $(i)$ the demand-side
is governed by a general non-parametric discrete choice framework
that can reproduce an important class of continuous and discrete demand
functions, $(ii)$ income heterogeneity is governed by a non-parametric
endowment distribution, $G_{i}(y)$, and $(iii)$ the supply-side
is identical to the baseline model.

\subsection*{General Discrete Choice Preference Structure }

Demand in all industries is characterized by a general discrete choice
framework that can reproduce an important class of well-known demand
systems. As in the baseline model, the economy is populated with $k=1,...,K$
industries, with $e_{k}$ denoting the constant share of expenditure
on industry $k$. Unless stated otherwise, the notation is similar
to the one used in Section \ref{sec:Theoretical-Framework}. The utility
individual $\iota$ extracts from choosing variety $\omega\in\Omega_{j,k}(\varphi)$
and consuming $q_{ji,k}(\omega)$ units of it plus $\tilde{q}$ hours
of industry $k$-related services is given by:
\[
U_{\iota}(q(\omega),\tilde{q};\varphi)=\mathcal{U}_{k}(\tilde{q},q(\omega);\varphi)\ +\ \xi_{\iota}(\omega),\ \ \ \omega\in\Omega_{j,k}(\varphi,y)
\]
where $\xi_{\iota}(\omega)$ is an \emph{individual$\times$variety}-specific
utility shifter that accounts for individual $\iota$'s personal taste
for variety $\omega$. $\Omega_{j,k}(\varphi,y)\subset\Omega_{j,k}(\varphi)$
denotes the subset of varieties affordable to an individual with endowment
$y$. An individual residing in country $i$ with endowed $y_{\iota}$
has a budget equal to $e_{k}y_{\iota}w_{i}$. The discrete choice
problem facing them, therefore, can be stated as
\begin{align*}
\max_{q(\omega),\tilde{q}}\ \ \  & \mathcal{U}_{k}(\tilde{q},q(\omega);\varphi)+\xi_{\iota}(\omega)\\
\textrm{s.t.}\ \ \  & \tilde{p}_{i,k}\ \tilde{q}+p_{ji,k}(\omega)\ q(\omega)\leq e_{k}y_{\iota}w_{i}\ .
\end{align*}
 The above problem can be solved in two separate steps. First, conditional
on choice $\omega$, the optimal choice for $q(\omega)$ and $\tilde{q}$
can be attained by maximizing $\mathcal{U}_{k}(.)$ subject to the
budget constraint. This stage delivers an indirect utility function
$V_{k}\left(yw_{i},\tilde{p}_{i,k},p_{ji,k}(\omega)\right)$, which
can be used to solve the second stage, which is choosing the optimal
variety $\omega$ given the menu of prices and individual taste. Namely,
\begin{equation}
\max_{\omega}\ \ \ V_{k}\left(yw_{i},\tilde{p}_{i,k},p_{ji,k}(\omega);\varphi\right)\ +\ \xi_{\iota}(\omega).\label{eq: indirect utility}
\end{equation}
To simplify the notation, I define $\mathcal{V}_{k}\left(.\right)\equiv\exp(V_{k}(.))$.
Moreover, to decouple the role of \emph{horizontal} and \emph{vertical}
differentiations, I assume that $\mathcal{V}_{k}(.)$ is separable
in quality:
\[
\mathcal{V}_{k}\left(yw_{i},\tilde{p}_{i,k},p_{ji,k}(\omega);\varphi\right)=v_{k}\left(yw_{i},\tilde{p}_{i,k},p_{ji,k}(\omega)\right)\ f_{k}(\varphi)
\]
As in the baseline model, $\boldsymbol{\xi}_{\iota,k}\equiv\left\{ \xi_{\iota}(\omega)\right\} $
is drawn independently from a General Extreme Value distribution,
$F_{k}(\boldsymbol{\xi}_{k})=\exp\left(-\sum_{\varphi\in\varPhi_{k}}\left(\sum_{\omega}\exp(\xi(\omega))^{-\theta_{k}}\right)^{\gamma_{k}/\theta_{k}}\right)$,
where $\theta_{k}\geq\gamma_{k}\geq0$. With this assumption, we can
invoke the theorem of \emph{General Extreme Value} (\citet{McFadden1978,ben1985discrete})
to produce the following lemma about \emph{individual-level} and \emph{market-level}
demand functions.\footnote{Proof provided in Appendix \ref{sec:Proof-of-Lemma 1}.} 
\begin{lem}
The probability of choosing variety $\omega\in\Omega_{i,k}\left(\varphi,y\right)$
by an individual with is
\[
\lambda_{k}\left(y,w_{i},p_{ji,k}(\omega),\varphi,\tilde{p}_{i,k},\boldsymbol{P}_{i}(y)\right)=\mathcal{V}_{k}\left(yw_{i},p_{ji,k}(\omega),\widetilde{p}_{i,k};\varphi\right)^{\theta_{k}}\ \frac{P_{i,k}\left(\varphi,y\right)^{\theta_{k}-\gamma_{k}}}{\sum_{\varphi'\in\varPhi_{k}}P_{i,k}\left(\varphi',y\right)^{-\gamma_{k}}},
\]
where $P_{i,k}\left(\varphi,y\right)\equiv\left[\sum_{j\in\mathbb{C}}\sum_{\omega\in\Omega_{i,k}\left(\varphi,y\right)}\mathcal{V}_{k}\left(yw_{i},p_{ji,k}(\omega),\widetilde{p}_{i,k};\varphi\right)^{\theta_{k}}\right]^{\frac{-1}{\theta_{k}}}$is
an income- and quality-specific demand shifter. Accordingly, the market-level
Marshallian demand facing variety $\omega\in\Omega_{i,k}\left(\varphi,y\right)$
can be expressed as\footnote{In order to make the notation more compact, I use $\boldsymbol{P}_{i}(y)$
instead $\boldsymbol{P}_{i}(yw_{i})$.}
\[
\hspace{-0.25in}q_{k}\left(w_{i},p_{ji,k}(\omega),\varphi,\tilde{p}_{i,k},\boldsymbol{P}_{i,k},G_{i}\right)=\left[\int_{y}\mathcal{D}_{k}\left(yw_{i},p_{ji,k}(\omega),\varphi,\tilde{p}_{i,k}\right)\lambda_{k}\left(y,p_{ji,k}(\omega),\varphi,\tilde{p}_{i,k},\boldsymbol{P}_{i}(y)\right)dG_{i}\left(y\right)\right]L_{i},
\]
where $\mathcal{D}_{k}\left(.\right)$ is the demand function associated
with $V_{k}\left(.\right)$.
\end{lem}
In the above lemma, $\boldsymbol{P}_{i,k}(y)=\left\{ P_{i,k}\left(\varphi,y\right)\right\} $
and $\boldsymbol{P}_{i,k}=\left\{ \boldsymbol{P}_{i,k}(y)\right\} $
denote the vector of demand shifters for various quality levels and
income groups. Likewise, I use $\boldsymbol{\Omega}_{i,k}\left(y\right)\equiv\left\{ \Omega_{i,k}\left(\varphi,y\right)\right\} _{\varphi\in\varPhi_{k}}$
to denote the complete set of varieties (from all quality levels)
affordable to individuals with endowment $y$ in market $i$. 

At this point, it is useful to formally define homothetic preferences
from the lens of the model. I say that individual-level preferences
are homothetic iff
\[
\partial\lambda_{k}\left(y,w_{i},p_{ji,k}(\omega),\varphi,\tilde{p}_{i,k},\boldsymbol{P}_{i}(y)\right)/\partial y=0,
\]
provided that $\partial$$\boldsymbol{\Omega}_{i,k}\left(y\right)/\partial y=0$.
The above definition outlines two sources of income-driven variation
in individual-level demand. First, conditional on choosing from the
same set of affordable varieties, $\lambda_{k}\left(.\right)$ may
vary with $y$ due to the functional-form of the indirect utility
function, $\mathcal{V}_{k}\left(.\right)$. Second, $\lambda_{k}\left(.\right)$
may vary with $y$ due to the expansion of set $\boldsymbol{\Omega}_{i,k}\left(y\right)$
with $y$.\footnote{Following the chain rule, if $\lambda_{k}(.)$ is invariant to income,
so is the individual-level demand elasticity, $\epsilon_{k}\left(.\right)\equiv\partial\ln\lambda_{k}\left(.\right)/\partial\ln p_{ji,k}(\omega)$.
Likewise, if preferences are homothetic (as defined above) the share
of individuals that chose a variety with quality $\varphi$, namely,
$\Lambda_{i,k}\left(y,\varphi\right)\equiv P_{i,k}\left(\varphi,y\right)^{-\gamma_{k}}/\sum_{\varphi'\in\varPhi_{k}}P_{i,k}\left(\varphi',y\right)^{-\gamma_{k}}$,
does not vary with $y$. } In what follows, I apply the demand system characterized by Lemma
1 to two special cases: $(i)$ continuous industries, and $(ii)$
discrete industries. 

\subsubsection*{(a) Continuous Industries}

Continuous industries involve goods that are $(i)$ infinitely divisible,
and $(ii)$ not directly substitutable with services. As a result,
all continuous goods are affordable to all individuals by assumption,
i.e., $\Omega_{i,k}(\varphi,y)=\Omega_{i,k}(\varphi)$ for all $y$.
I implement these features in my general discrete choice model by
assuming that individuals spend their entire budget on buying multiple
units of their preferred variety: $q(\omega)=e_{k}yw_{i}/p_{ji,k}(\omega)$.
As we will see shortly, at the market-level, this choice structure
is isomorphic to individuals having \emph{indirectly} additive preferences
and buying a fractional quantity of multiple varieties. The indirect
utility, therefore, can be formulated as\footnote{To be clear, the price of non-traded services does not enter the indirect
utility function, because of the implicit assumption that continuous
goods are not directly substitutable with services.}
\[
\mathcal{V}_{k}\left(.\right)=v_{k}\left(e_{k}yw_{i}/p_{ji,k}(\omega)\right)\ f_{k}\left(\varphi\right),\ \ \ \omega\in\Omega_{k}\left(\varphi\right).
\]
Plugging the above characterization of $\mathcal{V}_{k}\left(.\right)$
into Lemma 1, the probability that a consumer with income $y$ in
country $i$ chooses variety $\omega\in\Omega_{j,k}\left(\varphi\right)$
is given by
\[
\lambda_{k}\left(y,e_{k}w_{i}/p_{ji,k}(\omega),\varphi,\boldsymbol{P}_{i,k}(y)\right)=v_{k}\left(ye_{k}w_{i}/p_{ji,k}(\omega)\right)^{\theta_{k}}\ \frac{P_{i,k}\left(\varphi,y\right)^{\theta_{k}-\gamma_{k}}}{\sum_{\varphi'\in\varPhi_{k}}P_{i,k}\left(\varphi',y\right)^{-\gamma_{k}}},
\]
where $P_{i,k}\left(\varphi,y\right)=f_{k}\left(\varphi\right)\left[\sum_{j\in\mathbb{C}}\sum_{\omega\in\Omega_{j,k}\left(\varphi\right)}v_{k}\left(ye_{k}w_{i}/p_{ji,k}(\omega)\right)^{\theta_{k}}\right]^{\frac{-1}{\theta_{k}}}$.
Another characteristic of continuous industries is that all varieties
are accessible to all consumers---i.e., $\Omega_{i,k}\left(y,\varphi\right)=\Omega_{i,k}\left(\varphi\right)$
for all $y$. Correspondingly, the market-level demand facing firm
$\omega$ in market $i$ can be calculated as
\begin{equation}
q_{k}\left(e_{k}w_{i}/p_{ji,k}(\omega),\varphi,\boldsymbol{P}_{i,k},G_{i}\right)=\left[\int_{y}\left(ye_{k}w_{i}/p_{ji,k}(\omega)\right)\lambda_{k}\left(y,e_{k}w_{i}/p_{ji,k}(\omega),\varphi,\boldsymbol{P}_{i,k}(y)\right)dG_{i}\left(y\right)\right]L_{i},\label{eq: cont. demand}
\end{equation}
where the above expression uses the fact that $\mathcal{D}_{k}\left(.\right)\equiv ye_{k}w_{i}/p_{ji,k}(\omega)$.
The above demand function is identical to that arising from indirectly
additive preferences.\footnote{See \citet{bertoletti2018international} for an elaborate review of
indirectly additive preferences.} This class of preferences includes the CES demand system presented
in Section \ref{sec:Theoretical-Framework} as a special case, where
$v_{k}(x)=x$. The above demand function immediately characterizes
the market-level price elasticity of demand, both in the general case
and in the special case where preferences are homothetic (i.e., CES).
\begin{lem}
{[}\textbf{Demand Elasticity in Continuous Industries}{]}\\
$(a)$ The market-level Marshallian demand elasticity facing firm
$\omega\in\Omega_{i,k}(\varphi)$ depends primarily on the average
consumption level in market $i$ (i.e., $e_{k}w_{i}/p_{ji,k}(\omega)$):
\[
\varepsilon_{ji,k}(\omega)=\varepsilon_{k}\left(e_{k}w_{i}/p_{ji,k}(\omega),\varphi;\boldsymbol{P}_{i,k},G_{i}\right)\ \ \ ,k\in\mathbb{K}_{\mathcal{C}}.
\]
$(b)$ In the case where preferences are homothetic (i.e., $\mbox{x\ensuremath{v_{k}}'\ensuremath{\left(x\right)}/\ensuremath{v_{k}\left(x\right)}=1}$),
the market-level Marshallian demand elasticity is constant: $\varepsilon_{ji,k}(\omega)=1+\theta_{k}$.\footnote{More generally, individual-level preferences are invariant to income
if $xv_{k}'\left(x\right)/v_{k}\left(x\right)=constant_{k}$; but
here to simplify the exposition focus on the case where $xv_{k}'\left(x\right)/v_{k}\left(x\right)=1$.}
\end{lem}
A formal proof for the above lemma is provided in Appendix \ref{sec:Proof-Lemma2}.
But the basic intuition behind Lemma 2 can be provided by focusing
on the case without income heterogeneity. In this special case, the
market-level demand elasticity is given by 
\[
\varepsilon_{k}\left(\bar{q}_{ji,k}(\omega),\varphi\right)=1+\theta_{k}\frac{\bar{q}_{ji,k}(\omega)\ v_{k}'\left(\bar{q}_{ji,k}(\omega)\right)}{v_{k}\left(\bar{q}_{ji,k}(\omega)\right)}.
\]
where $\bar{q}_{ji,k}(\omega)\equiv e_{k}w_{i}/p_{ji,k}(\omega)$.
The above formula implies that the main source of variation in the
market-level demand elasticity is the level of consumption. Moreover,
$\varepsilon_{k}\left(.\right)$ is strictly increasing in the average
consumption level, $\bar{q}$, as long as individual-level demand
satisfies the sub-convexity assumption, $\frac{\partial}{\partial q}\frac{qv'_{k}\left(q\right)}{v_{k}\left(q\right)}<0$,
which is also known as Marshall's second law of demand (\citet{mrazova2017not}).
The above property is at the core of many existing theories of PTM
and IQS. That is, as individuals become richer and consume more, they
become less price-sensitive. Correspondingly, firms charge higher
markups in markets that are populated with high-income, less price-sensitive
individuals. The special case where preferences are homothetic (i.e.,
$qv_{k}'\left(q\right)/v_{k}\left(q\right)=1$), is isomorphic to
the CES model. The fact that the homotheticity and constant elasticity
of market-level demand coincide, is a basic corollary of Bergson\textquoteright s
theorem. This result is significant, as it states that the market-level
demand for continuous goods varies across low- and high-wage markets
only and only if individual-level preferences are non-homothetic.\footnote{There is another well-known class of continuous demand functions,
popularized by \citet{arkolakiselusive} (ACDR), where the market-level
demand elasticity varies not directly as a function of individual
consumption, but as function of price $p_{ji,k}(\omega)$ relative
to a market-level choke price. The ACDR class of preferences is closely
related to the indirectly additive preferences nested by the present
model. However, to formally reproduce the ACRD class of demand systems
with a discrete choice framework, one has to impose stronger functional
form assumptions on $\mathcal{V}_{k}\left(.\right)$---see \citet{thisse2016can}
who reproduce the ACDR demand system using a multinomial logit framework.}

\subsubsection*{(b) Discrete Industries}

Discrete industries involve goods that are $(i)$ indivisible, and
$(ii)$ substitutable with services. I implement these features in
my general discrete choice model by assuming that individuals acquire
\emph{one} unit of their preferred discrete variety, and spend the
rest of their budget on corresponding services. They can also forgo
paying an acquisition cost by spending their entire budget on services.\footnote{Theoretically, the choice of forgoing a discrete purchase can be thought
of as follows. There is a choice category $\tilde{\varphi}\in\varPhi_{k}$
(in addition to all existing categories), which consists of homogeneous
varieties (i.e., $\tilde{\theta}\rightarrow\infty$ in this particular
category). Each variety pertaining to this category is priced competitively
at $\tilde{p}_{i,k}$. Hence, choosing a variety from this category
is akin to spending one's entire budget on services.} Considering this, an individual with endowment $y$ who opts for
variety, $\omega\in\Omega_{k}\left(\varphi,y\right)$, purchases one
unit (i.e., $q(\omega)=1$) and is left with $e_{k}yw_{i}-p_{ji,k}(\omega)$
to spare on services. The indirect utility associated with $\mathcal{U}_{k}(.)$
can, therefore, be expressed as
\[
\mathcal{V}_{k}\left(.\right)=v_{k}\left([e_{k}w_{i}y-p_{ji,k}(\omega)]/\tilde{p}_{i,k}\right)\ f_{k}\left(\varphi\right),\ \ \ \omega\in\Omega_{i,k}\left(y,\varphi\right)
\]
Plugging the above formulation into Lemma 1, implies that the individual-level
probability of choosing variety $\omega\in\Omega_{i,k}\left(y,\varphi\right)$
is given by
\begin{equation}
\lambda_{k}\left(y,p_{ji,k}(\omega)/\tilde{p}_{i,k},\varphi,\boldsymbol{P}_{i,k}\left(y\right)\right)=v_{k}\left(\tilde{e}_{i,k}y-\frac{p_{ji,k}(\omega)}{\tilde{p}_{i,k}}\right)^{\theta_{k}}\frac{P_{i,k}\left(\varphi,y\right)^{\theta_{k}-\gamma_{k}}}{\sum_{\varphi'\in\varPhi_{k}}P_{i,k}\left(\varphi',y\right)^{-\gamma_{k}}},\label{eq: Demand}
\end{equation}
where $\tilde{e}_{i,k}\equiv e_{k}/\tilde{a}_{i,k}$, and $P_{i}\left(\varphi,y\right)=f_{k}\left(\varphi\right)\left[\sum_{j\in\mathbb{C}}\sum_{\omega'\in\Omega_{j,k}(y,\varphi)}v_{k}\left(\tilde{e}_{i,k}y-p_{ji,k}(\omega)/\tilde{p}_{i,k}\right)^{\theta_{k}}\right]^{\frac{1}{\theta_{k}}}$.
Noting that $\mathcal{D}_{k}\left(.\right)=1$, the market-level demand
facing variety $\omega$ can, therefore, be expressed as
\begin{equation}
q_{k}\left(p_{ji,k}(\omega)/e_{k}w_{i},p_{ji,k}(\omega)/\tilde{p}_{i,k},\varphi,\boldsymbol{\text{\ensuremath{P}}}_{i,k},G_{i}\right)=\left[\int_{\frac{p_{ji,k}(\omega)}{e_{k}w_{i}}}^{\infty}\lambda_{k}\left(y,p_{ji,k}(\omega)/\tilde{p}_{i,k},\varphi,\boldsymbol{P}_{i,k}\left(y\right)\right)dG_{i}(y)\right]L_{i}.\label{eq: Demand discrete}
\end{equation}
where, unlike the continuous case, the above aggregation takes into
account that variety $\omega$ is affordable only to consumers with
an endowment $y\geq p_{ji,k}(\omega)/e_{k}w_{i}$. 

Before moving forward, it should be noted that the above framework
differs from existing discrete choice trade models (e.g., \citet{verhoogen2008trade}
and \citet{Fajgelbaum2011}) in two key aspects. First, the present
setup allows for discrete goods to be combined with ``non-traded''
services, whereas prior discrete choice models do not feature a non-traded
sector.\footnote{To be specific, the present framework can replicate \citet{verhoogen2008trade}
under the following assumptions: $(i)$ $v(x)=e^{\alpha x}$ and $f\left(x\right)=e^{\beta x}$;
$(ii)$ all varieties are affordable; $(iii)$ services are costlessly
traded across markets, i.e., $\tilde{p}_{i}=1$ for all $i$ (\citet{verhoogen2008trade}
though allows $f\left(.\right)$ to be market-specific). Likewise,
the present framework can replicate the framework in \citet{Fajgelbaum2011}
under the following assumptions: $(i)$ $V\left(y,\tilde{p}_{i},p_{\omega i},\varphi\right)=\exp\left(\varphi\left(y-p_{\omega i}\right)/\tilde{p}_{i}\right)$,
$(ii)$ all varieties are affordable, and $(iii)$ services are costlessly
traded across markets, i.e., $\tilde{p}_{i}=1$ for all $i$. Note,
however, that the framework in \citet{Fajgelbaum2011} does not satisfy
Assumption A1, which as stated earlier is not consequential to any
of our claims that follow.} Second, the present model allows for affordability constraints to
bind. Using Equation \ref{eq: Demand discrete}, I can derive the
market-level demand elasticity both in the general case and in the
homothetic case.
\begin{lem}
{[}\textbf{ Demand Elasticity in Discrete Industries}{]}\\
$(a)$ The market-level Marshallian demand elasticity facing firm
$\omega\in\Omega_{i,k}(\varphi)$ is primarily a function of its overall
affordability (i.e., $p_{ji,k}(\omega)/e_{k}w_{i}$) and its price
relative to non-traded services (i.e., $p_{ji,k}(\omega)/\tilde{p}_{i,k}$)
:%
\begin{comment}
\begin{align*}
\varepsilon\left(p_{\omega}/w_{i},\varphi\right)= & \frac{p_{\omega}}{w_{i}}\frac{\lambda\left(p_{\omega}/w_{i};\varphi\right)}{\lambda\left(\bar{\chi};\varphi\right)}\frac{g_{i}\left(p_{\omega}/w_{i}\right)}{1-G_{i}\left(p_{\omega}/w_{i}\right)}\\
+ & \frac{\theta\frac{p_{\omega}}{w_{i}}\left(\int_{\frac{p_{\omega}}{w_{i}}}^{\infty}\frac{h'\left(\chi-\frac{p_{\omega}}{w_{i}}\right)}{h\left(\chi-\frac{p_{\omega}}{w_{i}}\right)}\lambda\left(p_{\omega}/w_{i},\omega\right)dG_{i}(\chi)\right)}{\int_{\frac{p_{\omega}}{w_{i}}}^{\infty}\lambda\left(p_{\omega}/w_{i},\omega\right)dG_{i}(\chi)}
\end{align*}
\end{comment}
\[
\varepsilon_{ji,k}(\omega)=\varepsilon_{k}\left(p_{ji,k}(\omega)/e_{k}w_{i},p_{ji,k}(\omega)/\tilde{p}_{i,k},\varphi;\boldsymbol{P}_{i},G_{i}\right),\ \ \ k\in\mathbb{K}_{\mathcal{D}}.
\]
$(b)$ In the case where preferences are homothetic (i.e., $v_{k}'\left(x\right)/v_{k}\left(x\right)=1)$,
the market-level demand elasticity is given by 
\[
\varepsilon_{ji,k}(\omega)=\theta_{k}p_{ji,k}(\omega)/\tilde{p}_{i,k}+\rho_{ji,k}(\omega)\mathcal{H}_{i}(p_{ji,k}(\omega)/e_{k}w_{i})p_{ji,k}(\omega)/e_{k}w_{i},\ \ \ k\in\mathbb{K}_{\mathcal{D}}.
\]
where $\rho_{ji,k}(\omega)$ reflects the importance of the marginal
consumers and $\mathcal{H}_{i}\left(.\right)$ denotes the hazard
rate of the income distribution in market $i$.
\end{lem}
The above lemma, which is proven in Appendix \ref{sec: Proof-Lemma3},
indicates that one source of cross-market variation in $\varepsilon_{k}\left(.\right)$
is the functional-form of $v_{k}\left(.\right)$. Specifically, if
$v_{k}'\left(x\right)/v_{k}\left(x\right)\neq1$, preferences will
be non-homothetic and will naturally produce cross-market heterogeneity
in demand. The demand elasticity corresponding to the homothetic case
differs from the formula presented in Section \ref{sec:Theoretical-Framework},
because it is derived under an arbitrary endowment distribution, $G_{i}(y)$.
The empirically-relevant case, though, is one where the income distribution
satisfies the IGFR property.\footnote{IGFR stands for Increasing Generalized Failure Rate, which is satisfied
by nearly all empirically-relevant income distributions (see \citet{lariviere2001selling}
and \citet{lariviere2006note}).} In that case, $\mathcal{H}_{i}(x)x$ is increasing and $\varepsilon_{ji,k}(\omega)$
is lower the more the affordable variety $\omega$. The term $\rho_{ji,k}(\omega)$,
is the purchase probability of variety $\omega$'s marginal consumers
relative to its average consumer. Theoretically, $\rho_{ji,k}(\omega)$
may be greater than 1, even when preferences are homothetic. That
is because the marginal consumer buying $\omega$ may be choosing
from a narrower set of affordable varieties. However, for any variety
$\omega$, we can set $\rho_{ji,k}(p_{ji,k}(\omega))\approx1$ to
a first-order approximation if $(i)$ $v'_{k}\left(.\right)/v_{k}\left(.\right)=1$,
and $(ii)$ $\Omega_{i,k}(\varphi)$ is composed of a discrete set
of varieties.%
\begin{comment}
Lemma 3 outlines two conditions that necessary and sufficient for
homotheticity, which reflects the two distinct drivers of non-homotheticiy
in discrete models. The condition $v'_{k}\left(.\right)/v_{k}\left(.\right)=1$
ensures that conditional on the menu of affordable varieties, the
probability if choosing quality $\varphi$ is independent of income.
However, due to affordability constraints, consumers with different
income levels may have access to a different menu of affordable varieties.
Condition $\gamma_{k}\rightarrow0$ (which is akin to the Cobb-Douglas
assumption), ensures that even then the probability of choosing quality
$\varphi$ is uniform across income levels.
\end{comment}
\begin{comment}
{[}The within-region Engel curve is flatter than the across-region
Engel curve?{]}
\end{comment}


\subsection*{Equilibrium}

Given Lemmas 2 and 3, the demand elasticity facing variety $\omega$
in market $i$ is uniquely determined by $p_{ji,k}(\omega)/w_{i}$
plus market-level demand shifters, $\boldsymbol{P}_{i,k}$. The symmetry
of firms within economy $j$ entails that they all charge the same
price $p_{ji,k}\left(\varphi\right)$ for industry $k$ goods as a
function of quality $\varphi$. Defining $\nu_{ji,k}\left(\varphi\right)\equiv p_{ji,k}\left(\varphi\right)/w_{i}$,
we can express the optimal monopoly pricing equation as follows:\footnote{In the discrete case the existence of a solution is implied by the
fact that $\varepsilon_{k}\left(.\right)$ is increasing in $\nu_{ji,k}(\varphi)$. }
\[
\nu_{ji,k}(\varphi)=\frac{\varepsilon_{k}\left(\nu_{ji,k}(\varphi),\varphi,\boldsymbol{P}_{i,k},G_{i}\right)}{\varepsilon_{k}\left(\nu_{ji,k}(\varphi),\varphi,\boldsymbol{P}_{i,k},G_{i}\right)-1}\frac{\tau_{ji}a_{j,k}c(\varphi)w_{j}}{w_{i}},
\]
where $\boldsymbol{P}_{i,k}=\boldsymbol{P}_{i,k}\left(\boldsymbol{w},\boldsymbol{N}_{k},\boldsymbol{\nu}_{i,k}\right)$
is uniquely determined by $\boldsymbol{\nu}_{i,k}\equiv\left\{ \nu_{ji,k}\left(\varphi\right)\right\} _{\omega\in\Omega_{k}\left(\varphi\right),\varphi\in\varPhi_{k}}$,
the vector of aggregate wage rates, $\boldsymbol{w}=\left\{ w_{i}\right\} _{i}$,
and number of firms $\boldsymbol{N}_{k}=\left\{ N_{i,k}\left(\varphi\right)\right\} _{i\in\mathbb{C},\varphi\in\varPhi_{k}}$.
Similarly, given $\boldsymbol{\nu}_{i,k}$, $\boldsymbol{w}$, and
$\boldsymbol{N}_{k}$ we can uniquely pin down the matrix of industry-level
bilateral trade values as
\begin{equation}
X_{ji,k}\left(\varphi,\boldsymbol{w},\boldsymbol{N}_{k},\boldsymbol{\nu}_{i,k},G_{i}\right)=\nu_{ji,k}\left(\varphi\right)N_{j,k}\left(\varphi\right)\bar{\lambda}_{k}\left(\nu_{ji,k}\left(\varphi\right),\varphi,\boldsymbol{P}_{i,k},G_{i}\right)Y_{i}\label{eq: EQ trade values-1}
\end{equation}
where $\bar{\lambda}_{k}\left(.\right)\equiv q_{k}\left(.\right)/L$
denotes the market-level probability of choosing an imported variety.
Subsequently, given the matrix $\boldsymbol{X}_{k}\equiv\left\{ X_{ji,k}\left(\varphi,\boldsymbol{w},\boldsymbol{N}_{k},\boldsymbol{\nu}_{i,k},G_{i}\right)\right\} _{j,i\in\mathbb{C},\varphi\in\varPhi_{k}}$,
we can calculate the variable profits in each country and industry.

So, altogether, we can define the equilibrium as a vector of normalized
prices $\boldsymbol{\nu}_{i,k}$, wage rates $\boldsymbol{w}$, number
of firms $\boldsymbol{N}_{k}$, and market-specific demand shifters,
$\boldsymbol{P}_{i,k}=\boldsymbol{P}_{i,k}\left(\boldsymbol{w},\boldsymbol{N}_{k},\boldsymbol{\nu}_{i,k}\right)$,
that satisfy $(i)$ the optimal pricing equation, $(ii)$ the balanced
trade condition (BT), and $(iii)$ the zero-profit condition (FE):
\[
\hspace{-0.25in}\begin{cases}
\nu_{ji,k}(\varphi)=\frac{\varepsilon_{k}\left(\nu_{ji,k}(\varphi),\varphi,\boldsymbol{P}_{i,k},G_{i}\right)}{\varepsilon_{k}\left(\nu_{ji,k}(\varphi),\varphi,\boldsymbol{P}_{i,k},G_{i}\right)-1}\frac{\tau_{ji}a_{j,k}c(\varphi)w_{j}}{w_{i}} & \hspace{-0.1in}\forall\varphi\in\varPhi_{k},k\in\mathbb{K}\\
\sum_{j}\sum_{k}\sum_{\varphi\in\varPhi_{k}}X_{ji,k}\left(\varphi,\boldsymbol{w},\boldsymbol{N}_{k},\boldsymbol{\nu}_{i,k},G_{i}\right)=\sum_{\ell\in}\sum_{k}\sum_{\varphi\in\varPhi_{k}}X_{i\ell,k}\left(\varphi,\boldsymbol{w},\boldsymbol{N}_{k},\boldsymbol{\nu}_{i,k},G_{i}\right) & \forall k\in\mathbb{K}\ \ \ (\mathbf{BT})\\
N_{j,k}\left(\varphi\right)=\left\{ \sum_{i}\frac{\nu_{ji,k}\left(\varphi\right)\ \bar{\lambda}_{k}\left(\nu_{ji,k}(\varphi),\varphi,\boldsymbol{P}_{i,k},G_{i}\right)}{\varepsilon_{k}\left(\nu_{ji,k}\left(\varphi\right),\varphi,\boldsymbol{P}_{i,k},G_{i}\right)}\frac{Y_{i}}{Y_{j}}\right\} \frac{L_{j}}{f_{k}^{e}(\varphi)} & \forall k\in\mathbb{K}\ \ \ (\mathbf{FE})
\end{cases},
\]
By Walras' law, the satisfaction of the BT and FE conditions ensures
that labor markets also clear in each country.%
\begin{comment}
{[}\emph{Note to self}: based on the present theory the price of tradables
increase with $w_{i}$ at a lower elasticity than non-tradables.{]}

{[}Remark: if preferences are ``unconditionally'' homothetic, then
aggregate demand is Gorman representable---probably not under conditionally
homothetic because $P\left(\varphi,y\right)$ varies on the extensive
margin!{]}
\end{comment}


\subsection*{IQS and PTM: \emph{Discrete vs. Continuous Industries}}

Recall that, when preferences are homothetic, the demand elasticity
is $(a)$ constant and uniform across markets in continuous industries,
but $(ii)$ variable and heterogeneous across markets in discrete
industries. This outcome immediately implies that we can produce the
following analog of Proposition \ref{prop: prop 1}, using the general
non-parametric model. 
\begin{prop}
$(a)$ Non-homothetic preferences and Ricardian technology differences
are the only drivers of international quality specialization and pricing-to-market
in continuous industries; but $(b)$ in discrete industries, international
quality specialization and pricing-to-market may be alternatively
driven by affordability constraints and cross-national heterogeneity
in the price of non-traded services.
\end{prop}
The proof of the above theorem follows trivially from Lemmas 2 and
3, but needs a basic qualification. There exist a class of continuous
preferences that are \emph{not} indirectly additive, but are homothetic
and imply a non-constant demand elasticity, e.g., Kimball preferences.
In that case, the above proposition should be restated as follows:
PTM may occur under homothetic preferences, but not PTM as a function
of income per capita, which is the focus of this paper.

\section{Proofs and Derivations}

\subsection{\label{sec:Proof-of-Lemma 1}Demand Shares under General Discrete
Choice Preferences}

Suppose the utility of the individual $\iota$ with income $Y$ can
be expressed as
\[
U_{\iota}(\omega)=\ln V_{k}(Y,p_{ji,k}(\omega),\tilde{p}_{i,k})+\xi_{\iota}(\omega),\ \ \ \omega\in\Omega_{j,k}(\varphi)
\]
In this appendix, I invoke the theorem of ``general extreme value''
to characterize the share of consumers who choose variety $\omega$
as their preferred variety. To this end, define $H_{k}(\boldsymbol{V})$
as follows
\[
H_{k}(\boldsymbol{V})=\sum_{\varphi\in\varPhi_{k}}\left(\sum_{j\in\mathbb{C}}\sum_{\omega\in\Omega_{j,k}\left(\varphi,y\right)}A(\varphi)V(\omega)^{\theta_{k}}\right)^{\gamma_{k}/\theta_{k}}
\]
where $\boldsymbol{V}\equiv\{V(\omega)\}$, with $V(\omega)\equiv V_{k}(Y,p_{ji,k}(\omega),\tilde{p}_{i,k})$.
Note that $H_{k}(.)$ is a continuous and differentiable function
of vector $\boldsymbol{V}$ and has the following properties:
\begin{enumerate}
\item $H_{k}(.)\geq0$;
\item $H_{k}(.)$ is a homogeneous function of rank $\theta_{k}$: $H_{k}(\rho\boldsymbol{V})=\rho^{\gamma_{k}}H_{k}(\boldsymbol{V})$;
\item $\lim_{V(\omega)\rightarrow\infty}H_{k}(\boldsymbol{V})=\infty,\ \forall\omega\in\cup_{j\in\mathbb{C}}\Omega_{j,k}$;
\item the $m$'th partial derivative of $H_{k}(.)$ with respect to a generic
combination of $m$ variables $V(\omega)$, is non-negative if $m$
is odd and non-positive if $m$ is even.
\end{enumerate}
Following \citet{ben1985discrete}, if $\xi_{\iota}(\omega)$'s (the
idiosyncratic taste parameters) are drawn from the following distribution
\[
F_{k}(\boldsymbol{\xi})=\exp\left(-H_{k}(\exp(\boldsymbol{\xi}))\right)=\exp\left(-\sum_{\varphi\in\varPhi_{k}}\left(\sum_{j\in\mathbb{C}}\sum_{\omega\in\Omega_{j,k}\left(\varphi,y\right)}A(\varphi)\exp(\xi(\omega))^{\theta_{k}}\right)^{\gamma_{k}/\theta_{k}}\right),
\]
then the probability of choosing variety $\omega\in\Omega_{i,k}\left(\varphi,y\right)$,
namely, $\lambda(\omega)\equiv\lambda\left(Y,\tilde{p}_{i,k},p_{\omega i},\varphi\right)$
is given by
\begin{align*}
\lambda(\omega) & =\frac{V(\omega)\frac{\partial H_{k}(\boldsymbol{V})}{\partial V(\omega)}}{\gamma_{k}H_{k}(\boldsymbol{V})}=\frac{A(\varphi)V(\omega)^{\theta_{k}}\ \left(\sum_{j\in\mathbb{C}}\sum_{\omega'\in\Omega_{j,k}\left(\varphi,y\right)}A(\varphi)V(\omega')^{\theta_{k}}\right)^{\frac{\gamma_{k}}{\theta_{k}}-1}}{\sum_{\varphi'\in\varPhi_{k}}\left(\sum_{\omega'\in\Omega_{i,k}\left(\varphi',y\right)}A(\varphi')V(\omega')^{\theta_{k}}\right)^{\frac{\gamma_{k}}{\theta_{k}}}}\\
 & =\frac{V(\omega)^{\theta_{k}}}{\sum_{j\in\mathbb{C}}\sum_{\omega'\in\Omega_{j,k}\left(\varphi,y\right)}V(\omega')^{\theta_{k}}}\cdot\frac{\left(A(\varphi)\ \sum_{j\in\mathbb{C}}\sum_{\omega'\in\Omega_{j,k}\left(\varphi,y\right)}V(\omega')^{\theta_{k}}\right)^{\frac{\gamma_{k}}{\theta_{k}}}}{\sum_{\varphi'\in\varPhi_{k}}\left(A(\varphi')\ \sum_{j\in\mathbb{C}}\sum_{\omega'\in\Omega_{j,k}\left(\varphi',y\right)}V(\omega')^{\theta_{k}}\right)^{\frac{\gamma_{k}}{\theta_{k}}}}.
\end{align*}
The expression for $q_{k}\left(.\right)$ follows trivially from the
above expression by integrating over all consumers in market $i$. 

\subsection{\label{sec:Proof-Lemma2}Demand Elasticity in Continuous Industries}

This appendix derives the demand elasticity in continuous industries,
which amounts to proving Lemma 2 from Appendix \ref{sec: Non-Parametric-Model}.
Claim (a) in Lemma 2 can be proven by taking the derivative of $q_{k}\left(.\right)$
w.r.t. $p_{ji,k}(\omega)$:
\begin{align*}
\hspace{-1in}\frac{\partial q_{k}\left(p_{ji,k}(\omega),...\right)}{\partial p_{ji,k}(\omega)}= & \frac{\partial}{\partial p_{ji,k}(\omega)}\left[\int_{y}\mathcal{D}_{k}\left(e_{k}yw_{i},p_{ji,k}(\omega)\right)\lambda_{k}\left(y,e_{k}w_{i}/p_{ji,k}(\omega),\varphi,\boldsymbol{P}_{i,k}(y)\right)dG_{i}\left(y\right)\right]L_{i}\\
\hspace{-1in}= & \left[\int_{y}\left(\frac{\partial\mathcal{D}_{k}\left(.\right)}{\partial p_{ji,k}(\omega)}\lambda_{k}\left(y,e_{k}w_{i}/p_{ji,k}(\omega),\varphi;\boldsymbol{P}_{i,k}(y)\right)+\mathcal{D}_{k}\left(e_{k}yw_{i},p_{ji,k}(\omega)\right)\frac{\partial\lambda_{k}\left(.\right)}{\partial p_{ji,k}(\omega)}\right)dG_{i}\left(y\right)\right]L_{i}\\
\hspace{-1in}= & -\left[\int_{y}\frac{1}{p_{ji,k}(\omega)}\left(1+\theta_{k}\frac{e_{k}yw_{i}}{p_{ji,k}(\omega)}\frac{v_{k}'\left(e_{k}yw_{i}/p_{ji,k}(\omega)\right)}{v_{k}\left(e_{k}yw_{i}/p_{ji,k}(\omega)\right)}\right)\mathcal{D}_{k}\left(e_{k}yw_{i},p_{ji,k}(\omega)\right)\lambda_{k}\left(.\right)dG_{i}\left(y\right)\right]L_{i}\\
\hspace{-1in}= & -\frac{q_{k}\left(.\right)}{p_{ji,k}(\omega)}\left[1+\theta_{k}\int_{y}\frac{e_{k}yw_{i}}{p_{ji,k}(\omega)}\frac{v_{k}'\left(yw_{i}/p_{ji,k}(\omega)\right)}{v_{k}\left(yw_{i}/p_{ji,k}(\omega)\right)}d\mathcal{G}_{i,k}\left(y,w_{i}/p_{ji,k}(\omega),\varphi;\boldsymbol{P}_{i}\right)\right],
\end{align*}
where \emph{$\mathcal{G}_{i,k}$ }denotes the\emph{ demand-weighted}
income distribution, which augments the actual income distribution,
$G_{i}\left(.\right)$, by weighing each income group by their demand:
\[
d\mathcal{G}_{i,k}\left(y,w_{i}/p_{ji,k}(\omega),\varphi;\boldsymbol{P}_{i}\right)\equiv\frac{\lambda_{k}\left(y,e_{k}w_{i}/p_{ji,k}(\omega),\varphi,\boldsymbol{P}_{i}(y)\right)dG_{i}\left(y\right)}{\int_{y'}\lambda_{k}\left(y,e_{k}w_{i}/p_{ji,k}(\omega),\varphi,\boldsymbol{P}_{i}(y')\right)\cdot dG_{i}\left(y'\right)}\cdot
\]
Using the above expressions, we can immediately calculate the demand
elasticity, $\varepsilon_{k}\left(.\right)=\mid\partial\ln q_{k}\left(.\right)/\partial\ln p_{ji,k}(\omega)\mid$,
as follows:
\[
\varepsilon_{k}\left(e_{k}w_{i}/p_{ji,k}(\omega),\varphi;\boldsymbol{P}_{i},G_{i}\right)=1+\theta_{k}\left[\int_{y}\frac{v_{k}'\left(ye_{k}w_{i}/p_{ji,k}(\omega)\right)}{v_{k}\left(ye_{k}w_{i}/p_{ji,k}(\omega)\right)}\ y\ d\mathcal{G}_{i,k}\left(y;e_{k}w_{i}/p_{ji,k}(\omega),\varphi)\right)\right]\frac{e_{k}w_{i}}{p_{ji,k}(\omega)}.
\]
The above equation clearly states that apart from the aggregate demand
shifters, $\boldsymbol{P}_{i}$, the income distribution, $G_{i}\left(.\right)$,
and the quality level, $\varphi$, which applies to all firms in set
$\Omega_{i,k}\left(\varphi\right)$, $\varepsilon_{k}\left(.\right)$
only depends on the average demand level, $e_{k}w_{i}/p_{ji,k}(\omega)$.

To prove Claim (b), I first need to characterize the conditions under
which $\lambda_{k}\left(.\right)$ is invariant to income. To this
end, defining $x_{ji,k}(\omega)\equiv e_{k}y/p_{ji,k}(\omega)$, we
can take the derivative of $\lambda_{k}\left(.\right)$ as follows:
\begin{align}
\hspace{-0.75in}\frac{\partial\ln\lambda_{k}\left(x_{\omega},\varphi,\boldsymbol{P}_{i}\right)}{\partial\ln y}= & \theta_{k}\left(\frac{\partial\ln v_{k}\left(x_{ji,k}(\omega)\right)}{\partial\ln y}\ -\sum_{j\in\mathbb{C}}\sum_{\omega'\in\Omega_{j,k}\left(y,\varphi\right)}\frac{\partial\ln v_{k}\left(x_{ji,k}(\omega')\right)}{\partial\ln y}\frac{\lambda_{k}\left(x_{ji,k}(\omega'),\varphi,\boldsymbol{P}_{i}\right)}{\Lambda_{k}\left(\varphi,y\right)}\right)\label{eq: Lemma 2 -- part b}\\
\hspace{-0.75in}- & \gamma_{k}\left[\sum_{j\in\mathbb{C}}\sum_{\omega'\in\Omega_{j,k}\left(y,\varphi\right)}\frac{\partial\ln v_{k}\left(x_{ji,k}(\omega')\right)}{\partial\ln y}\frac{\lambda_{k}\left(x_{ji,k}(\omega'),\varphi,\boldsymbol{P}_{i}\right)}{\Lambda_{k}\left(\varphi,y\right)}\right.\nonumber \\
\hspace{-0.75in} & -\left.\sum_{\varphi'\in\varPhi_{k}}\sum_{j\in\mathbb{C}}\sum_{\omega'\in\Omega_{j,k}\left(y,\varphi\right)}x_{ji,k}(\omega')\frac{\partial\ln v_{k}\left(x_{ji,k}(\omega')\right)}{\partial\ln y}\lambda_{k}\left(x_{ji,k}(\omega'),\varphi',\boldsymbol{P}_{i}\right)\right],\nonumber 
\end{align}
where $\Lambda_{k}\left(\varphi,y\right)\equiv P_{i,k}(\varphi,y)^{-\gamma_{k}}/\sum_{\varphi'\in\varPhi_{k}}\left[P_{i,k}(\varphi',y)^{-\gamma_{k}}\right]$
denotes the share of consumers choosing quality $\varphi$. Given
that the probability shares add up to one, i.e.,
\[
\sum_{j\in\mathbb{C}}\sum_{\omega'\in\Omega_{j,k}\left(y,\varphi\right)}\lambda_{k}\left(x_{ji,k}(\omega'),\varphi,\boldsymbol{P}_{i}\right)/\Lambda_{k}\left(\varphi,y\right)=\sum_{\varphi'\in\varPhi_{k}}\sum_{j\in\mathbb{C}}\sum_{\omega'\in\Omega_{j,k}\left(y,\varphi\right)}\lambda_{k}\left(x_{ji,k}(\omega'),\varphi',\boldsymbol{P}_{i}\right)=1,
\]
a sufficient condition for $\partial\ln\lambda_{k}\left(.\right)/\partial y=0$
is that $\partial\ln v_{k}\left(x\right)/\partial\ln y=\textrm{constant}$.
Given that $x\equiv ey/p$, this condition can be otherwise stated
as
\[
xv_{k}'\left(x\right)/v_{k}\left(x\right)=c_{k},
\]
where $c_{k}$ is an industry-level constant. It follows trivially
from Equation \ref{eq: Lemma 2 -- part b} that the above condition
is also necessary for the invariance of $\lambda_{k}\left(.\right)$
w.r.t. $y$. In particular, if $\partial\ln\lambda_{k}\left(x_{ji,k}(\omega),\varphi,\boldsymbol{P}_{i}\right)/\partial\ln y=0$
is to be satisfied simultaneously for all $\omega$, then it should
be the case that $xv_{k}'\left(x\right)/v_{k}\left(x\right)=c_{k}$.
Plugging this condition into the general demand elasticity formula,
yields the constant elasticity expression:
\[
\varepsilon_{k}\left(e_{k}w_{i}/p_{ji,k}(\omega),\varphi;\boldsymbol{P}_{i},G_{i}\right)=1+\theta_{k}c_{k},
\]
where the formula stated under Lemma 2 concerns the special case,
where $c_{k}=1$.

\subsection{\label{sec: Proof-Lemma3}Demand Elasticity in Discrete Industries}

Claim (a) can be proven by applying the Leibniz integral rule as follows:
\begin{align*}
\frac{\partial q_{k}\left(p_{ji,k}(\omega),...\right)}{\partial p_{ji,k}(\omega)} & =-\lambda_{k}\left(\frac{p_{ji,k}(\omega)}{e_{k}w_{i}},p_{ji,k}(\omega),\tilde{p}_{i,k},\varphi,\boldsymbol{P}_{i}\left(y\right)\right)g_{i}\left(p_{ji,k}(\omega)/e_{k}w_{i}\right)L_{i}/e_{k}w_{i}+\\
 & -\frac{\theta_{k}}{\tilde{p}_{i,k}}\left(\int_{\frac{p_{ji,k}(\omega)}{e_{k}w_{i}}}^{\infty}\frac{v_{k}'\left(y\frac{e_{k}w_{i}}{\tilde{p}_{i,k}}-\frac{p_{ji,k}(\omega)}{\tilde{p}_{i,k}}\right)}{v_{k}\left(y\frac{e_{k}w_{i}}{\tilde{p}_{i,k}}-\frac{p_{ji,k}(\omega)}{\tilde{p}_{i,k}}\right)}\lambda_{k}\left(y,p_{ji,k}(\omega),\tilde{p}_{i,k},\varphi,\boldsymbol{P}_{i}\left(y\right)\right)dG_{i}(y)\right)L_{i}.
\end{align*}
where, as defined in the mains text, $\overline{\lambda}_{k}(p_{ji,k}(\omega),\tilde{p}_{i},\varphi,\boldsymbol{P}_{i}\left(y\right))\equiv q_{k}\left(.\right)/\left[1-G_{i}\left(p_{ji,k}(\omega)/e_{k}w_{i}\right)\right]L_{i}$
denotes the purchase probability of the average consumer who can afford
$\omega$. Recalling the definition of $\mathcal{G}_{i,k}\left(.\right)$
from Appendix \ref{sec:Proof-Lemma2}, the above equation immediately
implies the following expression for the demand elasticity, $\varepsilon_{k}\left(.\right)=\mbox{\ensuremath{\mid\partial\ln q_{k}\left(.\right)}/\ensuremath{\partial\ln p_{ji,k}}(\ensuremath{\omega})\ensuremath{\mid}}$:
\begin{align*}
\hspace{-0.8in}\varepsilon_{k}\left(\frac{p_{ji,k}(\omega)}{e_{k}w_{i}},\frac{p_{ji,k}(\omega)}{\tilde{p}_{i,k}},\varphi,\boldsymbol{P}_{i},G_{i}\right)= & \left[\frac{\lambda_{k}\left(\frac{p_{ji,k}(\omega)}{e_{k}w_{i}},\frac{p_{ji,k}(\omega)}{\tilde{p}_{i,k}},\varphi,\boldsymbol{P}_{i}\left(y\right)\right)}{\overline{\lambda}_{k}(p_{ji,k}(\omega),\varphi,\tilde{p}_{i,k},\boldsymbol{P}_{i}\left(y\right))}\frac{g_{i}(p_{ji,k}(\omega)/e_{k}w_{i})}{1-G_{i}\left(p_{ji,k}(\omega)/e_{k}w_{i}\right)}\right]\frac{p_{ji,k}(\omega)}{e_{k}w_{i}}\\
\hspace{-0.8in}+ & \theta_{k}\left[\int_{\frac{p_{ji,k}(\omega)}{e_{k}w_{i}}}^{\infty}\frac{v_{k}'\left(\tilde{e}_{i,k}y-\frac{p_{ji,k}(\omega)}{\tilde{p}_{i,k}}\right)}{v_{k}\left(\tilde{e}_{i,k}y-\frac{p_{ji,k}(\omega)}{\tilde{p}_{i,k}}\right)}d\mathcal{G}_{i,k}\left(y;\frac{p_{ji,k}(\omega)}{\tilde{p}_{i,k}},\varphi,\boldsymbol{P}_{i}\right)\right]\frac{p_{ji,k}(\omega)}{\tilde{p}_{i,k}},
\end{align*}
where $\tilde{e}_{k}\equiv e_{k}/\tilde{a}_{i,k}$. The above equation
clearly states that apart from the aggregate demand shifters, $\boldsymbol{P}_{i}$,
the income distribution, $G_{i}\left(.\right)$, and the quality level,
$\varphi$, which applies to all firms in set $\Omega_{j,k}\left(\varphi,y\right)$,
$\varepsilon_{k}\left(.\right)$ only depends on $(i)$ the degree
of affordability $p_{ji,k}(\omega)/e_{k}w_{i}$, and $(ii)$ the price
relative to services, $p_{ji,k}(\omega)/\tilde{p}_{i,k}$.

To prove Claim (b), I first need to characterize the conditions under
which $\lambda_{k}\left(.\right)$ is invariant to income. To this
end, we can take the same exact steps outlined earlier in the proof
of Lemma 2, which imply that individual demand probabilities are (conditionally)
invariant to $y$ iff $\partial\ln v_{k}\left(x\right)/\partial\ln y=\textrm{constant}$.
In the case of discrete demand industries, where $x=\left(y-p_{ji,k}(\omega)\right)/\tilde{p}_{i,k}$,
this condition simply reduces to
\[
v_{k}'\left(x\right)/v_{k}\left(x\right)=c_{k},
\]
where $c_{k}$ is an industry-level constant. Plugging this condition
into the general demand elasticity formula, yields the following formulation:
\[
\varepsilon_{k}\left(p_{\omega i}/e_{k}w_{i},p_{ji,k}(\omega)/\tilde{p}_{i,k},\varphi,\boldsymbol{P}_{i},G_{i}\right)=c_{k}\theta_{k}\frac{p_{ji,k}(\omega)}{\tilde{p}_{i,k}}\ +\ \rho_{ji,k}(\omega)\mathcal{H}_{i}(\frac{p_{ji,k}(\omega)}{e_{k}w_{i}})\frac{p_{ji,k}(\omega)}{e_{k}w_{i}},
\]
where $\mathcal{H}_{i}\left(y\right)=g_{i}\left(y\right)/\left(1-G_{i}\left(y\right)\right)$
the hazard rate of the income distribution in market $i$. Also, note
that the formula stated under Lemma 3 concerns the special case, where
$c_{k}=1$. The demand elasticity in benchmark model presented in
\ref{sec:Theoretical-Framework} can be attained as a special case
of the above formula, where $(i)$ $c_{k}=1$, and $(ii)$ the endowment
distribution is exponential, with a constant hazard rate $\mathcal{H}_{i}(.)=\zeta_{i}$.
So, given that $\rho_{ji,k}(\omega)=\rho_{ji,k}\left(p_{ji,k}(\omega)\right)\approx1$
to a first-order approximation, the above elasticity formula in this
special case reduces to Equation \ref{eq: Eq epsilon} from Section
\ref{sec:Theoretical-Framework}:
\[
\varepsilon_{ji,k}(\varphi)=\theta_{k}\frac{p_{ji,k}(\omega)}{\tilde{p}_{i,k}}\ +\ \zeta_{i}\frac{p_{ji,k}(\omega)}{e_{k}w_{i}}.
\]


\subsection{\label{sec: Deriving-Equation for QS}Deriving Equation \ref{eq: IQS}}

\paragraph*{Continuous industries ($k\in\mathbb{K}_{\mathcal{C}}$):}

In the case of continuous industries, the CES demand implies that
\[
\text{\ensuremath{\frac{Q_{ji,k}\left(\varphi_{H}\right)/Q_{ji,k}\left(\varphi_{L}\right)}{Q_{ji,k}\left(\varphi_{H}\right)/Q_{ji,k}\left(\varphi_{L}\right)}}}=1\Longleftrightarrow\text{\ensuremath{\frac{p_{ji,k}\left(\varphi_{H}\right)/p_{ji,k}\left(\varphi_{L}\right)}{p_{ji,k}\left(\varphi_{H}\right)/p_{ji,k}\left(\varphi_{L}\right)}}}=1,\ \ \ k\in\mathbb{K}_{\mathcal{C}}.
\]
The right-hand side of the above condition holds due to the assumption
price exhibits a constant markup and the country-specific cost-shifter
is quality-neutral: $p_{ji,k}=\mu_{k}\tau_{ji,k}a_{j,k}c(\varphi)w_{i}$.

\paragraph*{Discrete industries ($k\in\mathbb{K}_{\mathcal{D}}$):}

The demand facing a firm supplying discrete varieties from country
$j$ to market $i$ is given by:
\begin{equation}
q_{ji,k}\left(\varphi\right)=\left[\int_{\frac{p_{ji,k}\left(\varphi\right)}{e_{k}w_{i}}}^{\infty}\exp\left(-p_{ji,k}\left(\varphi\right)/\theta_{k}\right)\varPsi_{i,k}\left(\varphi,y\right)e^{-\zeta_{i}\left(y-\chi_{i}\right)}dy\right]\ \ \ \varphi\in\varPhi_{k},\label{eq: appendix demand homo.}
\end{equation}
In the above equation, $\varPsi_{i,k}\left(\varphi,y\right)$ is a
\emph{quality$\times$income}-specific demand shifters, which can
be expressed as follows:
\[
\varPsi_{i,k}\left(\varphi,y\right)=\frac{\left[\sum_{\ell\in\mathbb{C}}N_{\ell,k}(\varphi)\exp\left(-p_{ji,k}\left(\varphi\right)/\theta_{k}\right)\right]^{\frac{\gamma_{k}}{\theta_{k}}-1}}{\sum_{\varphi\in\varPhi_{k}}\left[\sum_{\ell\in\mathbb{C}}N_{\ell,k}(\varphi)\exp\left(-p_{ji,k}\left(\varphi\right)/\theta_{k}\right)\right]^{\frac{\gamma_{k}}{\theta_{k}}}}L_{i}.
\]
Since, $\varPsi_{i,k}\left(\varphi,y\right)$ is piecewise differentiable
with $\partial\varPsi_{i,k}\left(\varphi,y\right)/\partial y=0$,
we can perform integration by parts (pice-by-piece) to re-write Equation
\ref{eq: appendix demand homo.} as follows:
\begin{align*}
q_{ji,k}\left(\varphi\right) & =\exp\left(-p_{ji,k}\left(\varphi\right)/\tilde{p}_{i}\theta_{k}\right)\left[\frac{1}{\zeta_{i}}e^{-\zeta_{i}\left(\frac{p_{ji,k}\left(\varphi\right)}{e_{k}w_{i}}-\bar{y}_{i}\right)}\right]\varPsi_{i,k}\left(\varphi,p_{ji,k}\left(\varphi\right)/e_{k}w_{i}\right)\\
 & =A_{i}\ \exp\left(p_{ji,k}\left(\varphi\right)/w_{i}\right)^{-\vartheta_{i,k}}\varPsi_{i,k}\left(\varphi,p_{ji,k}\left(\varphi\right)/e_{k}w_{i}\right),
\end{align*}
where $1/\vartheta_{i,k}\equiv\theta_{k}/\tilde{a}_{i}+\zeta_{i}/e_{k}$
and $A_{i}\equiv e^{\zeta_{i}\bar{y}_{i}}/\zeta_{i}$. Noting that
$(i)$ $Q_{ji,k}\left(\varphi\right)=N_{j,k}\left(\varphi\right)q_{ji,k}\left(\varphi\right)$
and $(ii)$ $p_{ji,k}\left(\varphi\right)=\tau_{ji,k}a_{j,k}w_{j}c_{k}\left(\varphi\right)+\vartheta_{i,k}w_{i}$
, the above equation yields the following expression:
\begin{align}
\frac{Q_{ji,k}\left(\varphi_{H}\right)}{Q_{ji,k}\left(\varphi_{L}\right)} & =\frac{N_{j,k}\left(\varphi_{H}\right)}{N_{j,k}\left(\varphi_{L}\right)}\exp\left(\left[p_{ji,k}\left(\varphi_{L}\right)-p_{ji,k}\left(\varphi_{H}\right)\right]\right)^{\frac{\vartheta_{i,k}}{w_{i}}}\frac{\varPsi_{i,k}\left(\varphi_{H},p_{ji,k}\left(\varphi_{H}\right)/e_{k}w_{i}\right)}{\varPsi_{i,k}\left(\varphi_{L},p_{ji,k}\left(\varphi_{L}\right)/e_{k}w_{i}\right)}\nonumber \\
 & =\frac{N_{j,k}\left(\varphi_{H}\right)}{N_{j,k}\left(\varphi_{L}\right)}\exp\left(\left[c_{k}\left(\varphi_{L}\right)-c_{k}\left(\varphi_{H}\right)\right]\tau_{ji,k}a_{j,k}w_{j}\right)^{\frac{\vartheta_{i,k}}{w_{i}}}\frac{\varPsi_{i,k}\left(\varphi_{H},p_{ji,k}\left(\varphi_{H}\right)/e_{k}w_{i}\right)}{\varPsi_{i,k}\left(\varphi_{L},p_{ji,k}\left(\varphi_{L}\right)/e_{k}w_{i}\right)}.\label{eq: appendix demand homo 2}
\end{align}
Finally, given that to a first-order approximation $\varPsi_{i,k}\left(\varphi,y\right)=\varPsi_{i,k}\left(\varphi,\bar{y}_{i}\right)$
for all $y$, then for any two suppliers $j$ and $\ell$,
\[
\frac{\varPsi_{i,k}\left(\varphi_{H},p_{ji,k}\left(\varphi_{H}\right)/e_{k}w_{i}\right)/\varPsi_{i,k}\left(\varphi_{L},p_{ji,k}\left(\varphi_{L}\right)/e_{k}w_{i}\right)}{\varPsi_{i,k}\left(\varphi_{H},p_{\ell i,k}\left(\varphi_{H}\right)/e_{k}w_{i}\right)/\varPsi_{i,k}\left(\varphi_{L},p_{\ell i,k}\left(\varphi_{L}\right)/e_{k}w_{i}\right)}\approx1.
\]
That being the case, Equation \ref{eq: appendix demand homo 2} immediately
implies Expression \ref{eq: IQS} for $k\in\mathbb{K}_{\mathcal{D}}$.

\subsection{\label{subsec: Proof Proposition 2}Proof of Proposition 2}

Let $r_{ji,k}(\varphi)=Q_{ji,k}(\varphi)/\sum_{\varphi'\in\varPhi_{k}}Q_{ji,k}(\varphi')$
and let us only focus on a discrete industry $k\in\mathbb{K}_{\mathcal{D}}$.
An immediate corollary of Equation \ref{eq: IQS} (derived in the
previous appendix) is that for any $\varphi_{H}$ and $\varphi_{L}\in\varPhi_{k}$
such that $\varphi_{H}>\varphi_{L}$:
\[
\frac{\partial r_{ji,k}(\varphi_{H})/r_{ji,k}(\varphi_{L})}{\partial\tilde{p}_{i}},\ \frac{\partial r_{ji,k}(\varphi_{H})/r_{ji,k}(\varphi_{L})}{\partial\tilde{p}_{j}},\ \frac{\partial r_{ji,k}(\varphi_{H})/r_{ji,k}(\varphi_{L})}{\partial a_{j,k}}>0.
\]
Combining the above relationship with the fact that $\sum_{\varphi\in\varPhi_{k}}r_{ji,k}(\varphi)=1$,
the above expression simply implies that
\[
\frac{\partial^{2}r_{ji,k}(\varphi)}{\partial\varphi\partial\tilde{p}_{i}},\ \frac{\partial^{2}r_{ji,k}(\varphi)}{\partial\varphi\partial\tilde{p}_{j}},\ \frac{\partial^{2}r_{ji,k}(\varphi)}{\partial\varphi\partial a_{j,k}}>0.
\]
Now let us focus on the effect of $\tilde{p}_{i}$ on the aggregate
export unit price, $\bar{p}_{ji,k}$. Given that $(a)$ $\frac{\partial p_{ji,k}(\varphi)}{\partial\tilde{p}_{i}}$,
$(b)$ $\frac{\partial p_{ji,k}(\varphi)}{\partial\varphi}>0$, and
$(c)$ $\frac{\partial^{2}r_{ji,k}(\varphi)}{\partial\varphi\partial\tilde{p}_{i}}>0$,
it immediately follows that
\[
\frac{\partial\bar{p}_{ji,k}}{\partial\tilde{p}_{i}}=\sum_{\varphi'\in\varPhi_{k}}\left(p_{ji,k}(\varphi)\frac{\partial r_{ji,k}(\varphi)}{\partial\tilde{p}_{i}}+\frac{\partial p_{ji,k}(\varphi)}{\partial\tilde{p}_{i}}r_{ji,k}(\varphi\right))>0,\ \ \ k\in\mathbb{K}_{\mathcal{D}}.
\]
Along the same line arguments we can establish that $\frac{\partial\bar{p}_{ji,k}}{\partial\tilde{p}_{j}}>0$
and $\frac{\partial\bar{\varphi}_{ji,k}}{\partial a_{j,k}}>0$ if
$k\in\mathbb{K}_{\mathcal{D}}$. Meanwhile, the fact $\frac{\partial\bar{p}_{ji,k}}{\partial\tilde{p}_{i}}=\frac{\partial\bar{p}_{ji,k}}{\partial\tilde{p}_{j}}=0$
if $k\in\mathbb{K}_{\mathcal{C}}$ follows trivially from all components
of $r_{ji,k}(\varphi)$ and $p_{ji,k}(\varphi)$ being invariant to$\tilde{p}_{j}$,
$\tilde{p}_{i}$. Relatedly, $\frac{\partial\bar{\varphi}_{ji,k}}{\partial a_{j,k}}=0$
follows from $r_{ji,k}(\varphi)$ being invariant to $a_{j,k}$.

\subsection{\label{subsec: GT Derivation}Gains from Trade Formulas}

This appendix derives the gains from trade formula presented in Section
\ref{sec: GT} under Equation \ref{eq: GT one-sector}. Recall that
this equation is derived under the assumption that $(a)$ the traded
sector is entirely discrete (labeled $\mathcal{D}$), and $(b)$ consumers
spend $e_{i,\mathcal{D}}$ of their total budget on consumption categories
where traded discrete goods are combined with non-traded services.
Based on these assumptions, Equation \ref{eq: GT} implies that
\[
GT_{i}=1-\left(\lambda_{ii,\mathcal{D}}/\hat{N}_{i,\mathcal{D}}\right)^{\frac{e_{i,\mathcal{D}}}{\theta}}.
\]
To arrive at Equation \ref{eq: GT one-sector} from the above expression,
we need to characterize $\hat{N}_{i,\mathcal{D}}$. To this end, we
can appeal to the free-entry condition:
\[
N_{i,\mathcal{D}}w_{i}f^{e}=\sum_{\ell\in\mathbb{C}}(p_{i\ell,\mathcal{D}}-c_{i\ell,\mathcal{D}})Q_{i\ell,\mathcal{D}}.
\]
where $Q_{i\ell,\mathcal{D}}=\lambda_{i\ell,\mathcal{D}}L_{\ell}$.
In the absence of affordability constraints, $p_{i\ell,\mathcal{D}}=c_{i\ell,\mathcal{D}}+w_{\ell}/\tilde{a}_{\ell}\theta$.
Assuming $\tilde{a}_{\ell}=1$ for all $\ell$, and substituting for
$p_{i\ell,\mathcal{D}}-c_{i\ell,\mathcal{D}}$ in the above expression
yields
\[
N_{i,\mathcal{D}}w_{i}f^{e}=\sum_{\ell\in\mathbb{C}}\frac{w_{\ell}}{\theta}\lambda_{i\ell,\mathcal{D}}L_{\ell}=\sum_{\ell\in\mathbb{C}}\left(\frac{1}{\theta}\lambda_{i\ell,\mathcal{D}}\frac{w_{\ell}L_{\ell}}{w_{i}L_{i}}\right)w_{i}L_{i}.
\]
Rearranging the above equation implies, $N_{i,\mathcal{D}}=\sum_{\ell\in\mathbb{C}}\left(\lambda_{i\ell,\mathcal{D}}\frac{w_{\ell}L_{\ell}}{w_{i}L_{i}}\right)\frac{L_{i}}{\theta f^{e}}$.
Accordingly, under autarky, $N_{i,\mathcal{D}}^{A}=\frac{L_{i}}{\theta f^{e}}$,
which yields 
\[
\hat{N}_{i,\mathcal{D}}=N_{i,\mathcal{D}}^{A}/N_{i,\mathcal{D}}=\sum_{\ell\in\mathbb{C}}\left(\lambda_{i\ell,\mathcal{D}}\frac{w_{\ell}L_{\ell}}{w_{i}L_{i}}\right).
\]
Plugging the above expression back in to expression for $GT_{i}$,
yields \ref{eq: GT one-sector}.

\subsection{\label{subsec: Equation 17}Derivation of Equation \ref{eq: N_i,D}}

As shown in the previous appendix, we can write the (Free Entry) equation
for country $1$ as follows:
\begin{align}
N_{1,\mathcal{D}} & =\left(\lambda_{12,\mathcal{D}}\frac{w_{2}L_{2}}{w_{1}L_{1}}+\lambda_{11,\mathcal{D}}\frac{w_{1}L_{1}}{w_{1}L_{1}}\right)\frac{L_{1}}{\theta f^{e}}=\left(\lambda_{12,\mathcal{D}}\frac{w_{2}L_{2}}{w_{1}L_{1}}+\lambda_{11,\mathcal{D}}\right)\frac{L_{1}}{\theta f^{e}}\label{eq: N1_D}
\end{align}
Rearranging the (Balanced Trade) equation, the term $\lambda_{12,\mathcal{D}}\frac{w_{2}L_{2}}{w_{1}L_{1}}$
(in Equation \ref{eq: N1_D}) can be alternatively expressed as
\begin{align*}
\lambda_{12,\mathcal{D}}\frac{w_{2}L_{2}}{w_{1}L_{1}} & =\left(\frac{\tau a\frac{w_{2}}{w_{1}}+\frac{1}{\theta}}{\tau a\frac{w_{1}}{w_{2}}+\frac{1}{\theta}}\right)\lambda_{21,\mathcal{D}}\\
 & =\lambda_{21,\mathcal{D}}\left[1+\frac{\tau a\frac{w_{2}}{w_{1}}-\tau a\frac{w_{1}}{w_{2}}}{\tau a\frac{w_{1}}{w_{2}}+\frac{1}{\theta}}\right]=\lambda_{21,\mathcal{D}}\left[1+\frac{\frac{w_{2}}{w_{1}}-\frac{w_{1}}{w_{2}}}{\frac{w_{1}}{w_{2}}+\frac{1}{\tau a\theta}}\right].
\end{align*}
Plugging the above expression for $\lambda_{12,\mathcal{D}}\frac{w_{2}L_{2}}{w_{1}L_{1}}$
back into Equation \ref{eq: N1_D} yields the following expression
for $N_{1,\mathcal{D}}$:
\begin{align}
N_{1,\mathcal{D}} & =\left(\lambda_{21,\mathcal{D}}\left[1+\frac{\frac{w_{2}}{w_{1}}-\frac{w_{1}}{w_{2}}}{\frac{w_{1}}{w_{2}}+\frac{1}{\tau a\theta}}\right]+\lambda_{11,\mathcal{D}}\right)\frac{L_{1}}{\theta f^{e}}=\left(1+\frac{\frac{w_{2}}{w_{1}}-\frac{w_{1}}{w_{2}}}{\frac{w_{1}}{w_{2}}+\frac{1}{\tau a\theta}}\lambda_{21,\mathcal{D}}\right)\frac{L_{1}}{\theta f^{e}},\label{eq: N1_D NEW}
\end{align}
where the second line follows from $\lambda_{21,\mathcal{D}}+\lambda_{11,\mathcal{D}}=1$.
Following the same exact steps, we can derive the following expression
for $N_{2,D}$:
\begin{equation}
N_{2,\mathcal{D}}=\left(1+\frac{\frac{w_{1}}{w_{2}}-\frac{w_{2}}{w_{1}}}{\frac{w_{2}}{w_{1}}+\frac{1}{\tau a\theta}}\lambda_{21,\mathcal{D}}\right)\frac{L_{2}}{\theta f^{e}}.\ \ \ \ \ \ \ \ \label{eq: N2_D NEW}
\end{equation}
Equations \ref{eq: N1_D NEW} and \ref{eq: N2_D NEW} can be alternatively
formulated as
\[
N_{i,\mathcal{D}}=\left(1+\frac{\frac{w_{-i}}{w_{i}}-\frac{w_{i}}{w_{-i}}}{\frac{w_{i}}{w_{-i}}+\frac{1}{a\tau\theta}}\lambda_{-ii,\mathcal{D}}\right)\frac{L_{i}}{\theta f^{e}}.
\]


\section{\label{sec: Discrete-Trade-with-Firm-Het}Discrete Trade with Firm-Level
Heterogeneity}

Below, I outline the implications of firm-heterogeneity for the model.
Before doing so, note that the firm-level predictions of the model,
with regards to PTM, are not sensitive to firm-level heterogeneity
by construction. The macro-level predictions, regarding IQS, however,
can change with the introduction of firm heterogeneity. 

In the baseline model, $c_{k}(\varphi,z$) denotes the marginal labor
requirement for producing quality $\varphi$ in industry $k$. In
the baseline case, $z$ differs across countries but is the same across
all firms in a given country. So, to model firm heterogeneity I make
the following amendments to the baseline model:
\begin{enumerate}
\item Firms are ex-ante identical in industry $k$. After paying the entry
cost, firms are assigned a productivity, $z$, which is the realization
of a random variable drawn independently across firms from a Pareto
distribution:
\[
F_{j,k}\left(z\right)=1-A_{j,k}z^{-\eta_{k}},
\]
where $A_{j,k}$ is a country$\times$industry-specific productivity
shifter.
\item Firms have to incur a fixed cost, $f_{i,k}\left(\varphi\right)$,
to serve market $i$ in quality level $\varphi\in\varPhi_{k}$.
\end{enumerate}
Extending the discussion in Appendix \ref{sec: Non-Parametric-Model},
the profits facing each firm with productivity $z$ can be expressed
as $\pi_{k}\left(\nu_{ji,k}\left(\varphi,z\right),\varphi,\boldsymbol{P}_{i,k}\right)$,
where $\nu_{ji,k}\left(\varphi,z\right)\equiv p_{ji,k}\left(\varphi,z\right)/w_{i}$.
The zero-profit cutoff is, therefore, determined by the following
equation:
\[
\pi_{k}\left(\nu_{ji,k}^{*}\left(\varphi,z\right),\varphi,\boldsymbol{P}_{i,k}\right)=w_{i}f_{i,k}\left(\varphi\right).
\]
Importantly, the above equation implies that $\nu_{ji,k}^{*}\left(\varphi,z\right)$
should be uniform across all exporting countries. Namely,
\[
\nu_{ji,k}\left(\varphi,z_{ji,k}^{*}\left(\varphi\right)\right)\equiv\nu_{i,k}^{*}\left(\varphi\right).
\]
Correspondingly, the total value of exports from country $j$ to market
$i$ in quality-level $\varphi\in\varPhi_{k}$ is given by
\begin{equation}
X_{ji,k}\left(\varphi\right)=N_{j,k}\left(\varphi\right)\left(\int_{z_{ji}^{*}\left(\varphi\right)}^{\infty}\nu_{ji,k}\left(\varphi,z\right)\lambda_{k}\left(\nu_{ji,k}\left(\varphi,z\right),\varphi,\boldsymbol{P}_{i,k}\right)w_{i}L_{i}\ \text{d}F_{j,k}\left(z\right)\right)w_{i}L_{i}.\label{eq: Xji Firm-Het}
\end{equation}
To present the next step transparently, let $m_{k}\left(\nu_{ji,k}\left(\varphi,z\right),\varphi,\boldsymbol{P}_{i,k}\right)\equiv\varepsilon/\varepsilon-1$
denote the markup charged by a firm with characteristics $\left(\varphi,z\right)$,
exporting industry $k$ goods from country $j$ to market $i$. Then,
by definition, $m_{k}\left(\nu_{ji,k}\left(\varphi,z\right),\varphi,\boldsymbol{P}_{i,k}\right)w_{j}\tau_{ji}c_{k}\left(\varphi\right)/w_{i}z=\nu_{ji,k}\left(\varphi,z\right)$;
and we can write $\text{d}\ln z$ in terms of $\text{d}\ln\nu$ as
follows
\[
\text{d}\ln z=\left(1+\frac{\partial\ln m_{k}}{\partial\ln\nu}\right)\text{d}\ln\nu.
\]
The above expression, in turn, implies that
\[
\text{d}z=\frac{w_{j}\tau_{ji}}{w_{i}}\frac{c_{k}\left(\varphi\right)m_{k}\left(\nu,\varphi,\boldsymbol{P}_{i,k}\right)}{\nu\left(1+\frac{\partial\ln m_{k}}{\partial\ln\nu}\right)^{-1}}\frac{\text{d}\nu}{\nu}.
\]
Given the Pareto productivity distribution, the above equation yields
the following formulation for $dF_{j,k}\left(z\right)$:
\begin{align*}
dF_{j,k}\left(z\right)= & \eta_{k}A_{j,k}z^{-\eta_{k}-1}\text{d}z\\
= & \eta_{k}A_{j,k}\left(\frac{w_{j}\tau_{ji}}{w_{i}}\frac{c_{k}\left(\varphi\right)m_{k}\left(\nu,\varphi,\boldsymbol{P}_{i,k}\right)}{\nu\left(1+\frac{\partial\ln m_{k}}{\partial\ln\nu}\right)^{-1}}\right)^{-\eta_{k}}\frac{\text{d}\nu}{\nu}.
\end{align*}
Plugging the above equation back into Equation \ref{eq: Xji Firm-Het},
yields the following expression for $X_{ji}\left(\varphi\right)$:
\[
X_{ji}\left(\varphi\right)=\psi_{k}A_{j,k}N_{j}\left(\varphi\right)\left(\frac{\tau_{ji,k}w_{j}}{w_{i}}\right)^{-\eta_{k}}w_{i}L_{i}\mathcal{Q}_{i,k}\left(\varphi\right),
\]
where
\[
\mathcal{Q}_{i,k}\left(\varphi\right)\equiv\int_{0}^{\nu_{i,k}^{*}\left(\varphi\right)}\lambda_{k}\left(\nu,\varphi,\boldsymbol{P}_{i,k}\right)\left[\frac{c_{k}\left(\varphi\right)m_{k}\left(\nu,\varphi,\boldsymbol{P}_{i,k}\right)}{\nu\left(1+\frac{\partial\ln m_{k}}{\partial\ln\nu}\right)^{-1}}\right]^{-\eta_{k}}\text{d}\ensuremath{\nu},
\]
depends only on the characteristics of the importing market, $i$.
Letting $E_{i}\left(\varphi\right)$ denote the total expenditure
on quality $\varphi\in\varPhi_{k}$, the budget constraint of consumers
in market $i$ entails that $\sum_{j\in\mathbb{C}}X_{ji}\left(\varphi\right)=E_{i}\left(\varphi\right)$.
This condition plus the above equation, imply the following gravity
equation describing macro-level export values:
\[
X_{ji}\left(\varphi\right)=\frac{A_{j,k}N_{j}\left(\varphi\right)\left(\tau_{ji,k}w_{j}\right)^{-\eta_{k}}}{\sum_{\ell\in\mathbb{C}}A_{\ell,k}N_{\ell}\left(\varphi\right)\left(\tau_{\ell i,k}w_{\ell}\right)^{-\eta_{k}}}E_{i}\left(\varphi\right).
\]
Considering the above expression, the prediction of model regarding
the \emph{home-market driven} quality specialization remains intact.
That is, high-wage countries, which have a greater home-market demand
for quality, would host a greater number of high-quality producing
firms. Accordingly, due to their higher $N_{j}\left(\varphi\right)$
in high-quality (high-$\varphi$) categories, they would export relatively
more high-quality goods. However, based on the above equation, the
prediction that affordability constraints induce IQS no longer holds.
One should nonetheless keep in mind that the break-down of affordability-driven
IQS is a mere artifact of the Pareto assumption, rather than firm
heterogeneity per se.

\section{\label{sec:Additional-Test-the}Additional Test the Model}

Section \ref{sec:Empirical-Evidence} tested Proposition 2 while restricting
the \noun{Baci} dataset to sub-samples of \emph{$(i)$ }only\emph{
final goods} and $(ii)$ only \emph{differentiated} \emph{goods} .
In this appendix I present an alternative test that groups HS6 product
by their underlying demand elasticity. To this end, I use the estimated
elasticities by \citet{Broda2006}. I classify HS6 products as less-
and more- differentiated based on the following criteria:
\begin{enumerate}
\item \emph{More-differentiated} HS6 products exhibit a lower-or-equal-to-median
elasticity, i.e., $\sigma_{BW}\leq3.6$.
\item \emph{Less-differentiated} HS6 products exhibit an above-median demand
elasticity, i.e., $\sigma_{BW}>3.6$. 
\end{enumerate}
I then re-estimate Equation \ref{eq: Baci Regression} separately
on these two sub-groups of products. The estimation results are reported
in Table \ref{tab:price-regression-robust_BW}, and corroborate the
model's prediction that export price levels --in discrete product
categories-- vary systematically with the price of non-traded services
in the origin and destination economy as well the origin economy's
degree of comparative advantage. At the same time, these same patterns
are either non-existent or weaker in continuous product categories---refer
to Section \ref{sec:Empirical-Evidence} for a more elaborate discussion
of these predictions.

\begin{table}[!htb] \def\sym#1{\ifmmode^{#1}\else\(^{#1}\)\fi}
\begin{centering}
\begin{adjustwidth}{-0.5in}{-0in}
\caption{Grouping products based on  Broda-Weinstein elasticities  (\emph{Dependent}: $\ln \bar{p}_{ji,k}$)}\label{tab:price-regression-robust_BW} 
\vspace{-0.1in}
\resizebox{1.2\textwidth}{!}{%  %
\begin{tabular}{l*{8}{D{.}{.}{-1}}}

\toprule                    
& \multicolumn{3}{c}{\Large{$\sigma_{BW}\leq$3.6}} & \phantom{abc}  & \multicolumn{3}{c}{\Large{$\sigma_{BW}$>3.6}} \\
\cmidrule{2-4} \cmidrule{6-8}

Regressor (log)  &\multicolumn{1}{c}{(1)} &\multicolumn{1}{c}{(2)} &\multicolumn{1}{c}{(3)} && &\multicolumn{1}{c}{(1)} &\multicolumn{1}{c}{(2)} &\multicolumn{1}{c}{(3)} \\
\toprule   

{\large $\tilde{p}_i$}           &                     &      -0.051         &                     &&                     &       0.021         &                     \\                     &                     &     (0.045)         &                     &&                     &     (0.025)         &                     \\ \addlinespace 
{\large $\tilde{p}_i$} $\times$ Discrete       &                     &       0.129\sym{***}&       0.134\sym{***}&&                     &       0.165\sym{***}&       0.158\sym{***}\\                     &                     &     (0.019)         &     (0.017)         &&                     &     (0.026)         &     (0.016)         \\


{\large $\tilde{p}_j$}             &       0.158\sym{***}&                     &&                     &       0.054\sym{***}&                     &                     \\                     &     (0.002)         &                     &&                     &     (0.001)         &                     &                     \\ \addlinespace 
{\large $\tilde{p}_j$} $\times$ Discrete        &       0.095\sym{***}&                     &       0.186\sym{***}&&       0.022\sym{***}&                     &       0.047\sym{**} \\                     &     (0.003)         &                     &     (0.022)         &&     (0.002)         &                     &     (0.021)         \\ \addlinespace


{\large $\textrm{RCA}_{j,k}$}             &       0.026\sym{***}&                     &       0.032\sym{***}&&       0.021\sym{***}&                     &       0.031\sym{***}\\                     &     (0.000)         &                     &     (0.005)         &     (0.000)         &&                     &     (0.004)         \\ \addlinespace {\large $\textrm{RCA}_{j,k}$} $\times$ Discrete        &       0.035\sym{***}&                     &       0.036\sym{***}&&       0.061\sym{***}&                     &       0.082\sym{***}\\                     &     (0.001)         &                     &     (0.011)         &&     (0.001)         &                     &     (0.010)         \\ \addlinespace 

Controls for GDP \& GDP p/c         & \multicolumn{1}{c}{Y} & \multicolumn{1}{c}{Y} & \multicolumn{1}{c}{...} &&   \multicolumn{1}{c}{Y} & \multicolumn{1}{c}{Y} & \multicolumn{1}{c}{...} \\
Controls for Distance        & \multicolumn{1}{c}{Y} & \multicolumn{1}{c}{Y} & \multicolumn{1}{c}{...} &&   \multicolumn{1}{c}{Y} & \multicolumn{1}{c}{Y} & \multicolumn{1}{c}{...} \\
HS6 FE       & \multicolumn{1}{c}{...} & \multicolumn{1}{c}{...} & \multicolumn{1}{c}{Y} && \multicolumn{1}{c}{...} & \multicolumn{1}{c}{...} & \multicolumn{1}{c}{Y} \\
Exporter$\times$HS6 FE       & \multicolumn{1}{c}{N} & \multicolumn{1}{c}{Y} & \multicolumn{1}{c}{N} && \multicolumn{1}{c}{N} & \multicolumn{1}{c}{Y} & \multicolumn{1}{c}{N} \\
Importer$\times$HS6 FE       & \multicolumn{1}{c}{Y} & \multicolumn{1}{c}{N} & \multicolumn{1}{c}{N} && \multicolumn{1}{c}{Y} & \multicolumn{1}{c}{N} & \multicolumn{1}{c}{N} \\
Exporter$\times$Importer FE & \multicolumn{1}{c}{...} & \multicolumn{1}{c}{...} & \multicolumn{1}{c}{Y} && \multicolumn{1}{c}{...} & \multicolumn{1}{c}{...} & \multicolumn{1}{c}{Y} \\
Observations                &    \multicolumn{3}{c}{2,204,000}        &&     \multicolumn{3}{c}{ 2,492,000}       \\
\bottomrule \addlinespace 
\end{tabular} }
\noindent\begin{minipage}[t]{1.2\columnwidth} \begin{spacing}{0.8} {\footnotesize \emph{Note:} This table corresponds to estimating Equation \ref{eq: Baci Regression} on more- and less differentiated goods as implied by Broda-Weinstein elasticity estimates, $\sigma_{BW}$. Robust standard errors are reported in parentheses. \sym{*} indicates statistically significant at 10\% level, \sym{**} indicates statistically significant at 5\% level, and  \sym{***} indicates statistically significant at 1\% level. Unit prices are constructed using the BACI database in 2011; the price of non-traded services are constructed using the ICP database in the same year; and each observation in the sample is weighted by trade quantity.}
\end{spacing} \end{minipage}
\end{adjustwidth}
\end{centering}
\end{table}

\section{\label{sec:Estimation-of-theta}Estimation of $\theta$}

This appendix describes how the triple difference procedure developed
by \citet{caliendo2014estimates} can be employed to estimate $\theta$.
This procedure relies on applied tariff data. So, before outlining
the procedure, I describe how the data on applied tariffs is compiled. 

\paragraph*{The UNCTAD-TRAINS Data}

The United Nations Statistical Division, Trade Analysis and Information
System (UNCTAD-TRAINS) reports data on applied tariff for 2008, spanning
31 two-digit (in ISIC rev.3) sectors, 181 importers, and 245 export
partners. In line with \citet{caliendo2014estimates}, I use the \emph{simple
tariff line average} of the \emph{effectively applied tariff} (AHS)
denoting it by $t_{ji,k}$. In instances where tariff data are missing
for 2008, I use tariff data for the nearest available year, giving
priority to earlier years. To aggregate the UNCTAD-TRAINS data into
individual WIOD industries, I closely follow the methodology outlined
in \citet{kucheryavyy2016grounded}. Importantly, individual European
Union (EU) member countries are not represented in the UNCTAD-TRAINS
data during the 2000-2014 period. So, I infer their applied tariff
rates from applied EU tariffs, given that intra-EU trade is subject
to zero tariffs while all EU members impose a common external tariff
on non-members.

\paragraph*{Continuous Model.}

Equation \ref{eq: expenditure share}, bilateral trade shares in continuous
traded industries can be characterized as
\[
\lambda_{ji,k}=N_{j,k}\Omega_{i,k}\left[c_{ji,k}(1+t_{ji,k})\right]{}^{-\theta}\ \ \ k\in\mathbb{K}_{\mathcal{C}},
\]
where $\Omega_{i,k}\equiv\sum_{n}\left(N_{n,k}\left[c_{ni,k}(1+t_{ni,k})\right]{}^{-\theta}\right)e_{i,k}$
can be treated as an importer$\times$industry fixed effect, and $c_{ji,k}$
denotes the marginal cost of producing and transporting industry $k$
goods from origin $j$ to destination $i$. Suppose $\ln c_{ji,k}=\ln\nu_{j,k}+\ln d_{ji,k}+\varepsilon_{ji,k}$,
where $d_{ji,k}=d_{ij,k}$ is a systematic and symmetric component
that accounts for the cost-equivalent effect of distance, common language,
and common border, and $\varepsilon_{ji,k}$ is a random disturbance
term that represents deviation from symmetry. Using this decomposition,
we can produce the following estimating equation for any triplet $(j,i,n)$:
\[
\ln\frac{\lambda_{ji,k}\lambda_{in,k}\lambda_{nj,k}}{\lambda_{ij,k}\lambda_{ni,k}\lambda_{jn,k}}=-\theta\ln\frac{(1+t_{ji,k})(1+t_{in,k})(1+t_{nj})}{(1+t_{ij,k})(1+t_{ni,k})(1+t_{jn,k})}+\tilde{\varepsilon}_{jin,k},
\]
where $\tilde{\varepsilon}_{jin,k}\equiv\theta_{k}(\varepsilon_{ij,k}-\varepsilon_{ji,k}+\varepsilon_{in,k}-\varepsilon_{ni,k}+\varepsilon_{nj,k}-\varepsilon_{jn,k})$.
The above equation can be used to attain unbiased and consistent estimates
for $\theta$ under the identifying assumption that $\textrm{Cov}(t_{ji,k},\varepsilon_{ji,k})=0$.
 I estimate the above equation using data on trade shares, $\lambda_{ji,k}$,
from the WIOD and matching data on applied tariffs, $t_{ji,k}$, from
the UNCTAD-TRAINS. As noted earlier, the original WIOD data features
16 traded industries. Since I am interested in attaining an economy-wide
estimate for $\theta$, I perform the estimation while pooling data
across all 16 industries, treating them as continuous.

\paragraph*{Discrete Model.}

Equation \ref{eq: expenditure share}, trade shares in discrete traded
industries can be characterized as:
\[
\lambda_{ji,k}=N_{j,k}\Omega_{i,k}\exp(c_{ji,k}+\bar{t}_{ji,k})^{-\theta/\tilde{p}_{j}},\ \ \ k\in\mathbb{K}_{\mathcal{D}}.
\]
where $\Omega_{i,k}\equiv\sum_{n}\left(N_{n,k}\exp(c_{ni,k}+\bar{t}_{ni,k})^{-\theta/\tilde{p}_{i}}\right)e_{i,k}$
can be once again treated as an importer$\times$industry fixed effect,
and $\bar{t}_{ji,k}$ is the per-unit equivalent of the ad-valorem
tariff rate, $t_{ji,k}$. Suppose $c_{ij,k}/\tilde{p}_{j}=\nu_{j,k}+\delta_{i,k}+d_{ij,k}+\varepsilon_{ij,k}$,
where $d_{ji,k}=d_{ij,k}$ is a systematic and symmetric component
that accounts for the cost-equivalent of distance, common language,
and common border, and $\varepsilon_{ji,k}$ is a random disturbance
term that represents deviation from symmetry. This decomposition produces
the following estimating equation for any triplet $(j,i,n)$:

\[
\ln\frac{\lambda_{ji,k}\lambda_{in,k}\lambda_{nj,k}}{\lambda_{ij,k}\lambda_{ni,k}\lambda_{jn,k}}=-\theta\left[\frac{\bar{t}_{ji,k}}{\tilde{p}_{i}}+\frac{\bar{t}_{in,k}}{\tilde{p}_{n}}+\frac{\bar{t}_{nj,k}}{\tilde{p}_{j}}-\frac{\bar{t}_{ij,k}}{\tilde{p}_{j}}-\frac{\bar{t}_{ni,k}}{\tilde{p}_{i}}-\frac{\bar{t}_{jn,k}}{\tilde{p}_{n}}\right]+\tilde{\varepsilon}_{jin,k},
\]
where $\tilde{\varepsilon}_{jin,k}\equiv\theta_{k}(\varepsilon_{ij,k}-\varepsilon_{ji,k}+\varepsilon_{in,k}-\varepsilon_{ni,k}+\varepsilon_{nj,k}-\varepsilon_{jn,k})$.
The above equation can be used to attain unbiased and consistent estimates
for $\theta$ under the identifying assumption that $\textrm{cov}(t_{ji,k},\varepsilon_{ji,k})=0$.
I estimate the above equation using data on trade shares and applied
tariffs from the the WIOD and UNCTAD-TRAINS. Since I am interested
in attaining an economy-wide estimate for $\theta$, I perform the
estimation while pooling data across the 16 traded WIOD industries,
treating them all as discrete.

\section{\label{sec: Construction-of-Figure}Out-of-Sample Performance of
the Discrete Model}

This appendix investigates the out-of-sample predictive power of the
discrete model that is calibrated to data in Section \ref{sec: GT}.
The data moment of interest is the share of expenditure on traded
versus non-trade sectors. Recall that the model is calibrated to fit
only the total share of expenditure on categories where traded and
non-traded services are bundled in the consumption basket.

In line with my analysis in Section \ref{sec: GT}, suppose the traded
sector consists of only discrete industries. In that case, the $L_{i}$
consumers in economy $i$ purchase in total $\mathcal{Q}_{i}=K\times L_{i}$
units of traded goods across all $K$ industries. Hence, letting $\bar{p}_{i}$
denote the average price of tradable consumption in economy $i$,
the total expenditure on non-traded services can be calculated as
$\widetilde{E}_{i}=w_{i}L_{i}-\bar{p}_{i}\mathcal{Q}_{i}$. Correspondingly,
the \emph{share} of service expenditure is given by
\[
\widetilde{e}_{i}\equiv\frac{\widetilde{E}_{i}}{Y_{i}}=\frac{w_{i}L_{i}-\bar{p}_{i}\mathcal{Q}_{i}}{w_{i}L_{i}}=1-\kappa\frac{p_{i}}{\widetilde{p}_{i}},
\]
where $\kappa\equiv K/\tilde{a}$. Given the above equation and data
on $p_{i}/\tilde{p}_{i}$, we can calculate the model's prediction
with regards to the share of service expenditure as $\ln\widehat{1-\tilde{e}_{i}}=\ln\kappa+\ln p_{i}/\widetilde{p}_{i}$.
The data on $p_{i}/\tilde{p}_{i}$ is readily accessible for 182 countries
from the 2011 ICP data, which was described in Section \ref{sec:Empirical-Evidence}.
The same data also reports the factual share of non-traded service
expenditure, $\tilde{e}_{i}$, for each economy. Using this information,
I contrast the predicted values for $\tilde{e}_{i}$ with factual
values in Figure \ref{fig:International-Expenditure-Share}. Encouragingly,
the discrete model performs remarkably well in predicting the cross-national
variation in service expenditure. Finally, note that by allowing for
$\tilde{a}_{i}$ to be a free-moving parameter, I can perfectly match
$\widetilde{e}_{i}$ in all countries. The gains from trade implied
in that case, though, closely resemble those implied by my benchmark
analysis.

\begin{figure}
\caption{\label{fig:International-Expenditure-Share}Discrete model's predictive
power w.r.t. relative service expenditure.}

\begin{centering}
\includegraphics[scale=0.5]{ICP/NT_Expenditure}\vspace{-0.2in}
\par\end{centering}
\noindent\begin{minipage}[t]{1\columnwidth}%
\begin{spacing}{0.8}
\emph{\footnotesize{}Note:}{\footnotesize{} This graph compares data
and predicted values for (relative) service expenditure across 182
countries, including the 31 countries in the WIOD sample. The predicted
values are obtained from a model that $(i)$ treats the entire traded
sector as discrete, and $(ii)$ the entire service sector as one that
complements discrete consumption. See Appendix \ref{sec: Construction-of-Figure}
for more details about the construction of each data point and predicted
value.}
\end{spacing}
%
\end{minipage}
\end{figure}

\section{Classification of the WIOD Industries }

This appendix describes my classification of the 35 industries in
the WIOD sample, which is used to compute $e_{i,\mathcal{C}}$ and
$e_{i,\mathcal{D}}$. My classification groups industries into three
different categories:
\begin{enumerate}
\item Traded industries (as classified by \citet{costinot2014trade});
\item Non-traded industries whose services are bundled with traded discrete
goods in the consumption basket; and
\item Other non-traded service industries. 
\end{enumerate}
The assignment of the WIOD industries into these categories is reported
in Table \ref{tab:List-of-industries}. To briefly explain my classification,
consider the industries supplying transportation services (i.e., industries
numbers 23-24 in Table \ref{tab:List-of-industries}). A consumer
may purchase a personal car, jet, or yacht, and spend the rest of
their transportation budget on other types of ground, air, or sea
transportation services that are infeasible to perform with their
personally-acquired transportation device. Similarly, a consumer may
purchase a personal grill or an RV (Recreational Vehicle) and spend
the rest of their food and travel budget on hotel and restaurant services
(i.e, industry number 22 in Table \ref{tab:List-of-industries}).
Similar examples can be constructed for household- and health-related
services (i.e., industries numbers 32-35 in Table \ref{tab:List-of-industries}).
But the same argument cannot be readily extended to real estate services
or services supplying residential water, electricity, and gas. Relatedly,
financial services are classified as solely non-traded, because they
concern a non-repetitive service that has to be frequently customized
to match the consumer's needs on a case-to-case basis.

\begin{table}[!tbh]
\caption{\label{tab:List-of-industries}Classification of Industries in the
WIOD.}

 \vspace{-0.1in} \begin{center}
\resizebox*{!}{1.25\textwidth}{
\begin{tabular}{|c|>{\raggedright}m{12cm}|c|c|} 
										\hline  
\textbf{WIOD Sector} & \hspace{4cm}  \textbf{Sector's Description} & \specialcell{\textbf{Category}} \\ 
										\toprule  
      1 & Agriculture, Hunting, Forestry and Fishing & \multirow{16}{*}{Traded} \\
	  2 & Mining and Quarrying &  \\
	  3 & Food, Beverages and Tobacco & \\
	  4 & Textiles and Textile Products & \\
	  5 & Leather and Footwear & \\
	  6 & Wood and Products of Wood and Cork & \\
	  7 & Pulp, Paper, Paper , Printing and Publishing & \\
	  8 & Coke, Refined Petroleum and Nuclear Fuel & \\
	  9 & Chemicals and Chemical Products & \\
	  10 & Rubber and Plastics & \\
	  11 & Other Non-Metallic Mineral & \\
	  12 & Basic Metals and Fabricated Metal & \\
	  13 & Machinery, Nec & \\
	  14 & Electrical and Optical Equipment & \\
 	 15 & Transport Equipment & \\
      16 & Manufacturing, Nec; Recycling & \\
\hline 17 & Electricity, Gas and Water Supply & \multirow{2}{*}{Non-Traded Service} \\
	   18 & Construction & \\
\hline \multirow{2}{*}{19} & \multirow{2}{*}{\specialcell{Sale, Maintenance and Repair of Motor Vehicles \\ and Motorcycles; Retail Sale of Fuel}} & \multirow{2}{*}{\specialcell{Non-Traded Service,\\ bundled w/ Tradables}} \\ & & \\
\hline 20 & Wholesale Trade and Commission Trade, Except of Motor Vehicles and Motorcycles & \multirow{2}{*}{Non-Traded Service} \\
 	  21 & Retail Trade, Except of Motor Vehicles and Motorcycles; Repair of Household Goods & \\
\hline 22 & Hotels and Restaurants & \multirow{6}{*}{\specialcell{Non-Traded Service, \\ bundled w/ Tradables}} \\
	   23 & Inland Transport & \\
  	 24 & Water Transport & \\
	   25 & Air Transport & \\
  	 26 & Other Supporting and Auxiliary Transport Activities; Activities of Travel Agencies & \\
 	  27 & Post and Telecommunications & \\
\hline 28 & Financial Intermediation & \multirow{4}{*}{Non-Traded Service} \\
  	 29 & Real Estate Activities & \\
	   30 & Renting of M\&Eq and Other Business Activities & \\
	   31 & Education & \\
\hline 32 & Health and Social Work & \multirow{4}{*}{\specialcell{Non-Traded Service, \\ bundled w/ Tradables}} \\
   	33 & Public Admin and Defence; Compulsory Social Security & \\
  	 34 & Other Community, Social and Personal Services & \\
  	 35 & Private Households with Employed Persons & \\
							\bottomrule 
\end{tabular} }
\end{center}
\end{table}

\end{document}
